\documentclass[11pt,a4paper,oneside,openany]{stacks-project-book}

% Non-rendering source-note environments retained from the Stacks sources.
\usepackage{verbatim}
\newenvironment{reference}{\comment}{\endcomment}
\newenvironment{slogan}{\comment}{\endcomment}
\newenvironment{history}{\comment}{\endcomment}

% Diagrams and chapter-support packages.
\usepackage[all]{xy}
\xyoption{2cell}
\UseAllTwocells
\usepackage{multicol}
\usepackage{graphicx}

% Mainland Simplified-Chinese scientific typography.
\usepackage{fontspec}
\usepackage{xeCJK}
\defaultfontfeatures{Ligatures=TeX}
\setmainfont{Latin Modern Roman}
\setsansfont{Latin Modern Sans}
\setmonofont{Latin Modern Mono}
\setCJKmainfont[AutoFakeSlant=0.2]{Noto Serif CJK SC}
\setCJKsansfont{Noto Sans CJK SC}
\setCJKmonofont{Noto Sans CJK SC}
\newfontfamily\stackszhgreek{Noto Serif CJK SC}
\DeclareRobustCommand{\zhdelta}{{\stackszhgreek δ}}
\xeCJKsetup{PunctStyle=kaiming}
\XeTeXlinebreaklocale "zh"
\usepackage[a4paper,left=22mm,right=22mm,top=24mm,bottom=26mm,headheight=14pt,headsep=7mm,footskip=14mm]{geometry}
\linespread{1.18}
\setlength{\parindent}{2em}
\setlength{\parskip}{0pt}
\setlength{\emergencystretch}{1em}

% Localized structural names.
\renewcommand{\contentsname}{目录}
\renewcommand{\bibname}{参考文献}
\renewcommand{\proofname}{证明}
\renewcommand{\figurename}{图}
\renewcommand{\tablename}{表}
\renewcommand{\appendixname}{附录}

% Chinese numbered chapter heads while keeping numerical identifiers stable.
\makeatletter
\def\@makechapterhead#1{\global\topskip 7.5pc\relax
  \begingroup
  \fontsize{17}{22}\bfseries\centering
    \ifnum\c@secnumdepth>\m@ne
      \leavevmode \hskip-\leftskip
      \rlap{\vbox to\z@{\vss
          \ifx\chaptername\appendixname
            \centerline{\normalsize\mdseries 附录\enspace\thechapter}
          \else
            \centerline{\normalsize\mdseries 第\thechapter 章}
          \fi
          \vskip 3pc}}\hskip\leftskip\fi
     #1\par \endgroup
  \skip@34\p@ \advance\skip@-\normalbaselineskip
  \vskip\skip@ }
\makeatother

% Hyperlinks: current-volume labels resolve locally; absent chapters use tags.
\usepackage{hyperref}
\hypersetup{
  unicode=true,
  hidelinks,
  pdfborder={0 0 0},
  pdftitle={Stacks Project 简体中文版 - 阶段性累积版},
  pdfauthor={The Stacks Project contributors; zh-Hans-CN translation production},
  pdfsubject={Mainland Simplified Chinese cumulative edition of selected Stacks Project chapters}
}
\urlstyle{same}
% Generated from the frozen Stacks permanent-tag registry.
% Existing labels resolve inside this PDF; only absent labels fall back to tags.
\expandafter\def\csname stacksTag@etale-cohomology-definition-0-cech\endcsname{03NR}
\expandafter\def\csname stacksTag@etale-cohomology-definition-Cr\endcsname{03R9}
\expandafter\def\csname stacksTag@etale-cohomology-definition-G-module-continuous\endcsname{04JP}
\expandafter\def\csname stacksTag@etale-cohomology-definition-additive-sheaf\endcsname{03P4}
\expandafter\def\csname stacksTag@etale-cohomology-definition-algebraic-geometric-point\endcsname{03QX}
\expandafter\def\csname stacksTag@etale-cohomology-definition-big-etale-site\endcsname{03PI}
\expandafter\def\csname stacksTag@etale-cohomology-definition-brauer-equivalent\endcsname{03R3}
\expandafter\def\csname stacksTag@etale-cohomology-definition-brauer-group\endcsname{03R5}
\expandafter\def\csname stacksTag@etale-cohomology-definition-c\endcsname{095W}
\expandafter\def\csname stacksTag@etale-cohomology-definition-category-sheaves\endcsname{03NO}
\expandafter\def\csname stacksTag@etale-cohomology-definition-cd\endcsname{0F0Q}
\expandafter\def\csname stacksTag@etale-cohomology-definition-cd-f\endcsname{0F0Z}
\expandafter\def\csname stacksTag@etale-cohomology-definition-cech-complex\endcsname{03OL}
\expandafter\def\csname stacksTag@etale-cohomology-definition-constructible\endcsname{03RW}
\expandafter\def\csname stacksTag@etale-cohomology-definition-ctf\endcsname{03TQ}
\expandafter\def\csname stacksTag@etale-cohomology-definition-descent-datum\endcsname{03O7}
\expandafter\def\csname stacksTag@etale-cohomology-definition-descent-datum-modules\endcsname{03OB}
\expandafter\def\csname stacksTag@etale-cohomology-definition-direct-image-presheaf\endcsname{03PW}
\expandafter\def\csname stacksTag@etale-cohomology-definition-direct-image-sheaf\endcsname{03PY}
\expandafter\def\csname stacksTag@etale-cohomology-definition-effective-modules\endcsname{03OC}
\expandafter\def\csname stacksTag@etale-cohomology-definition-etale-covering\endcsname{03PG}
\expandafter\def\csname stacksTag@etale-cohomology-definition-etale-covering-initial\endcsname{03N5}
\expandafter\def\csname stacksTag@etale-cohomology-definition-etale-local-rings\endcsname{03PS}
\expandafter\def\csname stacksTag@etale-cohomology-definition-etale-morphism\endcsname{03PB}
\expandafter\def\csname stacksTag@etale-cohomology-definition-etale-topos\endcsname{04HQ}
\expandafter\def\csname stacksTag@etale-cohomology-definition-extension-zero\endcsname{03S3}
\expandafter\def\csname stacksTag@etale-cohomology-definition-family-morphisms-fixed-target\endcsname{03NG}
\expandafter\def\csname stacksTag@etale-cohomology-definition-finite-locally-constant\endcsname{03RU}
\expandafter\def\csname stacksTag@etale-cohomology-definition-fpqc-covering\endcsname{03NW}
\expandafter\def\csname stacksTag@etale-cohomology-definition-free-abelian-presheaf\endcsname{03OO}
\expandafter\def\csname stacksTag@etale-cohomology-definition-galois-cohomology\endcsname{04JR}
\expandafter\def\csname stacksTag@etale-cohomology-definition-geometric-point\endcsname{03PO}
\expandafter\def\csname stacksTag@etale-cohomology-definition-henselian\endcsname{03QF}
\expandafter\def\csname stacksTag@etale-cohomology-definition-higher-direct-images\endcsname{04I2}
\expandafter\def\csname stacksTag@etale-cohomology-definition-inverse-image\endcsname{03Q0}
\expandafter\def\csname stacksTag@etale-cohomology-definition-inverse-system-sheaves\endcsname{0EZL}
\expandafter\def\csname stacksTag@etale-cohomology-definition-locally-acyclic\endcsname{0GJP}
\expandafter\def\csname stacksTag@etale-cohomology-definition-presheaf\endcsname{03NC}
\expandafter\def\csname stacksTag@etale-cohomology-definition-ringed-site\endcsname{03OH}
\expandafter\def\csname stacksTag@etale-cohomology-definition-sheaf\endcsname{03NK}
\expandafter\def\csname stacksTag@etale-cohomology-definition-sheaf-property-fpqc\endcsname{03X6}
\expandafter\def\csname stacksTag@etale-cohomology-definition-site\endcsname{03NH}
\expandafter\def\csname stacksTag@etale-cohomology-definition-stalk\endcsname{040R}
\expandafter\def\csname stacksTag@etale-cohomology-definition-standard-etale\endcsname{03PD}
\expandafter\def\csname stacksTag@etale-cohomology-definition-standard-tau\endcsname{03X9}
\expandafter\def\csname stacksTag@etale-cohomology-definition-strictly-henselian\endcsname{03QK}
\expandafter\def\csname stacksTag@etale-cohomology-definition-structure-sheaf\endcsname{04HS}
\expandafter\def\csname stacksTag@etale-cohomology-definition-support\endcsname{04FS}
\expandafter\def\csname stacksTag@etale-cohomology-definition-tau-covering\endcsname{03X8}
\expandafter\def\csname stacksTag@etale-cohomology-definition-tau-site\endcsname{03XB}
\expandafter\def\csname stacksTag@etale-cohomology-definition-trace-map\endcsname{03SE}
\expandafter\def\csname stacksTag@etale-cohomology-definition-variety\endcsname{03RE}
\expandafter\def\csname stacksTag@etale-cohomology-equation-j-p-shriek\endcsname{0F4K}
\expandafter\def\csname stacksTag@etale-cohomology-lemma-F-1\endcsname{0A3L}
\expandafter\def\csname stacksTag@etale-cohomology-lemma-V-C-all-n-etale-fppf\endcsname{0DDS}
\expandafter\def\csname stacksTag@etale-cohomology-lemma-V-C-all-n-etale-ph\endcsname{0DE4}
\expandafter\def\csname stacksTag@etale-cohomology-lemma-affine-only-closed-points\endcsname{09AY}
\expandafter\def\csname stacksTag@etale-cohomology-lemma-all-modules-quasi-coherent\endcsname{0D1W}
\expandafter\def\csname stacksTag@etale-cohomology-lemma-category-is-colimit\endcsname{09YU}
\expandafter\def\csname stacksTag@etale-cohomology-lemma-category-loc-cst-is-colimit\endcsname{0GL2}
\expandafter\def\csname stacksTag@etale-cohomology-lemma-cech-complex\endcsname{03OZ}
\expandafter\def\csname stacksTag@etale-cohomology-lemma-characterize-locally-constant-module\endcsname{0GKC}
\expandafter\def\csname stacksTag@etale-cohomology-lemma-check-constructible\endcsname{095Q}
\expandafter\def\csname stacksTag@etale-cohomology-lemma-cofinal-etale\endcsname{03PQ}
\expandafter\def\csname stacksTag@etale-cohomology-lemma-cohomological-descent-complex-modules\endcsname{0H0U}
\expandafter\def\csname stacksTag@etale-cohomology-lemma-cohomological-dimension-proper\endcsname{095U}
\expandafter\def\csname stacksTag@etale-cohomology-lemma-cohomology-with-support-quasi-coherent\endcsname{0A46}
\expandafter\def\csname stacksTag@etale-cohomology-lemma-compare-cohomology\endcsname{0DDH}
\expandafter\def\csname stacksTag@etale-cohomology-lemma-compare-cohomology-big-small\endcsname{03YX}
\expandafter\def\csname stacksTag@etale-cohomology-lemma-compare-fppf-etale\endcsname{0F0H}
\expandafter\def\csname stacksTag@etale-cohomology-lemma-compare-h-etale\endcsname{0F0J}
\expandafter\def\csname stacksTag@etale-cohomology-lemma-compare-injectives\endcsname{0758}
\expandafter\def\csname stacksTag@etale-cohomology-lemma-compare-ph-etale\endcsname{0F0I}
\expandafter\def\csname stacksTag@etale-cohomology-lemma-compare-structure-sheaves\endcsname{07AJ}
\expandafter\def\csname stacksTag@etale-cohomology-lemma-connected-ctf-locally-constant\endcsname{09BI}
\expandafter\def\csname stacksTag@etale-cohomology-lemma-connected-locally-constant\endcsname{09BF}
\expandafter\def\csname stacksTag@etale-cohomology-lemma-constructible-abelian\endcsname{03RZ}
\expandafter\def\csname stacksTag@etale-cohomology-lemma-constructible-quasi-compact-quasi-separated\endcsname{095E}
\expandafter\def\csname stacksTag@etale-cohomology-lemma-decompose-quasi-finite-morphism\endcsname{095K}
\expandafter\def\csname stacksTag@etale-cohomology-lemma-descent-sheaf-fppf-etale\endcsname{0DEU}
\expandafter\def\csname stacksTag@etale-cohomology-lemma-describe-etale-local-ring\endcsname{04HX}
\expandafter\def\csname stacksTag@etale-cohomology-lemma-describe-pullback\endcsname{09XN}
\expandafter\def\csname stacksTag@etale-cohomology-lemma-describe-pullback-pi-fppf\endcsname{0DDL}
\expandafter\def\csname stacksTag@etale-cohomology-lemma-describe-pullback-pi-ph\endcsname{0DDW}
\expandafter\def\csname stacksTag@etale-cohomology-lemma-direct-sum-bounded-below-pushforward\endcsname{0GIW}
\expandafter\def\csname stacksTag@etale-cohomology-lemma-directed-colimit-cohomology\endcsname{03Q6}
\expandafter\def\csname stacksTag@etale-cohomology-lemma-etale-topos-independent-partial-universe\endcsname{0958}
\expandafter\def\csname stacksTag@etale-cohomology-lemma-ext-modules-hom\endcsname{0DVE}
\expandafter\def\csname stacksTag@etale-cohomology-lemma-faithful\endcsname{04LW}
\expandafter\def\csname stacksTag@etale-cohomology-lemma-finite-dim-group-cohomology\endcsname{0DVF}
\expandafter\def\csname stacksTag@etale-cohomology-lemma-finite-pushforward-constructible\endcsname{095R}
\expandafter\def\csname stacksTag@etale-cohomology-lemma-glue-etale-sheaf-fppf-cover\endcsname{0GF2}
\expandafter\def\csname stacksTag@etale-cohomology-lemma-glue-etale-sheaf-integral-surjective\endcsname{0GEZ}
\expandafter\def\csname stacksTag@etale-cohomology-lemma-glue-etale-sheaf-proper-surjective\endcsname{0GF0}
\expandafter\def\csname stacksTag@etale-cohomology-lemma-h0-proper-over-henselian-pair\endcsname{0A0C}
\expandafter\def\csname stacksTag@etale-cohomology-lemma-jshriek-constructible\endcsname{03S8}
\expandafter\def\csname stacksTag@etale-cohomology-lemma-jshriek-open\endcsname{0F70}
\expandafter\def\csname stacksTag@etale-cohomology-lemma-morphism-locally-ringed\endcsname{04I5}
\expandafter\def\csname stacksTag@etale-cohomology-lemma-points-fppf\endcsname{06VX}
\expandafter\def\csname stacksTag@etale-cohomology-lemma-points-small-etale-site\endcsname{04HU}
\expandafter\def\csname stacksTag@etale-cohomology-lemma-projection-formula-proper-mod-n\endcsname{0F0G}
\expandafter\def\csname stacksTag@etale-cohomology-lemma-proper-base-change\endcsname{0DDE}
\expandafter\def\csname stacksTag@etale-cohomology-lemma-proper-base-change-f-star\endcsname{0A3U}
\expandafter\def\csname stacksTag@etale-cohomology-lemma-proper-base-change-mod-n\endcsname{0F0C}
\expandafter\def\csname stacksTag@etale-cohomology-lemma-proper-base-change-stalk\endcsname{0DDF}
\expandafter\def\csname stacksTag@etale-cohomology-lemma-proper-mod-n-direct-sums\endcsname{0F0D}
\expandafter\def\csname stacksTag@etale-cohomology-lemma-proper-pushforward-stalk\endcsname{0A3T}
\expandafter\def\csname stacksTag@etale-cohomology-lemma-pullback-filtered-modules\endcsname{0GJ0}
\expandafter\def\csname stacksTag@etale-cohomology-lemma-pullback-on-h2-curve\endcsname{0AMB}
\expandafter\def\csname stacksTag@etale-cohomology-lemma-relative-colimit-general\endcsname{0EYM}
\expandafter\def\csname stacksTag@etale-cohomology-lemma-review-quasi-coherent\endcsname{0DEW}
\expandafter\def\csname stacksTag@etale-cohomology-lemma-sections-upstairs\endcsname{0A3H}
\expandafter\def\csname stacksTag@etale-cohomology-lemma-ses-associated-to-open\endcsname{095L}
\expandafter\def\csname stacksTag@etale-cohomology-lemma-shriek-into-star-separated-etale\endcsname{0F4L}
\expandafter\def\csname stacksTag@etale-cohomology-lemma-sp-isom-proper-loc-cst-torsion\endcsname{0GKD}
\expandafter\def\csname stacksTag@etale-cohomology-lemma-sp-isom-proper-torsion-loc-ac\endcsname{0GJW}
\expandafter\def\csname stacksTag@etale-cohomology-lemma-stalk-exact\endcsname{03PT}
\expandafter\def\csname stacksTag@etale-cohomology-lemma-stalk-pullback\endcsname{03Q1}
\expandafter\def\csname stacksTag@etale-cohomology-lemma-stalk-pushforward-closed-immersion\endcsname{04FX}
\expandafter\def\csname stacksTag@etale-cohomology-lemma-support-section-closed\endcsname{04FT}
\expandafter\def\csname stacksTag@etale-cohomology-lemma-supported-in-closed-points\endcsname{0F1F}
\expandafter\def\csname stacksTag@etale-cohomology-lemma-tensor-c\endcsname{0961}
\expandafter\def\csname stacksTag@etale-cohomology-lemma-tensor-constructible\endcsname{0GKB}
\expandafter\def\csname stacksTag@etale-cohomology-lemma-torsion-direct-image\endcsname{0DDD}
\expandafter\def\csname stacksTag@etale-cohomology-lemma-what-integral\endcsname{04C2}
\expandafter\def\csname stacksTag@etale-cohomology-lemma-when-ctf\endcsname{03TT}
\expandafter\def\csname stacksTag@etale-cohomology-lemma-zero-over-image\endcsname{04FR}
\expandafter\def\csname stacksTag@etale-cohomology-proposition-closed-immersion-pushforward\endcsname{04CA}
\expandafter\def\csname stacksTag@etale-cohomology-proposition-constructible-over-noetherian\endcsname{09BH}
\expandafter\def\csname stacksTag@etale-cohomology-proposition-describe-jshriek\endcsname{03S5}
\expandafter\def\csname stacksTag@etale-cohomology-proposition-finite-higher-direct-image-zero\endcsname{03QP}
\expandafter\def\csname stacksTag@etale-cohomology-proposition-integral-universally-injective-pushforward\endcsname{04FZ}
\expandafter\def\csname stacksTag@etale-cohomology-proposition-quasi-coherent-sheaf-fpqc\endcsname{03OG}
\expandafter\def\csname stacksTag@etale-cohomology-remark-affine-inside-equivalence\endcsname{05YX}
\expandafter\def\csname stacksTag@etale-cohomology-remark-constant-locally-constant-maps\endcsname{03P5}
\expandafter\def\csname stacksTag@etale-cohomology-remark-stalk-pullback\endcsname{04JN}
\expandafter\def\csname stacksTag@etale-cohomology-remarks-enough-points\endcsname{040S}
\expandafter\def\csname stacksTag@etale-cohomology-section-Dc\endcsname{095V}
\expandafter\def\csname stacksTag@etale-cohomology-section-artin-schreier\endcsname{0A3J}
\expandafter\def\csname stacksTag@etale-cohomology-section-base-change-preliminaries\endcsname{0EZQ}
\expandafter\def\csname stacksTag@etale-cohomology-section-cohomology-support\endcsname{09XP}
\expandafter\def\csname stacksTag@etale-cohomology-section-compare\endcsname{0757}
\expandafter\def\csname stacksTag@etale-cohomology-section-compare-topologies\endcsname{09XL}
\expandafter\def\csname stacksTag@etale-cohomology-section-computation\endcsname{03N8}
\expandafter\def\csname stacksTag@etale-cohomology-section-continuous-group-cohomology\endcsname{0DVG}
\expandafter\def\csname stacksTag@etale-cohomology-section-extension-by-zero\endcsname{03S2}
\expandafter\def\csname stacksTag@etale-cohomology-section-fppf-etale\endcsname{0DDK}
\expandafter\def\csname stacksTag@etale-cohomology-section-functoriality\endcsname{04I0}
\expandafter\def\csname stacksTag@etale-cohomology-section-galois-action-stalks\endcsname{03QW}
\expandafter\def\csname stacksTag@etale-cohomology-section-glueing-etale\endcsname{0GEX}
\expandafter\def\csname stacksTag@etale-cohomology-section-introduction\endcsname{03N2}
\expandafter\def\csname stacksTag@etale-cohomology-section-ph-etale\endcsname{0DDV}
\expandafter\def\csname stacksTag@etale-cohomology-section-proper-base-change\endcsname{095S}
\expandafter\def\csname stacksTag@etale-cohomology-section-stalks\endcsname{03PN}
\expandafter\def\csname stacksTag@etale-cohomology-section-stalks-structure-sheaf\endcsname{04HW}
\expandafter\def\csname stacksTag@etale-cohomology-section-trace-method\endcsname{03SH}
\expandafter\def\csname stacksTag@etale-cohomology-section-vanishing-torsion\endcsname{03SB}
\expandafter\def\csname stacksTag@etale-cohomology-theorem-colimit\endcsname{09YQ}
\expandafter\def\csname stacksTag@etale-cohomology-theorem-equivalence-sheaves-point\endcsname{03QT}
\expandafter\def\csname stacksTag@etale-cohomology-theorem-exactness-stalks\endcsname{03PU}
\expandafter\def\csname stacksTag@etale-cohomology-theorem-fully-faithful\endcsname{04I7}
\expandafter\def\csname stacksTag@etale-cohomology-theorem-higher-direct-images\endcsname{03Q9}
\expandafter\def\csname stacksTag@etale-cohomology-theorem-proper-base-change\endcsname{095T}
\expandafter\def\csname stacksTag@etale-cohomology-theorem-quasi-coherent\endcsname{03OJ}
\expandafter\def\csname stacksTag@etale-cohomology-theorem-topological-invariance\endcsname{04DZ}
\expandafter\def\csname stacksTag@etale-cohomology-theorem-vanishing-affine-curves\endcsname{03SC}
\expandafter\def\csname stacksTag@etale-cohomology-theorem-vanishing-affine-curves-coefficients\endcsname{0GJI}
\expandafter\def\csname stacksTag@pione-definition-G-set-continuous\endcsname{04JJ}
\expandafter\def\csname stacksTag@pione-definition-fundamental-group\endcsname{0BNC}
\expandafter\def\csname stacksTag@pione-definition-galois-category\endcsname{0BMY}
\expandafter\def\csname stacksTag@pione-lemma-inertial-invariants-unramified\endcsname{0BSU}
\expandafter\def\csname stacksTag@pione-section-finite-etale\endcsname{0BL6}
\expandafter\def\csname stacksTag@pione-section-group-actions-integral\endcsname{0BSN}
\expandafter\def\csname stacksTag@pione-section-introduction\endcsname{0BQ7}
\expandafter\def\csname stacksTag@spaces-more-groupoids-definition-split-at-point\endcsname{04RK}
\expandafter\def\csname stacksTag@spaces-more-groupoids-lemma-composition-is-addition\endcsname{0CKF}
\expandafter\def\csname stacksTag@spaces-more-groupoids-lemma-group-scheme-over-field-separated\endcsname{08BH}
\expandafter\def\csname stacksTag@spaces-more-groupoids-lemma-group-space-scheme-locally-finite-type-over-k\endcsname{0B8F}
\expandafter\def\csname stacksTag@spaces-more-groupoids-lemma-groupoid-on-field-dimension-equal-stabilizer\endcsname{06FF}
\expandafter\def\csname stacksTag@spaces-more-groupoids-lemma-no-nonconstant-morphism-from-P1-to-group\endcsname{0AEN}
\expandafter\def\csname stacksTag@spaces-more-groupoids-lemma-property-G-invariant\endcsname{06R4}
\expandafter\def\csname stacksTag@spaces-more-groupoids-lemma-quasi-splitting-affine-scheme\endcsname{04S0}
\expandafter\def\csname stacksTag@spaces-more-groupoids-lemma-restrict-preserves-type\endcsname{04RP}
\expandafter\def\csname stacksTag@spaces-more-groupoids-lemma-splitting-affine-scheme\endcsname{0DTE}
\expandafter\def\csname stacksTag@spaces-more-groupoids-lemma-strong-splitting-affine-scheme\endcsname{04RZ}
\expandafter\def\csname stacksTag@spaces-more-groupoids-proposition-finite-algebraic-space\endcsname{04QH}
\expandafter\def\csname stacksTag@spaces-more-groupoids-proposition-group-space-scheme-over-field\endcsname{0B8G}
\expandafter\def\csname stacksTag@spaces-more-groupoids-section-etale-localize\endcsname{04RJ}
\expandafter\def\csname stacksTag@spaces-more-groupoids-section-finite\endcsname{04PB}
\expandafter\def\csname stacksTag@spaces-more-groupoids-section-groupoid-sections\endcsname{0CKB}
\expandafter\def\csname stacksTag@varieties-definition-absolute-frobenius\endcsname{03SM}
\expandafter\def\csname stacksTag@varieties-definition-algebraic-scheme\endcsname{06LG}
\expandafter\def\csname stacksTag@varieties-definition-curve\endcsname{0A23}
\expandafter\def\csname stacksTag@varieties-definition-degree\endcsname{0BEW}
\expandafter\def\csname stacksTag@varieties-definition-degree-invertible-sheaf\endcsname{0AYR}
\expandafter\def\csname stacksTag@varieties-definition-delta-invariant\endcsname{0C1T}
\expandafter\def\csname stacksTag@varieties-definition-delta-invariant-algebra\endcsname{0C3T}
\expandafter\def\csname stacksTag@varieties-definition-dual-numbers\endcsname{0B29}
\expandafter\def\csname stacksTag@varieties-definition-embed-dim\endcsname{0C1Q}
\expandafter\def\csname stacksTag@varieties-definition-embedding-dimension\endcsname{0C2H}
\expandafter\def\csname stacksTag@varieties-definition-euler-characteristic\endcsname{0BEJ}
\expandafter\def\csname stacksTag@varieties-definition-geometrically-connected\endcsname{0362}
\expandafter\def\csname stacksTag@varieties-definition-geometrically-integral\endcsname{020H}
\expandafter\def\csname stacksTag@varieties-definition-geometrically-irreducible\endcsname{0365}
\expandafter\def\csname stacksTag@varieties-definition-geometrically-normal\endcsname{038M}
\expandafter\def\csname stacksTag@varieties-definition-geometrically-reduced\endcsname{035V}
\expandafter\def\csname stacksTag@varieties-definition-geometrically-regular\endcsname{038T}
\expandafter\def\csname stacksTag@varieties-definition-hilbert-polynomial\endcsname{08AD}
\expandafter\def\csname stacksTag@varieties-definition-intersection-number\endcsname{0BEP}
\expandafter\def\csname stacksTag@varieties-definition-regularity\endcsname{08A3}
\expandafter\def\csname stacksTag@varieties-definition-relative-frobenius\endcsname{0CC9}
\expandafter\def\csname stacksTag@varieties-definition-tangent-space\endcsname{0B2C}
\expandafter\def\csname stacksTag@varieties-definition-variety\endcsname{020D}
\expandafter\def\csname stacksTag@varieties-definition-variety-type\endcsname{04L1}
\expandafter\def\csname stacksTag@varieties-definition-wedge\endcsname{0C41}
\expandafter\def\csname stacksTag@varieties-equation-tangent-space-fibre\endcsname{0BEA}
\expandafter\def\csname stacksTag@varieties-lemma-CM-base-change\endcsname{045P}
\expandafter\def\csname stacksTag@varieties-lemma-Galois-action-quasi-compact-open\endcsname{04KU}
\expandafter\def\csname stacksTag@varieties-lemma-affine-geometrically-connected\endcsname{0386}
\expandafter\def\csname stacksTag@varieties-lemma-affine-space-over-field\endcsname{055T}
\expandafter\def\csname stacksTag@varieties-lemma-algebraic-scheme-dim-0\endcsname{06LH}
\expandafter\def\csname stacksTag@varieties-lemma-ample-after-field-extension\endcsname{0BDC}
\expandafter\def\csname stacksTag@varieties-lemma-ample-curve\endcsname{0B5X}
\expandafter\def\csname stacksTag@varieties-lemma-ample-positive\endcsname{0BEV}
\expandafter\def\csname stacksTag@varieties-lemma-ampleness-in-terms-of-degrees-components\endcsname{0B5Y}
\expandafter\def\csname stacksTag@varieties-lemma-automorphism\endcsname{0G05}
\expandafter\def\csname stacksTag@varieties-lemma-baby-stein\endcsname{0FD1}
\expandafter\def\csname stacksTag@varieties-lemma-base-change-complete-local-ring\endcsname{0C54}
\expandafter\def\csname stacksTag@varieties-lemma-bertini\endcsname{0FD6}
\expandafter\def\csname stacksTag@varieties-lemma-bijection-connected-components\endcsname{0385}
\expandafter\def\csname stacksTag@varieties-lemma-bijection-irreducible-components\endcsname{038F}
\expandafter\def\csname stacksTag@varieties-lemma-bound-quotients-free\endcsname{08AG}
\expandafter\def\csname stacksTag@varieties-lemma-change-fields-pic\endcsname{0CC5}
\expandafter\def\csname stacksTag@varieties-lemma-char-p-differentials-free-smooth\endcsname{04QP}
\expandafter\def\csname stacksTag@varieties-lemma-char-zero-differentials-free-smooth\endcsname{04QN}
\expandafter\def\csname stacksTag@varieties-lemma-characterize-geometrically-connected\endcsname{0387}
\expandafter\def\csname stacksTag@varieties-lemma-characterize-geometrically-disconnected\endcsname{0389}
\expandafter\def\csname stacksTag@varieties-lemma-check-invertible-sheaf-trivial\endcsname{0B40}
\expandafter\def\csname stacksTag@varieties-lemma-chi-tensor-finite\endcsname{0AYT}
\expandafter\def\csname stacksTag@varieties-lemma-closed-fixed-by-Galois\endcsname{038B}
\expandafter\def\csname stacksTag@varieties-lemma-closure-image-product-map\endcsname{04Q0}
\expandafter\def\csname stacksTag@varieties-lemma-closure-of-product\endcsname{047B}
\expandafter\def\csname stacksTag@varieties-lemma-complement-codim-1-closed-points\endcsname{09NA}
\expandafter\def\csname stacksTag@varieties-lemma-complete-local-ring-structure-as-algebra\endcsname{0C52}
\expandafter\def\csname stacksTag@varieties-lemma-connectedness-ample-divisor\endcsname{0FD9}
\expandafter\def\csname stacksTag@varieties-lemma-degree-additive\endcsname{0AYS}
\expandafter\def\csname stacksTag@varieties-lemma-degree-base-change\endcsname{0B59}
\expandafter\def\csname stacksTag@varieties-lemma-degree-birational-pullback\endcsname{0AYU}
\expandafter\def\csname stacksTag@varieties-lemma-degree-effective-Cartier-divisor\endcsname{0AYY}
\expandafter\def\csname stacksTag@varieties-lemma-degree-finite-morphism-in-terms-degrees\endcsname{0BEX}
\expandafter\def\csname stacksTag@varieties-lemma-degree-in-terms-of-components\endcsname{0AYW}
\expandafter\def\csname stacksTag@varieties-lemma-degree-on-proper-curve\endcsname{0AYV}
\expandafter\def\csname stacksTag@varieties-lemma-degree-pullback-map-proper-curves\endcsname{0AYZ}
\expandafter\def\csname stacksTag@varieties-lemma-degree-tensor-product\endcsname{0AYX}
\expandafter\def\csname stacksTag@varieties-lemma-delta-invariant-and-change-of-fields\endcsname{0C3X}
\expandafter\def\csname stacksTag@varieties-lemma-delta-invariant-and-change-of-fields-better\endcsname{0C3Y}
\expandafter\def\csname stacksTag@varieties-lemma-delta-invariant-is-zero\endcsname{0C3U}
\expandafter\def\csname stacksTag@varieties-lemma-delta-number-branches-inequality\endcsname{0C43}
\expandafter\def\csname stacksTag@varieties-lemma-delta-number-branches-inequality-sh\endcsname{0C42}
\expandafter\def\csname stacksTag@varieties-lemma-delta-same-after-completion\endcsname{0C3W}
\expandafter\def\csname stacksTag@varieties-lemma-dim-1-projective-completion\endcsname{0BXV}
\expandafter\def\csname stacksTag@varieties-lemma-dim-1-projective-completion-after-insep\endcsname{0GK5}
\expandafter\def\csname stacksTag@varieties-lemma-dim-1-proper-projective\endcsname{0A26}
\expandafter\def\csname stacksTag@varieties-lemma-dim-1-quasi-projective\endcsname{0A25}
\expandafter\def\csname stacksTag@varieties-lemma-dimension-fibre-in-higher-codimension\endcsname{0B2K}
\expandafter\def\csname stacksTag@varieties-lemma-dimension-fibres-locally-algebraic\endcsname{0B2L}
\expandafter\def\csname stacksTag@varieties-lemma-dimension-locally-algebraic\endcsname{0A21}
\expandafter\def\csname stacksTag@varieties-lemma-dimension-product-locally-algebraic\endcsname{0B2M}
\expandafter\def\csname stacksTag@varieties-lemma-divisible\endcsname{0C6P}
\expandafter\def\csname stacksTag@varieties-lemma-dominate-valuation-ring-dimension-fibres\endcsname{0B2J}
\expandafter\def\csname stacksTag@varieties-lemma-equation-codim-1-in-projective-space\endcsname{0BXU}
\expandafter\def\csname stacksTag@varieties-lemma-euler-characteristic-additive\endcsname{08AA}
\expandafter\def\csname stacksTag@varieties-lemma-exact-sequence-induction\endcsname{08A0}
\expandafter\def\csname stacksTag@varieties-lemma-finite-extension-geometrically-irreducible-components\endcsname{054R}
\expandafter\def\csname stacksTag@varieties-lemma-finite-extension-geometrically-normal\endcsname{0BXT}
\expandafter\def\csname stacksTag@varieties-lemma-finite-extension-geometrically-reduced\endcsname{04KT}
\expandafter\def\csname stacksTag@varieties-lemma-finite-in-codim-1\endcsname{0AB7}
\expandafter\def\csname stacksTag@varieties-lemma-flat-under-smooth\endcsname{05AX}
\expandafter\def\csname stacksTag@varieties-lemma-general-degree-g-line-bundle\endcsname{0B8Z}
\expandafter\def\csname stacksTag@varieties-lemma-generate-over-complement\endcsname{0FD5}
\expandafter\def\csname stacksTag@varieties-lemma-generic-points-geometrically-reduced\endcsname{04KS}
\expandafter\def\csname stacksTag@varieties-lemma-geometric-branches-and-change-of-fields\endcsname{0C40}
\expandafter\def\csname stacksTag@varieties-lemma-geometric-structure-unramified\endcsname{0CBJ}
\expandafter\def\csname stacksTag@varieties-lemma-geometrically-connected-check-after-extension\endcsname{054N}
\expandafter\def\csname stacksTag@varieties-lemma-geometrically-connected-criterion\endcsname{056R}
\expandafter\def\csname stacksTag@varieties-lemma-geometrically-connected-if-connected-and-point\endcsname{04KV}
\expandafter\def\csname stacksTag@varieties-lemma-geometrically-integral\endcsname{038K}
\expandafter\def\csname stacksTag@varieties-lemma-geometrically-irreducible-check-after-extension\endcsname{054P}
\expandafter\def\csname stacksTag@varieties-lemma-geometrically-irreducible-function-field\endcsname{054Q}
\expandafter\def\csname stacksTag@varieties-lemma-geometrically-normal-in-codim-1\endcsname{0C3P}
\expandafter\def\csname stacksTag@varieties-lemma-geometrically-normal-stein\endcsname{0FD3}
\expandafter\def\csname stacksTag@varieties-lemma-geometrically-normal-upstairs\endcsname{038P}
\expandafter\def\csname stacksTag@varieties-lemma-geometrically-reduced\endcsname{035X}
\expandafter\def\csname stacksTag@varieties-lemma-geometrically-reduced-any-base-change\endcsname{035Z}
\expandafter\def\csname stacksTag@varieties-lemma-geometrically-reduced-at-point\endcsname{035W}
\expandafter\def\csname stacksTag@varieties-lemma-geometrically-reduced-dense-smooth-open\endcsname{056V}
\expandafter\def\csname stacksTag@varieties-lemma-geometrically-reduced-upstairs\endcsname{0384}
\expandafter\def\csname stacksTag@varieties-lemma-geometrically-regular\endcsname{038V}
\expandafter\def\csname stacksTag@varieties-lemma-geometrically-regular-smooth\endcsname{038X}
\expandafter\def\csname stacksTag@varieties-lemma-geometrically-regular-upstairs\endcsname{038W}
\expandafter\def\csname stacksTag@varieties-lemma-geometrically-unibranch-and-change-of-fields\endcsname{0C55}
\expandafter\def\csname stacksTag@varieties-lemma-globally-generated-base-change\endcsname{0B57}
\expandafter\def\csname stacksTag@varieties-lemma-image-connected-component\endcsname{04PZ}
\expandafter\def\csname stacksTag@varieties-lemma-image-irreducible\endcsname{04KX}
\expandafter\def\csname stacksTag@varieties-lemma-immersion-into-affine\endcsname{0CBL}
\expandafter\def\csname stacksTag@varieties-lemma-injective-tangent-spaces-unramified\endcsname{0B2G}
\expandafter\def\csname stacksTag@varieties-lemma-inseparable-deg-p-smooth\endcsname{0CCF}
\expandafter\def\csname stacksTag@varieties-lemma-integral-closure-ground-field\endcsname{04MI}
\expandafter\def\csname stacksTag@varieties-lemma-intersection-in-projective-space\endcsname{0B2R}
\expandafter\def\csname stacksTag@varieties-lemma-intersection-number-additive\endcsname{0BER}
\expandafter\def\csname stacksTag@varieties-lemma-intersection-number-and-pullback\endcsname{0BET}
\expandafter\def\csname stacksTag@varieties-lemma-intersection-numbers-and-degrees-on-curves\endcsname{0BEY}
\expandafter\def\csname stacksTag@varieties-lemma-irreducible-components-geometrically-irreducible\endcsname{054S}
\expandafter\def\csname stacksTag@varieties-lemma-kunneth\endcsname{0BED}
\expandafter\def\csname stacksTag@varieties-lemma-locally-Noetherian-base-change\endcsname{038R}
\expandafter\def\csname stacksTag@varieties-lemma-locally-finite-type-Jacobson\endcsname{0478}
\expandafter\def\csname stacksTag@varieties-lemma-m-regular-globally-generated\endcsname{08A8}
\expandafter\def\csname stacksTag@varieties-lemma-m-regular-up\endcsname{08A6}
\expandafter\def\csname stacksTag@varieties-lemma-make-Jacobson\endcsname{0479}
\expandafter\def\csname stacksTag@varieties-lemma-map-tangent-spaces\endcsname{0B2F}
\expandafter\def\csname stacksTag@varieties-lemma-modification-normal-iso-over-codimension-1\endcsname{0BFP}
\expandafter\def\csname stacksTag@varieties-lemma-no-sections-dual-nef\endcsname{0E22}
\expandafter\def\csname stacksTag@varieties-lemma-normalization-locally-algebraic\endcsname{0BXR}
\expandafter\def\csname stacksTag@varieties-lemma-normalization-same-after-completion\endcsname{0C3V}
\expandafter\def\csname stacksTag@varieties-lemma-numerical-intersection-effective-Cartier-divisor\endcsname{0BEU}
\expandafter\def\csname stacksTag@varieties-lemma-numerical-polynomial-from-euler\endcsname{0BEM}
\expandafter\def\csname stacksTag@varieties-lemma-numerical-polynomial-leading-term\endcsname{0BEN}
\expandafter\def\csname stacksTag@varieties-lemma-orbit-irreducible-components\endcsname{04KY}
\expandafter\def\csname stacksTag@varieties-lemma-perfect-reduced\endcsname{020I}
\expandafter\def\csname stacksTag@varieties-lemma-pre-delta-invariant\endcsname{0C3S}
\expandafter\def\csname stacksTag@varieties-lemma-prepare-delta-invariant\endcsname{0C1R}
\expandafter\def\csname stacksTag@varieties-lemma-product-varieties\endcsname{05P3}
\expandafter\def\csname stacksTag@varieties-lemma-proper-geometrically-reduced-global-sections\endcsname{0BUG}
\expandafter\def\csname stacksTag@varieties-lemma-quasi-finite-in-codim-1\endcsname{0AB6}
\expandafter\def\csname stacksTag@varieties-lemma-rational-equivalence-for-Pic\endcsname{0BEH}
\expandafter\def\csname stacksTag@varieties-lemma-reduced-dim-1-projective-completion\endcsname{0BXW}
\expandafter\def\csname stacksTag@varieties-lemma-regular-functions-proper-variety\endcsname{04L2}
\expandafter\def\csname stacksTag@varieties-lemma-regular-point-on-curve\endcsname{0B8Y}
\expandafter\def\csname stacksTag@varieties-lemma-relative-frobenius\endcsname{0CCB}
\expandafter\def\csname stacksTag@varieties-lemma-relative-frobenius-endomorphism-identity\endcsname{0CCA}
\expandafter\def\csname stacksTag@varieties-lemma-relative-frobenius-finite\endcsname{0CCD}
\expandafter\def\csname stacksTag@varieties-lemma-relative-frobenius-omega\endcsname{0CCC}
\expandafter\def\csname stacksTag@varieties-lemma-scheme-theoretic-image-product-map\endcsname{04Q1}
\expandafter\def\csname stacksTag@varieties-lemma-separably-closed-field-connected-components\endcsname{0363}
\expandafter\def\csname stacksTag@varieties-lemma-separably-closed-field-irreducible-components\endcsname{038H}
\expandafter\def\csname stacksTag@varieties-lemma-separably-closed-irreducible\endcsname{020J}
\expandafter\def\csname stacksTag@varieties-lemma-smooth-geometrically-normal\endcsname{056T}
\expandafter\def\csname stacksTag@varieties-lemma-smooth-regular\endcsname{056S}
\expandafter\def\csname stacksTag@varieties-lemma-smooth-separable-closed-points-dense\endcsname{056U}
\expandafter\def\csname stacksTag@varieties-lemma-smooth-stratification\endcsname{0H3W}
\expandafter\def\csname stacksTag@varieties-lemma-surjective-pic-birational-finite\endcsname{0C0T}
\expandafter\def\csname stacksTag@varieties-lemma-tangent-space-cotangent-space\endcsname{0B2D}
\expandafter\def\csname stacksTag@varieties-lemma-tangent-space-product\endcsname{0BEB}
\expandafter\def\csname stacksTag@varieties-lemma-vanishin-h0-negative\endcsname{0FD7}
\expandafter\def\csname stacksTag@varieties-lemma-vanishing-degree-2g-and-1-line-bundle\endcsname{0B90}
\expandafter\def\csname stacksTag@varieties-lemma-variant-separate-points-tangent-vectors\endcsname{0E8T}
\expandafter\def\csname stacksTag@varieties-lemma-variety-with-smooth-rational-point\endcsname{0CDW}
\expandafter\def\csname stacksTag@varieties-lemma-very-ample-vanish-at-point\endcsname{0B58}
\expandafter\def\csname stacksTag@varieties-proposition-finite-set-of-points-of-codim-1-in-affine\endcsname{09NN}
\expandafter\def\csname stacksTag@varieties-proposition-unique-base-field\endcsname{04MK}
\expandafter\def\csname stacksTag@varieties-remark-n-fold-relative-frobenius\endcsname{0CCG}
\expandafter\def\csname stacksTag@varieties-section-CM\endcsname{045O}
\expandafter\def\csname stacksTag@varieties-section-algebraic-schemes\endcsname{06LF}
\expandafter\def\csname stacksTag@varieties-section-bertini\endcsname{0FD4}
\expandafter\def\csname stacksTag@varieties-section-curves\endcsname{0A22}
\expandafter\def\csname stacksTag@varieties-section-dimension-one\endcsname{09N7}
\expandafter\def\csname stacksTag@varieties-section-divisors-curves\endcsname{0AYQ}
\expandafter\def\csname stacksTag@varieties-section-embedding-dimension\endcsname{0C2G}
\expandafter\def\csname stacksTag@varieties-section-frobenius\endcsname{0CC6}
\expandafter\def\csname stacksTag@varieties-section-generically-finite\endcsname{0AB5}
\expandafter\def\csname stacksTag@varieties-section-kunneth\endcsname{0BEC}
\expandafter\def\csname stacksTag@varieties-section-normalization\endcsname{0BXQ}
\expandafter\def\csname stacksTag@varieties-section-normalization-one-dimensional\endcsname{0C44}
\expandafter\def\csname stacksTag@varieties-section-num\endcsname{0BEL}
\expandafter\def\csname stacksTag@varieties-section-number-of-branches\endcsname{0C3Z}
\expandafter\def\csname stacksTag@varieties-section-picard-groups\endcsname{0BEG}
\expandafter\def\csname stacksTag@varieties-section-separating-points-tangent-vectors\endcsname{0E8R}
\expandafter\def\csname stacksTag@varieties-section-smooth\endcsname{04QM}
\expandafter\def\csname stacksTag@varieties-section-varieties\endcsname{020C}
\expandafter\def\csname stacksTag@varieties-section-varieties-rational-maps\endcsname{0BXM}
\expandafter\def\csname stacksTag@varieties-situation-family-divisors\endcsname{0G47}
\expandafter\def\csname stacksTag@varieties-subsection-hilbert\endcsname{08A9}
\expandafter\def\csname stacksTag@varieties-subsection-regularity\endcsname{08A2}
\expandafter\def\csname stacksTag@varieties-theorem-varieties-rational-maps\endcsname{0BXN}
\expandafter\def\csname stacksManual@foo-lemma-bar\endcsname{\href{https://github.com/stacks/stacks-project/blob/a04446e57ec1fbc252a871afcec7752fb2807b14/coding.tex}{\texttt{源码示例}}}
\DeclareRobustCommand{\StacksResolvedRef}[1]{%
  \ifcsname r@#1\endcsname
    \StacksOriginalRef{#1}%
  \else
    \ifcsname stacksTag@#1\endcsname
      \href{https://stacks.math.columbia.edu/tag/\csname stacksTag@#1\endcsname}{\texttt{\csname stacksTag@#1\endcsname}}%
    \else
      \ifcsname stacksManual@#1\endcsname
        \csname stacksManual@#1\endcsname%
      \else
        \StacksOriginalRef{#1}%
      \fi
    \fi
  \fi
}
\AtBeginDocument{%
  \let\StacksOriginalRef\ref
  \let\ref\StacksResolvedRef
}


% Theorem environments.
%
\theoremstyle{plain}
\newtheorem{theorem}[subsection]{定理}
\newtheorem{proposition}[subsection]{命题}
\newtheorem{lemma}[subsection]{引理}

\theoremstyle{definition}
\newtheorem{definition}[subsection]{定义}
\newtheorem{example}[subsection]{例}
\newtheorem{exercise}[subsection]{习题}
\newtheorem{situation}[subsection]{情形}

\theoremstyle{remark}
\newtheorem{remark}[subsection]{注}
\newtheorem{remarks}[subsection]{诸注}

\numberwithin{equation}{subsection}

% Macros
%
\def\lim{\mathop{\mathrm{lim}}\nolimits}
\def\colim{\mathop{\mathrm{colim}}\nolimits}
\def\Spec{\mathop{\mathrm{Spec}}}
\def\Hom{\mathop{\mathrm{Hom}}\nolimits}
\def\Ext{\mathop{\mathrm{Ext}}\nolimits}
\def\SheafHom{\mathop{\mathcal{H}\!\mathit{om}}\nolimits}
\def\SheafExt{\mathop{\mathcal{E}\!\mathit{xt}}\nolimits}
\def\Sch{\mathit{Sch}}
\def\Mor{\mathop{\mathrm{Mor}}\nolimits}
\def\Ob{\mathop{\mathrm{Ob}}\nolimits}
\def\Sh{\mathop{\mathit{Sh}}\nolimits}
\def\NL{\mathop{N\!L}\nolimits}
\def\CH{\mathop{\mathrm{CH}}\nolimits}
\def\proetale{{pro\text{-}\acute{e}tale}}
\def\etale{{\acute{e}tale}}
\def\QCoh{\mathit{QCoh}}
\def\Ker{\mathop{\mathrm{Ker}}}
\def\Im{\mathop{\mathrm{Im}}}
\def\Coker{\mathop{\mathrm{Coker}}}
\def\Coim{\mathop{\mathrm{Coim}}}

% Boxtimes
%
\DeclareMathSymbol{\boxtimes}{\mathbin}{AMSa}{"02}

%
% Macros for moduli stacks/spaces
%
\def\QCohstack{\mathcal{QC}\!\mathit{oh}}
\def\Cohstack{\mathcal{C}\!\mathit{oh}}
\def\Spacesstack{\mathcal{S}\!\mathit{paces}}
\def\Quotfunctor{\mathrm{Quot}}
\def\Hilbfunctor{\mathrm{Hilb}}
\def\Curvesstack{\mathcal{C}\!\mathit{urves}}
\def\Polarizedstack{\mathcal{P}\!\mathit{olarized}}
\def\Complexesstack{\mathcal{C}\!\mathit{omplexes}}
% \Pic is the operator that assigns to X its picard group, usage \Pic(X)
% \Picardstack_{X/B} denotes the Picard stack of X over B
% \Picardfunctor_{X/B} denotes the Picard functor of X over B
\def\Pic{\mathop{\mathrm{Pic}}\nolimits}
\def\Picardstack{\mathcal{P}\!\mathit{ic}}
\def\Picardfunctor{\mathrm{Pic}}
\def\Deformationcategory{\mathcal{D}\!\mathit{ef}}


% OK, start here.
%
\begin{document}

\title{代数空间的上同调}

\maketitle

\phantomsection
\label{section-phantom}

\tableofcontents




\section{引言}
\label{section-introduction}

\noindent
本章讨论代数空间的上同调。虽然我们会证明一些关于阿贝尔层上同调的结果，
但重点是拟凝聚层的上同调；换言之，我们将证明“概形的上同调”一章中
若干结果的类似命题。本章的一些结果可见 \cite{Kn}。

\medskip\noindent
本章缺少的一项重要工具是{\it 归纳原理}，也就是“概形的上同调”中
引理 \ref{coherent-lemma-induction-principle} 对拟紧拟分离代数空间的类似形式。
“空间的导出范畴”第 \ref{spaces-perfect-section-induction} 节对它作了精确表述并给出了详细证明。
本章不用归纳原理，而采用交错 {\v C}ech 复形，见
第 \ref{section-alternating-cech} 节。这个复形旨在证明
命题 \ref{proposition-vanishing} 一类的消失性命题；不过在某些情形下，
归纳原理更为有力，而且或许也更为“标准”。我们建议读者在读完本节的
部分内容之后，再去了解这一归纳原理。



\section{约定}
\label{section-conventions}

\noindent
我们始终假设所有概形都包含在一个大 fppf 位点 $\Sch_{fppf}$ 中，
并且所考虑的每个环 $A$ 都满足：$\Spec(A)$（同构于）这个大位点的一个对象。

\medskip\noindent
设 $S$ 为概形，$X$ 为 $S$ 上的代数空间。在本章及后续各章中，
我们把 $X$ 与自身的积（取于 $S$ 上代数空间的范畴中）记作
$X \times_S X$，而不记作 $X \times X$。








\section{高阶正像}
\label{section-higher-direct-image}

\noindent
设 $S$ 为概形，$X$ 为 $S$ 上的可表代数空间，并设 $\mathcal{F}$ 为 $X$ 上的
拟凝聚模（见“代数空间的性质”第
\ref{spaces-properties-section-quasi-coherent} 节）。由“下降”命题
\ref{descent-proposition-same-cohomology-quasi-coherent}，上同调群
$H^i(X, \mathcal{F})$ 与通常的上同调群一致；后者是在表示 $X$ 的概形上，
对相应拟凝聚模按 Zariski 拓扑计算得到的。

\medskip\noindent
更一般地，设 $f : X \to Y$ 是可表代数空间 $X$ 与 $Y$ 之间的拟紧拟分离态射，
$\mathcal{F}$ 是 $X$ 上的拟凝聚模。由“下降”引理
\ref{descent-lemma-higher-direct-images-small-etale}，层
$R^if_*\mathcal{F}$ 与通常的高阶正像一致；后者是在表示 $X$ 的概形到表示
$Y$ 的概形的态射上，对拟凝聚模按 Zariski 拓扑计算得到的。

\medskip\noindent
再一般些，设 $f : X \to Y$ 是 $S$ 上代数空间之间的可表、拟紧且拟分离的态射。
设 $V$ 为概形，$V \to Y$ 为平展满态射。令 $U = V \times_Y X$，并以
$f' : U \to V$ 表示 $f$ 的基变换。于是对任意拟凝聚
$\mathcal{O}_X$-模 $\mathcal{F}$，有
\begin{equation}
\label{equation-representable-higher-direct-image}
R^if'_*(\mathcal{F}|_U) = (R^if_*\mathcal{F})|_V,
\end{equation}
见“代数空间的性质”引理
\ref{spaces-properties-lemma-pushforward-etale-base-change-modules}。
又因为 $f' : U \to V$ 是概形之间的拟紧拟分离态射，根据上一段的说明，
我们可以把 $\mathcal{F}|_U$ 看成概形 $U$ 上的拟凝聚层，把 $f'$ 看成概形之间
的态射，从而计算 $R^if'_*(\mathcal{F}|_U)$。下文将反复使用这一事实而不再说明。

\medskip\noindent
下面证明：对于代数空间的任意拟紧拟分离态射，拟凝聚层的高阶正像仍为拟凝聚层。
证明中会用到一个技巧；另一种“更好”的证明会使用相对 {\v C}ech 复形，
参见“代数栈上的层”第 \ref{stacks-sheaves-section-cech} 节、
第 \ref{stacks-sheaves-section-sheaf-cech-complex} 节及其后各节。

\begin{lemma}
\label{lemma-higher-direct-image}
设 $S$ 为概形，$f : X \to Y$ 为 $S$ 上代数空间的态射。若 $f$ 拟紧且拟分离，
则 $R^if_*$ 把拟凝聚 $\mathcal{O}_X$-模变为拟凝聚 $\mathcal{O}_Y$-模。
\end{lemma}

\begin{proof}
取平展态射 $V \to Y$，其中 $V$ 为仿射概形。令 $U = V \times_Y X$，并以
$f' : U \to V$ 表示诱导态射。设 $\mathcal{F}$ 为拟凝聚
$\mathcal{O}_X$-模。由“代数空间的性质”引理
\ref{spaces-properties-lemma-pushforward-etale-base-change-modules}，有
$R^if'_*(\mathcal{F}|_U) = (R^if_*\mathcal{F})|_V$.
由于“是拟凝聚模”这一性质对于 $Y$ 上的平展拓扑是局部的（见“代数空间的性质”
引理 \ref{spaces-properties-lemma-characterize-quasi-coherent}），
我们可以用 $V$ 代替 $Y$；也就是说，可以假设 $Y$ 是仿射概形。

\medskip\noindent
假设 $Y$ 仿射。由于 $f$ 拟紧，$X$ 也是拟紧的。因此可以选取一个仿射概形
$U$ 以及一个平展满态射 $g : U \to X$，见“代数空间的性质”引理
\ref{spaces-properties-lemma-quasi-compact-affine-cover}。图示如下：
$$
\xymatrix{
U \ar[r]_g \ar[rd]_{f \circ g} & X \ar[d]^f \\
& Y
}
$$
因为 $X$ 拟分离，态射 $g : U \to X$ 可表、分离且拟紧。因此本引理对 $g$
成立（由引理前的讨论）。它对 $f \circ g : U \to Y$ 也成立，因为这是仿射概形
之间的态射。

\medskip\noindent
在上一段所述情形中，我们将对 $n$ 作归纳，证明命题 $IH_n$：对 $X$ 上任意
拟凝聚层 $\mathcal{F}$，当 $i \leq n$ 时，层 $R^if\mathcal{F}$ 拟凝聚。
$n = 0$ 的情形由“代数空间的态射”引理
\ref{spaces-morphisms-lemma-pushforward} 得到。假设 $IH_n$ 成立；以下证明
$IH_{n + 1}$。

\medskip\noindent
设 $\mathcal{H}$ 为拟凝聚 $\mathcal{O}_U$-模。考虑 Leray 谱序列
$$
E_2^{p, q} = R^pf_* R^qg_* \mathcal{H}
\Rightarrow
R^{p + q}(f \circ g)_*\mathcal{H}
$$
见“位点上的上同调”引理 \ref{sites-cohomology-lemma-relative-Leray}。
由 $IH_n$，$R^qg_*\mathcal{H}$ 拟凝聚，因而当 $p \leq n$ 时，所有层
$R^pf_*R^qg_*\mathcal{H}$ 都拟凝聚。各层
$R^{p + q}(f \circ g)_*\mathcal{H}$ 也都拟凝聚（事实上，当
$p + q > 0$ 时它们为零，但这里不需要这一点）。在总次数不超过
$n + 1$ 的项中，唯一尚不知是否拟凝聚的模是
$E_2^{n + 1, 0} = R^{n + 1}f_*g_*\mathcal{H}$。此外，微分
$d_r^{n + 1, 0} : E_r^{n + 1, 0} \to E_r^{n + 1 + r, 1 - r}$
为零，因为其目标为零。再利用 $\QCoh(\mathcal{O}_X)$ 是
$\textit{Mod}(\mathcal{O}_X)$ 的弱 Serre 子范畴（“代数空间的性质”引理
\ref{spaces-properties-lemma-properties-quasi-coherent}），可知
$R^{n + 1}f_*g_*\mathcal{H}$ 拟凝聚（细节从略）。

\medskip\noindent
设 $\mathcal{F}$ 为拟凝聚 $\mathcal{O}_X$-模，并令
$\mathcal{H} = g^*\mathcal{F}$。由于 $U \to X$ 是平展满态射，伴随映射
$\mathcal{F} \to g_*g^*\mathcal{F} = g_*\mathcal{H}$ 为单射。考虑正合列
$$
0 \to \mathcal{F} \to g_*\mathcal{H} \to \mathcal{G} \to 0
$$
其中 $\mathcal{G}$ 是第一个映射的余核，因而是拟凝聚的。应用上同调长正合列，得到
$$
R^nf_*g_*\mathcal{H} \to
R^nf_*\mathcal{G} \to
R^{n + 1}f_*\mathcal{F} \to
R^{n + 1}f_*g_*\mathcal{H} \to
R^{n + 1}f_*\mathcal{G}
$$
第一个箭头的余核拟凝聚，而上面已经证明
$R^{n + 1}f_*g_*\mathcal{H}$ 拟凝聚。因此 $R^{n + 1}f_*\mathcal{F}$
有一个两步滤过，其中第一步拟凝聚，第二步是某个拟凝聚层的子模。
由于 $\mathcal{F}$ 是任意拟凝聚 $\mathcal{O}_X$-模，这个结论也适用于
$\mathcal{G}$。所以可以选取正合列
$0 \to \mathcal{A} \to R^{n + 1}f_*\mathcal{G} \to \mathcal{B}$
其中 $\mathcal{A}$、$\mathcal{B}$ 都是拟凝聚 $\mathcal{O}_Y$-模。
于是复合映射
$R^{n + 1}f_*g_*\mathcal{H} \to R^{n + 1}f_*\mathcal{G}
\to \mathcal{B}$ 的核 $\mathcal{K}$ 拟凝聚，并得到映射
$\mathcal{K} \to \mathcal{A}$；其核 $\mathcal{K}'$ 也拟凝聚。因此
$R^{n + 1}f_*\mathcal{F}$ 位于正合列
$$
R^nf_*g_*\mathcal{H} \to
R^nf_*\mathcal{G} \to
R^{n + 1}f_*\mathcal{F} \to \mathcal{K}' \to 0
$$
中；其中除 $R^{n + 1}f_*\mathcal{F}$ 外，其余各模均拟凝聚。由此得出
$R^{n + 1}f_*\mathcal{F}$ 拟凝聚，即 $IH_{n + 1}$ 如期成立。
\end{proof}

\begin{lemma}
\label{lemma-quasi-coherence-higher-direct-images-application}
设 $S$ 为概形，$f : X \to Y$ 为 $S$ 上代数空间的拟分离拟紧态射。
对任意拟凝聚 $\mathcal{O}_X$-模 $\mathcal{F}$ 以及 $Y_\etale$ 的任意仿射对象
$V$，有
$$
H^q(V \times_Y X, \mathcal{F}) = H^0(V, R^qf_*\mathcal{F})
$$
其中 $q \in \mathbf{Z}$ 任意。
\end{lemma}

\begin{proof}
由于 $Rf_*$ 的形成与平展局部化可交换（“代数空间的性质”引理
\ref{spaces-properties-lemma-pushforward-etale-base-change-modules}），
我们可以用 $V$ 代替 $Y$，并假设 $Y = V$ 仿射。考虑 Leray 谱序列
$E_2^{p, q} = H^p(Y, R^qf_*\mathcal{F})$
它收敛到 $H^{p + q}(X, \mathcal{F})$，见“位点上的上同调”引理
\ref{sites-cohomology-lemma-Leray}。由引理
\ref{lemma-higher-direct-image}，各层 $R^qf_*\mathcal{F}$ 均拟凝聚。
再由“概形的上同调”引理
\ref{coherent-lemma-quasi-coherent-affine-cohomology-zero}，当 $p > 0$ 时
$E_2^{p, q} = 0$。所以该谱序列在 $E_2$ 页退化，结论成立。
\end{proof}




\section{有限态射}
\label{section-finite-morphisms}

\noindent
下面给出若干对所有阿贝尔层都成立的结果（特别地，也适用于拟凝聚模）。
我们要{\bf 提醒}读者：对于概形的有限态射和 Zariski 拓扑，以下引理并不成立。

\begin{lemma}
\label{lemma-finite-higher-direct-image-zero}
设 $S$ 为概形，$f : X \to Y$ 为代数空间的整态射（例如有限态射）。则
$f_* : \textit{Ab}(X_\etale) \to \textit{Ab}(Y_\etale)$
是正合函子，并且当 $p > 0$ 时 $R^pf_* = 0$。
\end{lemma}

\begin{proof}
由“代数空间的性质”引理
\ref{spaces-properties-lemma-pushforward-etale-base-change}，
可以在 $Y$ 的一个平展覆盖上计算高阶正像。因此可以假设 $Y$ 是概形，
进而 $X$ 也是概形（见“代数空间的态射”引理
\ref{spaces-morphisms-lemma-integral-local}）。此时可应用“平展上同调”引理
\ref{etale-cohomology-lemma-what-integral}。对于有限情形，读者也可以参阅
技术性较弱的“平展上同调”命题
\ref{etale-cohomology-proposition-finite-higher-direct-image-zero}。
\end{proof}

\begin{lemma}
\label{lemma-stalk-push-finite}
设 $S$ 为概形，$f : X \to Y$ 为 $S$ 上代数空间的有限态射。设
$\overline{y}$ 是 $Y$ 的一个几何点，并且它在 $X$ 中的提升为
$\overline{x}_1, \ldots, \overline{x}_n$。则
$$
(f_*\mathcal{F})_{\overline{y}} =
\prod\nolimits_{i = 1, \ldots, n}
\mathcal{F}_{\overline{x}_i}
$$
对 $X_\etale$ 上任意层 $\mathcal{F}$ 都成立。
\end{lemma}

\begin{proof}
选取 $\overline{y}$ 的一个平展邻域 $(V, \overline{v})$。于是茎
$(f_*\mathcal{F})_{\overline{y}}$ 就是 $f_*\mathcal{F}|_V$ 在
$\overline{v}$ 处的茎。由“代数空间的性质”引理
\ref{spaces-properties-lemma-pushforward-etale-base-change}，
可以用 $V$ 代替 $Y$，并用 $X \times_Y V$ 代替 $X$。
% P09-ERR-0001: frozen source literally has Z -> X here; see corrections.md.
于是 $Z \to X$ 是概形之间的有限态射，而所需结论就是“平展上同调”命题
\ref{etale-cohomology-proposition-finite-higher-direct-image-zero}。
\end{proof}

\begin{lemma}
\label{lemma-finite-rings}
设 $S$ 为概形，$\pi : X \to Y$ 为 $S$ 上代数空间的有限态射。设
$\mathcal{A}$ 是 $X_\etale$ 上的环层，$\mathcal{B}$ 是 $Y_\etale$ 上的环层，
并设 $\varphi : \mathcal{B} \to \pi_*\mathcal{A}$ 为环层同态，从而得到
带环拓扑斯的态射
$$
f = (\pi, \varphi) :
(\Sh(X_\etale), \mathcal{A})
\longrightarrow
(\Sh(Y_\etale), \mathcal{B}).
$$
若 $\mathcal{F}$ 是 $\mathcal{A}$-模层，$\mathcal{G}$ 是 $\mathcal{B}$-模层，
则典范映射
$$
\mathcal{G} \otimes_\mathcal{B} f_*\mathcal{F}
\longrightarrow
f_*(f^*\mathcal{G} \otimes_\mathcal{A} \mathcal{F}).
$$
是同构。
\end{lemma}

\begin{proof}
这个映射是下列映射的伴随映射：
$$
f^*\mathcal{G} \otimes_\mathcal{A}
f^* f_*\mathcal{F} =
f^*(\mathcal{G} \otimes_\mathcal{B} f_*\mathcal{F})
\longrightarrow
f^*\mathcal{G} \otimes_\mathcal{A} \mathcal{F}
$$
它由 $\text{id} : f^*\mathcal{G} \to f^*\mathcal{G}$ 和伴随映射
$f^* f_*\mathcal{F} \to \mathcal{F}$ 给出。要证明它是同构，只需在茎上检验
（见“代数空间的性质”定理
\ref{spaces-properties-theorem-exactness-stalks}）。设 $\overline{y}$ 为 $Y$
的几何点，$\overline{x}_1, \ldots, \overline{x}_n$ 为位于
$\overline{y}$ 上方的 $X$ 的几何点。把这些映射在茎上的作用写出后，只需证明
$$
\mathcal{G}_{\overline{y}}
\otimes_{\mathcal{B}_{\overline{y}}}
\left(
\bigoplus\nolimits_{i = 1, \ldots, n} \mathcal{F}_{\overline{x}_i}
\right) =
\bigoplus\nolimits_{i = 1, \ldots, n}
% P09-ERR-0002: frozen source literally has B_{xbar}; see corrections.md.
(\mathcal{G}_{\overline{y}}
\otimes_{\mathcal{B}_{\overline{x}}}
\mathcal{A}_{\overline{x}_i}) \otimes_{\mathcal{A}_{\overline{x}_i}}
\mathcal{F}_{\overline{x}_i}
$$
确实成立。这里使用了张量积与取茎可交换，以及“代数空间的性质”引理
\ref{spaces-properties-lemma-stalk-pullback} 所述拉回下茎的行为；还使用了
% P09-ERR-0003: frozen source says closed immersion; see corrections.md.
引理 \ref{lemma-stalk-push-finite} 所述沿闭浸入作正像时茎的行为。
\end{proof}

\noindent
最后给出一个对有限态射而言极其一般的投影公式。

\begin{lemma}
\label{lemma-projection-formula-finite}
沿用引理 \ref{lemma-finite-rings} 中的
$S$、$X$、$Y$、$\pi$、$\mathcal{A}$、$\mathcal{B}$、$\varphi$ 和 $f$，则
$$
K \otimes_\mathcal{B}^\mathbf{L} Rf_*M =
Rf_*(Lf^*K \otimes_\mathcal{A}^\mathbf{L} M)
$$
对任意 $K \in D(\mathcal{B})$ 和 $M \in D(\mathcal{A})$，上述等式在
$D(\mathcal{B})$ 中成立。
\end{lemma}

\begin{proof}
由于 $f_*$ 正合（引理 \ref{lemma-finite-higher-direct-image-zero}），
对任意代表复形逐项应用 $f_*$ 即可计算函子 $Rf_*$。选取一个代表 $K$ 的
$\mathcal{B}$-模复形 $\mathcal{K}^\bullet$，使其 K-平坦且各项平坦；见
“位点上的上同调”引理 \ref{sites-cohomology-lemma-K-flat-resolution}。
于是 $f^*\mathcal{K}^\bullet$ 也 K-平坦且各项平坦；见“位点上的上同调”引理
\ref{sites-cohomology-lemma-pullback-K-flat}。任取一个代表 $M$ 的
$\mathcal{A}$-模复形 $\mathcal{M}^\bullet$。现在只需证明
$$
\text{Tot}(\mathcal{K}^\bullet \otimes_\mathcal{B} f_*\mathcal{M}^\bullet)
=
f_*\text{Tot}(f^*\mathcal{K}^\bullet \otimes_\mathcal{A} \mathcal{M}^\bullet)
$$
因为按照上述选择，这两个复形分别代表引理中公式的右端和左端。
又因为 $f_*$ 与直和可交换（例如由引理
\ref{lemma-stalk-push-finite} 中对茎的描述可知），问题归结为等式
$$
\mathcal{K}^n \otimes_\mathcal{B} f_*\mathcal{M}^m
=
f_*(f^*\mathcal{K}^n \otimes_\mathcal{A} \mathcal{M}^m)
$$
而这些等式由引理 \ref{lemma-finite-rings} 成立。
\end{proof}









\section{余极限与上同调}
\label{section-colimits}

\noindent
下面的引理特别适用于拟凝聚层的图式。

\begin{lemma}
\label{lemma-colimits}
设 $S$ 为概形，$X$ 为 $S$ 上的代数空间。若 $X$ 拟紧且拟分离，则
$$
\colim_i H^p(X, \mathcal{F}_i)
\longrightarrow
H^p(X, \colim_i \mathcal{F}_i)
$$
对 $X_\etale$ 上阿贝尔层的每个滤过图式都是同构。
\end{lemma}

\begin{proof}
这由“位点上的上同调”引理
\ref{sites-cohomology-lemma-colim-works-over-collection} 得到。具体而言，
令 $\mathcal{B} \subset \Ob(X_{spaces, \etale})$ 为所有在 $X$ 上平展、
拟紧且拟分离的空间组成的集合。注意，若 $U \in \mathcal{B}$，则由于 $U$
拟紧，由 $U_i \in \mathcal{B}$ 构成的有限覆盖 $\{U_i \to U\}$ 的全体，
在 $X_{spaces, \etale}$ 中 $U$ 的所有覆盖组成的集合里是共尾的。由
“代数空间的态射”引理
\ref{spaces-morphisms-lemma-quasi-compact-quasi-separated-permanence}，
集合 $\mathcal{B}$ 满足“位点上的上同调”引理
\ref{sites-cohomology-lemma-colim-works-over-collection} 的全部假设。
因为 $X \in \mathcal{B}$，结论成立。
\end{proof}

\begin{lemma}
\label{lemma-colimit-cohomology}
\begin{slogan}
拟紧拟分离态射的高阶正像与层的滤过余极限可交换。
\end{slogan}
设 $S$ 为概形，$f : X \to Y$ 为 $S$ 上代数空间的拟紧拟分离态射。设
$\mathcal{F} = \colim \mathcal{F}_i$ 是 $X_\etale$ 上阿贝尔层的滤过余极限。
则对任意 $p \geq 0$，有
$$
R^pf_*\mathcal{F} = \colim R^pf_*\mathcal{F}_i.
$$
\end{lemma}

\begin{proof}
我们要用到如下事实：拓扑斯态射 $f_{small} : X_{small} \to Y_{small}$
来自位点态射
$f_{spaces, \etale} : X_{spaces, \etale} \to Y_{spaces, \etale}$
它对应于连续函子 $V \longmapsto X \times_Y V$；见“代数空间的性质”引理
\ref{spaces-properties-lemma-functoriality-etale-site}。我们将对这个位点态射
应用“位点上的上同调”引理
\ref{sites-cohomology-lemma-higher-direct-image-colimit}。由于
$Y_{spaces, \etale}$ 的每个对象都有一个由仿射对象组成的覆盖，只需证明：
当 $V$ 仿射且在 $Y$ 上平展时，有
$H^p(X \times_Y V, \mathcal{F}) = \colim H^p(X \times_Y V, \mathcal{F}_i)$.
由于 $V$ 仿射，代数空间 $X \times_Y V$ 拟紧且拟分离。因此应用引理
\ref{lemma-colimits} 即得结论。
\end{proof}

\noindent
下面的引理说明，有限表示模在拟紧拟分离代数空间上具有预期的性质。

\begin{lemma}
\label{lemma-finite-presentation-quasi-compact-colimit}
设 $S$ 为概形，$X$ 为 $S$ 上拟紧拟分离的代数空间。设 $I$ 为有向集，
$(\mathcal{F}_i, \varphi_{ii'})$ 为由 $I$ 索引的 $\mathcal{O}_X$-模系统，
并设 $\mathcal{G}$ 为有限表示的 $\mathcal{O}_X$-模。则
$$
\colim_i \Hom_X(\mathcal{G}, \mathcal{F}_i)
=
\Hom_X(\mathcal{G}, \colim_i \mathcal{F}_i).
$$
特别地，$\Hom_X(\mathcal{G}, -)$ 在 $\QCoh(\mathcal{O}_X)$ 中与滤过余极限可交换。
\end{lemma}

\begin{proof}
所示等式是“位点上的模”引理
\ref{sites-modules-lemma-finite-presentation-quasi-compact-colimit}
的一个特例。为应用该引理，需要对位点 $X_\etale$ 验证“位点”引理
\ref{sites-lemma-directed-colimits-global-sections} 第 (4) 部分的假设。
为此，我们验证“位点”注记 \ref{sites-remark-stronger-conditions}
中的假设 (2)(a)、(2)(b)、(2)(c)。令
$\mathcal{B} \subset \Ob(X_\etale)$ 为仿射对象的集合，则：
\begin{enumerate}
\item 由于 $X$ 拟紧，存在 $U \in \mathcal{B}$，使 $U \to X$ 为满射
（见“代数空间的性质”引理
\ref{spaces-properties-lemma-quasi-compact-affine-cover}），
所以 $h_U^\# \to *$ 为满射。
\item 对 $U \in \mathcal{B}$，$U$ 的每个平展覆盖
$\{U_i \to U\}_{i \in I}$ 都可由一个有限平展覆盖
$\{U_j \to U\}_{j = 1, \ldots, m}$ 加细，其中 $U_j \in \mathcal{B}$
（见“拓扑”引理 \ref{topologies-lemma-etale-affine}）。
\item 对 $U, U' \in \Ob(X_\etale)$，有
$h_U^\# \times h_{U'}^\# = h_{U \times_X U'}^\#$.
若 $U, U' \in \mathcal{B}$，则因为 $X$ 拟分离，$U \times_X U'$ 拟紧；
例如见“代数空间的态射”引理
\ref{spaces-morphisms-lemma-quasi-compact-quasi-separated-permanence}。
因此可找到平展满态射 $U'' \to U \times_X U'$，其中
$U'' \in \mathcal{B}$（见“代数空间的性质”引理
\ref{spaces-properties-lemma-quasi-compact-affine-cover}）。换言之，存在态射
$U'' \to U$ 和 $U'' \to U'$，使得
% P09-ERR-0004: frozen source literally has lower-case u'; see corrections.md.
$h_{U''}^\# \to h_U^\# \times h_{u'}^\#$ 为满射。
\end{enumerate}
关于最后一个断言，注意包含函子
$\QCoh(\mathcal{O}_X) \to \textit{Mod}(\mathcal{O}_X)$ 与余极限可交换，
且有限表示模拟凝聚。参见“代数空间的性质”引理
\ref{spaces-properties-lemma-properties-quasi-coherent}。
\end{proof}





\section{交错 {\v C}ech 复形}
\label{section-alternating-cech}

\noindent
设 $S$ 为概形，$f : U \to X$ 为 $S$ 上代数空间的平展态射。函子
$$
j : U_{spaces, \etale} \longrightarrow X_{spaces, \etale},\quad
V/U \longmapsto V/X
$$
诱导出 $U_{spaces, \etale}$ 与局部化 $X_{spaces, \etale}/U$ 之间的等价；
见“代数空间的性质”第 \ref{spaces-properties-section-localize} 节。因此存在函子
$$
f_! : \textit{Ab}(U_\etale) \longrightarrow
\textit{Ab}(X_\etale),\quad
f_! : \textit{Mod}(\mathcal{O}_U) \longrightarrow \textit{Mod}(\mathcal{O}_X),
$$
它们分别左伴随于
$$
f^{-1} : \textit{Ab}(X_\etale) \longrightarrow
\textit{Ab}(U_\etale),\quad
f^* : \textit{Mod}(\mathcal{O}_X) \longrightarrow \textit{Mod}(\mathcal{O}_U)
$$
见“位点上的模”第 \ref{sites-modules-section-localize} 节。
注意：先验地说，这个函子与紧支撑上同调毫无关系！在上述文献中，
我们把它称为“零延拓”函子。还要注意，对 $\mathcal{O}_X$-模层有
$f^* = f^{-1}$，所以 $f_!$ 的两个版本彼此一致。

\medskip\noindent
下面会用到这个构造，先回顾它的若干性质。给定 $U_\etale$ 上的阿贝尔层
$\mathcal{G}$，
% P09-ERR-0005: frozen source literally omits G after f_!; see ERRATA_CANDIDATES.jsonl.
层 $f_!$ 是下列预层的层化：
$$
V/X \longmapsto
f_!\mathcal{G}(V) =
\bigoplus\nolimits_{\varphi \in \Mor_X(V, U)}
\mathcal{G}(V \xrightarrow{\varphi} U),
$$
见“位点上的模”引理 \ref{sites-modules-lemma-extension-by-zero}。
此外，若 $\mathcal{G}$ 是 $\mathcal{O}_U$-模，则 $f_!\mathcal{G}$ 是同一个
阿贝尔群预层的层化，而该预层以显然的方式带有 $\mathcal{O}_X$-模结构
（见同处）。设 $\overline{x} : \Spec(k) \to X$ 为几何点，则有典范等同
$$
(f_!\mathcal{G})_{\overline{x}} =
\bigoplus\nolimits_{\overline{u}} \mathcal{G}_{\overline{u}}
$$
其中直和遍历所有满足 $f \circ \overline{u} = \overline{x}$ 的
$\overline{u} : \Spec(k) \to U$；见“位点上的模”引理
\ref{sites-modules-lemma-stalk-j-shriek} 和“代数空间的性质”引理
\ref{spaces-properties-lemma-points-small-etale-site}。下文将研究层
$f_!\underline{\mathbf{Z}}$。这里 $\underline{\mathbf{Z}}$ 表示
$X_\etale$ 或 $U_\etale$ 上的常值层。

\begin{lemma}
\label{lemma-product-is-tensor-product}
设 $S$ 为概形，$f_i : U_i \to X$ 为 $S$ 上代数空间的平展态射。则有同构
$$
f_{1, !}\underline{\mathbf{Z}} \otimes_{\mathbf{Z}}
f_{2, !}\underline{\mathbf{Z}}
\longrightarrow
f_{12, !}\underline{\mathbf{Z}}
$$
其中 $f_{12} : U_1 \times_X U_2 \to X$ 是结构态射；并且有
$$
(f_1 \amalg f_2)_! \underline{\mathbf{Z}}
\longrightarrow
f_{1, !}\underline{\mathbf{Z}} \oplus
f_{2, !}\underline{\mathbf{Z}}
$$
\end{lemma}

\begin{proof}
一旦定义了这个映射，由上面对茎的描述可知它是同构。定义该映射时，
只需在预层层次上工作。因此需要定义映射
$$
\left(\bigoplus\nolimits_{\varphi_1 \in \Mor_X(V, U_1)} \mathbf{Z}\right)
\otimes_{\mathbf{Z}}
\left(\bigoplus\nolimits_{\varphi_2 \in \Mor_X(V, U_2)} \mathbf{Z}\right)
\longrightarrow
\bigoplus\nolimits_{\varphi \in \Mor_X(V, U_1 \times_X U_2)}
\mathbf{Z}
$$
使用显然的记号，把元素 $1_{\varphi_1} \otimes 1_{\varphi_2}$ 映到元素
$1_{\varphi_1 \times \varphi_2}$。第二个等式的证明从略。
\end{proof}

\noindent
另一个重要性质是迹映射
$$
\text{Tr}_f : f_!\underline{\mathbf{Z}} \longrightarrow \underline{\mathbf{Z}}.
$$
迹映射伴随于映射 $\mathbf{Z} \to f^{-1}\underline{\mathbf{Z}}$
（后者是同构）。若 $\overline{x}$ 如上，则 $\text{Tr}_f$ 在
$\overline{x}$ 处的茎上是映射
$$
(\text{Tr}_f)_{\overline{x}} :
(f_!\underline{\mathbf{Z}})_{\overline{x}} =
\bigoplus\nolimits_{\overline{u}} \mathbf{Z}
\longrightarrow
\mathbf{Z} = \underline{\mathbf{Z}}_{\overline{x}}
$$
它把给定的整数相加。这是因为它伴随于映射
$1 : \mathbf{Z} \to f^{-1}\underline{\mathbf{Z}}$。特别地，若 $f$ 既平展又
为满射，则 $\text{Tr}_f$ 也是满射。

\medskip\noindent
假设 $f : U \to X$ 是代数空间的平展满态射。考虑与上述迹映射相伴的
{\it Koszul 复形}
$$
\ldots \to \wedge^3f_!\underline{\mathbf{Z}} \to
\wedge^2f_!\underline{\mathbf{Z}} \to f_!\underline{\mathbf{Z}} \to
\underline{\mathbf{Z}} \to 0
$$
这里的外幂取在环层 $\underline{\mathbf{Z}}$ 上。映射由规则
$$
e_1 \wedge \ldots \wedge e_n \longmapsto
\sum\nolimits_{i = 1, \ldots, n} (-1)^{i + 1}
\text{Tr}_f(e_i)
e_1 \wedge \ldots \wedge \widehat{e_i} \wedge \ldots \wedge e_n
$$
定义，其中 $e_1, \ldots, e_n$ 是 $f_!\underline{\mathbf{Z}}$ 的局部截面。
设 $\overline{x}$ 为 $X$ 的几何点，并令
$M_{\overline{x}} = (f_!\underline{\mathbf{Z}})_{\overline{x}} =
\bigoplus_{\overline{u}} \mathbf{Z}$。则上述复形在 $\overline{x}$ 处的茎为复形
$$
\ldots \to \wedge^3 M_{\overline{x}} \to \wedge^2 M_{\overline{x}}
\to M_{\overline{x}} \to \mathbf{Z} \to 0
$$
它是正合的，因为 $M_{\overline{x}} \to \mathbf{Z}$ 为满射；见
“关于代数的更多内容”引理
\ref{more-algebra-lemma-homotopy-koszul-abstract}。因此，若以
$K^\bullet = K^\bullet(f)$ 表示满足
$K^i = \wedge^{i + 1}f_!\underline{\mathbf{Z}}$ 的复形，就得到拟同构
\begin{equation}
\label{equation-quasi-isomorphism}
K^\bullet \longrightarrow \underline{\mathbf{Z}}[0]
\end{equation}
我们用复形 $K^\bullet$ 定义与 $f : U \to X$ 相伴的所谓交错 {\v C}ech 复形。

\begin{definition}
\label{definition-alternating-cech-complex}
设 $S$ 为概形，$f : U \to X$ 为 $S$ 上代数空间的平展满态射，并设
$\mathcal{F}$ 为 $\textit{Ab}(X_\etale)$ 的对象。与 $\mathcal{F}$ 和 $f$
相伴的{\it 交错 {\v C}ech 复形}\footnote{这可能不是标准记号}
$\check{\mathcal{C}}^\bullet_{alt}(f, \mathcal{F})$ 是复形
$$
\Hom(K^0, \mathcal{F}) \to \Hom(K^1, \mathcal{F}) \to
\Hom(K^2, \mathcal{F}) \to \ldots
$$
其中 Hom 群在 $\textit{Ab}(X_\etale)$ 中计算。
\end{definition}

\noindent
读者可以验证：若 $U = \coprod U_i$，且 $f|_{U_i} : U_i \to X$ 是某个子空间
的开浸入，则 $\check{\mathcal{C}}_{alt}^\bullet(f, \mathcal{F})$ 与“上同调”
第 \ref{cohomology-section-alternating-cech} 节针对 Zariski 覆盖
$X = \bigcup U_i$ 及 $\mathcal{F}$ 在 $X$ 的 Zariski 位点上的限制所引入的
复形一致。不过更重要的是，把交错 {\v C}ech 复形的上同调与层上同调联系起来。

\begin{lemma}
\label{lemma-alternating-cech-to-cohomology}
设 $S$ 为概形，$f : U \to X$ 为 $S$ 上代数空间的平展满态射，并设
$\mathcal{F}$ 为 $\textit{Ab}(X_\etale)$ 的对象。则在 $D(\textit{Ab})$ 中
存在典范映射
$$
\check{\mathcal{C}}^\bullet_{alt}(f, \mathcal{F})
\longrightarrow
R\Gamma(X, \mathcal{F})
$$
此外，存在一个 $E_1$ 页为
$$
E_1^{p, q} =
\Ext_{\textit{Ab}(X_\etale)}^q(K^p, \mathcal{F})
$$
的谱序列，它收敛到 $H^{p + q}(X, \mathcal{F})$，其中
$K^p = \wedge^{p + 1}f_!\underline{\mathbf{Z}}$。
\end{lemma}

\begin{proof}
回顾拟同构 $K^\bullet \to \underline{\mathbf{Z}}[0]$，见
(\ref{equation-quasi-isomorphism})。在 $\textit{Ab}(X_\etale)$ 中选取单射
消解 $\mathcal{F} \to \mathcal{I}^\bullet$。考虑二重复形
$\Hom(K^\bullet, \mathcal{I}^\bullet)$，其项为
$\Hom(K^p, \mathcal{I}^q)$。微分
$d_1^{p, q} : A^{p, q} \to A^{p + 1, q}$
来自微分 $K^{p + 1} \to K^p$，而微分
$d_2^{p, q} : A^{p, q} \to A^{p, q + 1}$ 来自微分
$\mathcal{I}^q \to \mathcal{I}^{q + 1}$。以
$\text{Tot}(\Hom(K^\bullet, \mathcal{I}^\bullet))$ 表示相应的全复形；见
“同调代数”第 \ref{homology-section-double-complexes} 节。我们将使用与这个
二重复形相伴的两个谱序列 $({}'E_r, {}'d_r)$ 和 $({}''E_r, {}''d_r)$；见
“同调代数”第 \ref{homology-section-double-complex} 节。

\medskip\noindent
因为 $K^\bullet$ 是 $\underline{\mathbf{Z}}$ 的消解，复形
$$
\Hom(K^\bullet, \mathcal{I}^q) :
\Hom(K^0, \mathcal{I}^q) \to
\Hom(K^1, \mathcal{I}^q) \to
\Hom(K^2, \mathcal{I}^q) \to \ldots
$$
在正次数上无上同调，并且其 $H^0$ 等于 $\Gamma(X, \mathcal{I}^q)$。
因此由“同调代数”引理
\ref{homology-lemma-double-complex-gives-resolution}，自然映射
$$
\mathcal{I}^\bullet(X) \longrightarrow
\text{Tot}(\Hom(K^\bullet, \mathcal{I}^\bullet))
$$
是阿贝尔群复形的拟同构。特别地，有
$H^n(\text{Tot}(\Hom(K^\bullet, \mathcal{I}^\bullet))) = H^n(X, \mathcal{F})$.

\medskip\noindent
本引理中的映射 $\check{\mathcal{C}}^\bullet_{alt}(f, \mathcal{F}) \to
R\Gamma(X, \mathcal{F})$，是映射
$\check{\mathcal{C}}^\bullet_{alt}(f, \mathcal{F}) \to
\text{Tot}(\Hom(K^\bullet, \mathcal{I}^\bullet))$ 与上述拟同构之逆的复合。

\medskip\noindent
最后考虑谱序列 $({}'E_r, {}'d_r)$。有
$$
E_1^{p, q} = q\text{th cohomology of }
\Hom(K^p, \mathcal{I}^0) \to
\Hom(K^p, \mathcal{I}^1) \to
\Hom(K^p, \mathcal{I}^2) \to \ldots
$$
这就证明了引理。
\end{proof}

\noindent
由该引理可见，理解 Ext 群
$\Ext_{\textit{Ab}(X_\etale)}(K^p, \mathcal{F})$ 很重要；也就是说，
要理解函子 $\mathcal{F} \mapsto \Hom(K^p, \mathcal{F})$ 的右导出函子。

\begin{lemma}
\label{lemma-compute}
设 $S$ 为概形，$f : U \to X$ 为 $S$ 上代数空间的满、平展且分离的态射。
对 $p \geq 0$，令
$$
W_p = U \times_X \ldots \times_X U \setminus \text{all diagonals}
$$
其中纤维积有 $p + 1$ 个因子。在 $X$ 上，$S_{p + 1}$ 自由作用于 $W_p$，并且
$$
\Hom(K^p, \mathcal{F}) = S_{p + 1}\text{-anti-invariant elements of }
\mathcal{F}(W_p)
$$
这一等式关于 $\mathcal{F}$ 函子性地成立，其中
$K^p = \wedge^{p + 1}f_!\underline{\mathbf{Z}}$。
\end{lemma}

\begin{proof}
因为 $U \to X$ 分离，对角态射 $U \to U \times_X U$ 是闭浸入。又因为
$U \to X$ 平展，该对角态射也是开浸入；见“代数空间的态射”引理
\ref{spaces-morphisms-lemma-etale-unramified} 和
\ref{spaces-morphisms-lemma-diagonal-unramified-morphism}。因此 $W_p$ 是
$U^{p + 1} = U \times_X \ldots \times_X U$ 的既开又闭子空间。我们已经删去
这个作用的所有不动点，所以 $S_{p + 1}$ 在 $W_p$ 上的作用是自由的。
由引理 \ref{lemma-product-is-tensor-product}，有
$$
(f_!\underline{\mathbf{Z}})^{\otimes p + 1} =
f^{p + 1}_!\underline{\mathbf{Z}} = (W_p \to X)_!\underline{\mathbf{Z}}
\oplus Rest
$$
其中 $f^{p + 1} : U^{p + 1} \to X$ 是结构态射。考察 $X$ 的几何点
$\overline{x}$ 上方的茎，可见
$$
\left(
\bigoplus\nolimits_{\overline{u} \mapsto \overline{x}} \mathbf{Z}
\right)^{\otimes p + 1}
\longrightarrow
(W_p \to X)_!\underline{\mathbf{Z}}_{\overline{x}}
$$
是一个商映射，其核由所有满足某个 $i \not = j$ 时
$\overline{u}_i = \overline{u}_j$ 的张量
$1_{\overline{u}_0} \otimes \ldots \otimes 1_{\overline{u}_p}$ 生成。
因此商映射
$$
(f_!\underline{\mathbf{Z}})^{\otimes p + 1}
\longrightarrow
\wedge^{p + 1}f_!\underline{\mathbf{Z}}
$$
经由 $(W_p \to X)_!\underline{\mathbf{Z}}$ 分解；也就是说，得到
$$
(f_!\underline{\mathbf{Z}})^{\otimes p + 1}
\longrightarrow
(W_p \to X)_!\underline{\mathbf{Z}}
\longrightarrow
\wedge^{p + 1}f_!\underline{\mathbf{Z}}
$$
这已经证明 $\Hom(K^p, \mathcal{F})$ 函子性地是下列群的子群：
$$
\Hom((W_p \to X)_!\underline{\mathbf{Z}}, \mathcal{F}) = \mathcal{F}(W_p)
$$
要把它认作 $S_{p + 1}$-反不变元，需要证明满射
$(W_p \to X)_!\underline{\mathbf{Z}}
\to \wedge^{p + 1}f_!\underline{\mathbf{Z}}$ 给出最大的
$S_{p + 1}$-反不变商。换言之，需要证明
$\wedge^{p + 1}f_!\underline{\mathbf{Z}}$ 是
$(W_p \to X)_!\underline{\mathbf{Z}}$ 对下述子层的商：该子层由局部截面
$s - \text{sign}(\sigma)\sigma(s)$ 生成，其中 $s$ 是
$(W_p \to X)_!\underline{\mathbf{Z}}$ 的局部截面。这可以在茎上检验，
而在茎上结论显然。
\end{proof}

\begin{lemma}
\label{lemma-twist}
设 $S$ 为概形，$W$ 为 $S$ 上的代数空间，并设有限群 $G$ 自由作用于 $W$。
令 $U = W/G$，见“代数空间的性质”引理
\ref{spaces-properties-lemma-quotient}。设
$\chi : G \to \{+1, -1\}$ 为特征标。则 $U_\etale$ 上存在秩为 $1$ 的局部自由
$\mathbf{Z}$-模层 $\underline{\mathbf{Z}}(\chi)$，使得对 $U_\etale$ 上任意
阿贝尔层 $\mathcal{F}$ 都有
$$
H^0(W, \mathcal{F}|_W)^\chi =
H^0(U, \mathcal{F} \otimes_{\mathbf{Z}} \underline{\mathbf{Z}}(\chi))
$$
\end{lemma}

\begin{proof}
商态射 $q : W \to U$ 是 $G$-挠子；也就是说，存在平展满态射 $U' \to U$，
使得作为 $U'$ 上带 $G$-作用的空间，有
$W \times_U U' = \coprod_{g \in G} U'$（取 $U' = W$ 即可）。因此
$q_*\underline{\mathbf{Z}}$ 是带 $G$-作用的有限局部自由 $\mathbf{Z}$-模。
对 $U$ 的任意几何点 $\overline{u}$，得到 $G$-等变同构
$$
(q_*\underline{\mathbf{Z}})_{\overline{u}}
= \bigoplus\nolimits_{\overline{w} \mapsto \overline{u}} \mathbf{Z}
= \bigoplus\nolimits_{g \in G} \mathbf{Z} = \mathbf{Z}[G]
$$
第二个等号用到位于 $\overline{u}$ 上方的一个几何点 $\overline{w}_0$，
它把对应于 $g \in G$ 的直和项映到对应于 $g(\overline{w}_0)$ 的直和项。有
$$
H^0(W, \mathcal{F}|_W) =
H^0(U, \mathcal{F} \otimes_\mathbf{Z} q_*\underline{\mathbf{Z}})
$$
这是因为
$q_*\mathcal{F}|_W = \mathcal{F} \otimes_\mathbf{Z} q_*\underline{\mathbf{Z}}$；
限制到 $U'$ 后即可检验。令
$$
\underline{\mathbf{Z}}(\chi) =
(q_*\underline{\mathbf{Z}})^\chi \subset
q_*\underline{\mathbf{Z}}
$$
为按 $\chi$ 变换的截面所成的子层。对 $U$ 的任意几何点 $\overline{u}$，有
$$
\underline{\mathbf{Z}}(\chi)_{\overline{u}} =
\mathbf{Z} \cdot \sum\nolimits_g \chi(g) g
\subset
\mathbf{Z}[G] = (q_*\underline{\mathbf{Z}})_{\overline{u}}
$$
由此可知 $\underline{\mathbf{Z}}(\chi)$ 是秩为 $1$ 的局部自由层
（更确切地说，应限制到 $U'$ 后检验）。注意，对任意 $\mathbf{Z}$-模 $M$，
$M[G]$ 的 $\chi$-半不变元恰为形如
$m \cdot \sum\nolimits_g \chi(g) g$ 的元素。因此，对 $U$ 上任意阿贝尔层
$\mathcal{F}$，有
$$
\left(\mathcal{F} \otimes_\mathbf{Z} q_*\underline{\mathbf{Z}}\right)^\chi
=
\mathcal{F} \otimes_\mathbf{Z} \underline{\mathbf{Z}}(\chi)
$$
因为该等式在所有茎上成立。取整体截面即得引理结论。
\end{proof}

\noindent
现在把上述结果汇集起来，得到下面这个优美的结论。

\begin{lemma}
\label{lemma-alternating-spectral-sequence}
设 $S$ 为概形，$f : U \to X$ 为 $S$ 上代数空间的满、平展且分离的态射。
对 $p \geq 0$，令
$$
W_p = U \times_X \ldots \times_X U \setminus \text{all diagonals}
$$
如引理 \ref{lemma-compute}，其中有 $p + 1$ 个因子。设
$\chi_p : S_{p + 1} \to \{+1, -1\}$ 为符号特征标，并令
$U_p = W_p/S_{p + 1}$；$\underline{\mathbf{Z}}(\chi_p)$ 的含义同引理
\ref{lemma-twist}。则引理 \ref{lemma-alternating-cech-to-cohomology}
中的谱序列，其 $E_1$ 页为
$$
E_1^{p, q} =
H^q(U_p, \mathcal{F}|_{U_p} \otimes_\mathbf{Z} \underline{\mathbf{Z}}(\chi_p))
$$
并收敛到 $H^{p + q}(X, \mathcal{F})$。
\end{lemma}

\begin{proof}
注意，$S_{p + 1}$ 在 $W_p$ 上的作用是 $X$ 上的作用，所以确实得到态射
$U_p \to X$。由于 $W_p \to X$ 平展，且 $W_p \to U_p$ 是平展满态射，
$U_p \to X$ 也平展；见“代数空间的态射”引理
\ref{spaces-morphisms-lemma-etale-local}。因此
$\textit{Ab}(X_\etale)$ 的单射对象限制为
$\textit{Ab}(U_{p, \etale})$ 的单射对象；见“位点上的上同调”引理
\ref{sites-cohomology-lemma-cohomology-of-open}。此外，函子
$\mathcal{G} \mapsto
% P09-ERR-0006: frozen source has an extra closing parenthesis; see ERRATA_CANDIDATES.jsonl.
\mathcal{G} \otimes_\mathbf{Z} \underline{\mathbf{Z}}(\chi_p))$
% P09-ERR-0007: frozen source writes Ab(U_p), unlike the surrounding etale categories.
是 $\textit{Ab}(U_p)$ 的自等价，因而把单射对象变为单射对象，并且是正合的
（因为 $\underline{\mathbf{Z}}(\chi_p)$ 是可逆
$\underline{\mathbf{Z}}$-模）。因此，给定 $\textit{Ab}(X_\etale)$ 中的单射
消解 $\mathcal{F} \to \mathcal{I}^\bullet$，复形
$$
\Gamma(U_p,
\mathcal{I}^0|_{U_p} \otimes_\mathbf{Z} \underline{\mathbf{Z}}(\chi_p))
\to
\Gamma(U_p,
\mathcal{I}^1|_{U_p} \otimes_\mathbf{Z} \underline{\mathbf{Z}}(\chi_p))
\to
\Gamma(U_p,
\mathcal{I}^2|_{U_p} \otimes_\mathbf{Z} \underline{\mathbf{Z}}(\chi_p))
\to \ldots
$$
计算
$H^*(U_p,
\mathcal{F}|_{U_p} \otimes_\mathbf{Z} \underline{\mathbf{Z}}(\chi_p))$.
另一方面，由引理 \ref{lemma-twist}，它等于下列复形中
$S_{p + 1}$-反不变元所成的复形：
$$
\Gamma(W_p, \mathcal{I}^0) \to
\Gamma(W_p, \mathcal{I}^1) \to
\Gamma(W_p, \mathcal{I}^2) \to \ldots
$$
由引理 \ref{lemma-compute}，后者又等于复形
$$
\Hom(K^p, \mathcal{I}^0) \to
\Hom(K^p, \mathcal{I}^1) \to
\Hom(K^p, \mathcal{I}^2) \to \ldots
$$
它计算
$\Ext^*_{\textit{Ab}(X_\etale)}(K^p, \mathcal{F})$.
汇总以上各点即得结论。
\end{proof}





\section{拟凝聚层的高阶消失性}
\label{section-higher-vanishing}

\noindent
本节证明：给定拟紧拟分离代数空间 $X$，存在整数 $n = n(X)$，使得 $X$ 上
任意拟凝聚层的上同调在次数超过 $n$ 时消失。

\begin{lemma}
\label{lemma-quasi-coherent-twist}
沿用引理 \ref{lemma-twist} 中的 $S$、$W$、$G$、$U$、$\chi$。
若 $\mathcal{F}$ 是拟凝聚 $\mathcal{O}_U$-模，则
$\mathcal{F} \otimes_{\mathbf{Z}} \underline{\mathbf{Z}}(\chi)$ 也是。
\end{lemma}

\begin{proof}
$\mathcal{O}_U$-模结构是显然的。要检验
$\mathcal{F} \otimes_{\mathbf{Z}} \underline{\mathbf{Z}}(\chi)$ 拟凝聚，
只需平展局部地检验。由于 $\underline{\mathbf{Z}}(\chi)$ 作为
$\underline{\mathbf{Z}}$-模是有限局部自由的，引理随即成立。
\end{proof}

\noindent
即使 $X$ 是概形，下面的命题也很有意义。它是“概形的上同调”引理
\ref{coherent-lemma-vanishing-nr-affines} 的自然推广。在陈述之前先注意：
给定从仿射概形到拟分离代数空间 $X$ 的平展态射 $f : U \to X$，$f$ 的纤维
普遍有界。特别地，存在整数 $d$，使得 $|U| \to |X|$ 的每个纤维至多有
$d$ 个点；这就是“合宜代数空间”引理
\ref{decent-spaces-lemma-bounded-fibres} 中的蕴含
$(\eta) \Rightarrow (\delta)$。

\begin{proposition}
\label{proposition-vanishing}
设 $S$ 为概形，$X$ 为 $S$ 上拟紧且分离的代数空间。设 $U$ 为仿射概形，
$f : U \to X$ 为平展满态射，并设 $d$ 是 $|U| \to |X|$ 各纤维大小的上界。
则对任意拟凝聚 $\mathcal{O}_X$-模 $\mathcal{F}$，当 $q \geq d$ 时有
$H^q(X, \mathcal{F}) = 0$。
\end{proposition}

\begin{proof}
使用引理 \ref{lemma-alternating-spectral-sequence} 的谱序列。由于 $U$ 分离，
$f$ 也分离，所以该引理适用；见“代数空间的态射”引理
\ref{spaces-morphisms-lemma-compose-after-separated}。因为 $X$ 分离，概形
$U \times_X \ldots \times_X U$ 是
$U \times_{\Spec(\mathbf{Z})} \ldots \times_{\Spec(\mathbf{Z})} U$ 的闭子概形，
因而仿射。所以 $W_p$ 仿射。进而由“群胚”命题
\ref{groupoids-proposition-finite-flat-equivalence}，
$U_p = W_p/S_{p + 1}$ 是仿射概形。
% P09-ERR-0008: frozen source switches from U_p to W_p in the next three loci.
第 \ref{section-higher-direct-image} 节的讨论说明，$W_p$ 上拟凝聚层的上同调
（把 $W_p$ 看作代数空间）与底层仿射概形上相应拟凝聚层的上同调一致；
所以由“概形的上同调”引理
\ref{coherent-lemma-quasi-coherent-affine-cohomology-zero}，它在正次数上消失。
由引理 \ref{lemma-quasi-coherent-twist}，层
$\mathcal{F}|_{U_p} \otimes_\mathbf{Z} \underline{\mathbf{Z}}(\chi_p)$
拟凝聚。因此
$H^q(W_p,
\mathcal{F}|_{U_p} \otimes_\mathbf{Z} \underline{\mathbf{Z}}(\chi_p))$
在 $q > 0$ 时为零。由整数 $d$ 的定义，当 $p \geq d$ 时
$W_p = \emptyset$。所以
$H^0(W_p,
\mathcal{F}|_{U_p} \otimes_\mathbf{Z} \underline{\mathbf{Z}}(\chi_p))$
在 $p \geq d$ 时也为零。这就证明了命题。
\end{proof}

\noindent
下面的引理证明：拟紧拟分离代数空间对于拟凝聚模具有有限上同调维数。
我们明确写出这个界，只是因为稍后要用它证明高阶正像的类似结果。

\begin{lemma}
\label{lemma-vanishing-quasi-separated}
设 $S$ 为概形，$X$ 为 $S$ 上拟紧拟分离的代数空间。则可以选取：
\begin{enumerate}
\item 一个仿射概形 $U$；
\item 一个平展满态射 $f : U \to X$；
\item 一个整数 $d$，它界定 $U \to X$ 各纤维的次数；
\item 对每个 $p = 0, 1, \ldots, d$，从仿射概形 $V_p$ 出发的平展满态射
$V_p \to U_p$，其中 $U_p$ 的含义同引理
\ref{lemma-alternating-spectral-sequence}；以及
\item 一个整数 $d_p$，它界定 $V_p \to U_p$ 各纤维的次数。
\end{enumerate}
此外，只要有 (1)--(5)，则对任意拟凝聚 $\mathcal{O}_X$-模
$\mathcal{F}$，当 $q \geq \max(d_p + p)$ 时有
$H^q(X, \mathcal{F}) = 0$。
\end{lemma}

\begin{proof}
由于 $X$ 拟紧，可以找到平展满态射 $U \to X$，其中 $U$ 仿射；见
“代数空间的性质”引理
\ref{spaces-properties-lemma-quasi-compact-affine-cover}。由“合宜代数空间”引理
\ref{decent-spaces-lemma-bounded-fibres}，$f$ 的纤维普遍有界，故可找到 $d$。
有 $U_p = W_p/S_{p + 1}$，且
$W_p \subset U \times_X \ldots \times_X U$ 既开又闭。由于 $X$ 拟分离，
概形 $W_p$ 拟紧，因此 $U_p$ 拟紧。由于 $U$ 分离，概形 $W_p$ 分离，
所以由“代数空间”引理 \ref{spaces-lemma-quotient-finite-separated}
的绝对形式，$U_p$ 分离。由“代数空间的性质”引理
\ref{spaces-properties-lemma-quasi-compact-affine-cover}，
% P09-ERR-0009: frozen source says V_p -> W_p, contrary to the lemma statement.
可以找到态射 $V_p \to W_p$。再由“合宜代数空间”引理
\ref{decent-spaces-lemma-bounded-fibres}，可以找到各整数 $d_p$。

\medskip\noindent
此时，证明使用谱序列
$$
E_1^{p, q} =
H^q(U_p, \mathcal{F}|_{U_p} \otimes_\mathbf{Z} \underline{\mathbf{Z}}(\chi_p))
\Rightarrow
H^{p + q}(X, \mathcal{F})
$$
见引理 \ref{lemma-alternating-spectral-sequence}。由整数 $d$ 的定义，
% P09-ERR-0010: frozen source writes U_p = 0 for a space.
当 $p \geq d$ 时 $U_p = 0$。由命题 \ref{proposition-vanishing} 和引理
\ref{lemma-quasi-coherent-twist}，可知
$H^q(U_p,
\mathcal{F}|_{U_p} \otimes_\mathbf{Z} \underline{\mathbf{Z}}(\chi_p))$
当 $p = 0, \ldots, d$ 且 $q \geq d_p$ 时为零。引理随之成立。
\end{proof}









\section{高阶正像的消失性}
\label{section-vanishing-higher-direct-images}

\noindent
应用第 \ref{section-higher-vanishing} 节的结果，可得拟紧拟分离态射下
拟凝聚层高阶正像的消失性。这很有用，因为在某些情形下，它允许我们对上同调
次数作降归纳。

\begin{lemma}
\label{lemma-vanishing-higher-direct-images}
设 $S$ 为概形，$f : X \to Y$ 为 $S$ 上代数空间的态射。假设：
\begin{enumerate}
\item $f$ 拟紧且拟分离；
\item $Y$ 拟紧。
\end{enumerate}
则存在整数 $n(X \to Y)$，使得对任意代数空间 $Y'$、任意态射
$Y' \to Y$，以及 $X' = Y' \times_Y X$ 上任意拟凝聚层
$\mathcal{F}'$，当 $i \geq n(X \to Y)$ 时，高阶正像
$R^if'_*\mathcal{F}'$ 为零。
\end{lemma}

\begin{proof}
取平展满态射 $V \to Y$，其中 $V$ 是仿射概形；见“代数空间的性质”引理
\ref{spaces-properties-lemma-quasi-compact-affine-cover}。假设已经对基变换
$f_V : V \times_Y X \to V$ 证明了结论，那么结论对 $f$ 也成立，并可取
$n(X \to Y) = n(X_V \to V)$。事实上，若 $Y' \to Y$ 与
$\mathcal{F}'$ 如引理所述，则 $R^if'_*\mathcal{F}'|_{V \times_Y Y'}$
等于 $R^if'_{V, *}\mathcal{F}'|_{X'_V}$，其中
$f'_V : X'_V = V \times_Y Y' \times_Y X \to V \times_Y Y' = Y'_V$；
见“代数空间的性质”引理
\ref{spaces-properties-lemma-pushforward-etale-base-change-modules}。
因此可以假设 $Y$ 是仿射概形。

\medskip\noindent
此外，要对所有 $Y' \to Y$ 与 $\mathcal{F}'$ 证明消失性，只需处理 $Y'$
为仿射概形的情形。此时由引理 \ref{lemma-higher-direct-image}，
$R^if'_*\mathcal{F}'$ 拟凝聚。因此只需证明
$H^i(X', \mathcal{F}') = 0$；因为由“位点上的上同调”引理
\ref{sites-cohomology-lemma-apply-Leray} 以及仿射代数空间上拟凝聚层高阶
上同调的消失性（命题 \ref{proposition-vanishing}），有
$H^i(X', \mathcal{F}') = H^0(Y', R^if'_*\mathcal{F}')$。

\medskip\noindent
按引理 \ref{lemma-vanishing-quasi-separated} 选取
$U \to X$、$d$、$V_p \to U_p$ 与 $d_p$。对任意仿射概形 $Y'$ 及态射
$Y' \to Y$，记
$X' = Y' \times_Y X$、$U' = Y' \times_Y U$、
$V'_p = Y' \times_Y V_p$。则
% P09-ERR-0011: frozen source sets d'_p = d instead of d_p.
$U' \to X'$、$d' = d$、$V'_p \to U'_p$ 以及 $d'_p = d$，构成代数空间
$X'$ 的一组如引理 \ref{lemma-vanishing-quasi-separated} 所述的选择
（细节从略）。所以当 $i \geq \max(p + d_p)$ 时，
$H^i(X', \mathcal{F}') = 0$，结论成立。
\end{proof}

\begin{lemma}
\label{lemma-affine-vanishing-higher-direct-images}
设 $S$ 为概形，$f : X \to Y$ 为 $S$ 上代数空间的仿射态射。则对任意拟凝聚
$\mathcal{O}_X$-模 $\mathcal{F}$，当 $i > 0$ 时
$R^if_*\mathcal{F} = 0$。
\end{lemma}

\begin{proof}
回顾代数空间的仿射态射是可表的。因此结论由
(\ref{equation-representable-higher-direct-image}) 和“概形的上同调”引理
\ref{coherent-lemma-relative-affine-vanishing} 得到。
\end{proof}

\begin{lemma}
\label{lemma-relative-affine-cohomology}
设 $S$ 为概形，$f : X \to Y$ 为 $S$ 上代数空间的仿射态射，并设
$\mathcal{F}$ 为拟凝聚 $\mathcal{O}_X$-模。则对所有 $i \geq 0$，有
$H^i(X, \mathcal{F}) = H^i(Y, f_*\mathcal{F})$。
\end{lemma}

\begin{proof}
这由引理 \ref{lemma-affine-vanishing-higher-direct-images} 和 Leray 谱序列得到。
参见“位点上的上同调”引理 \ref{sites-cohomology-lemma-apply-Leray}。
\end{proof}






\section{支撑在闭子空间上的上同调}
\label{section-cohomology-support}

\noindent
本节是“上同调”第 \ref{cohomology-section-cohomology-support}、
\ref{cohomology-section-cohomology-support-bis} 节以及“平展上同调”第
\ref{etale-cohomology-section-cohomology-support} 节针对代数空间上阿贝尔层的
类似形式。

\medskip\noindent
设 $S$ 为概形，$X$ 为 $S$ 上的代数空间，$Z \subset X$ 为闭子空间，并设
$\mathcal{F}$ 为 $X_\etale$ 上的阿贝尔层。令
$$
\Gamma_Z(X, \mathcal{F}) =
\{s \in \mathcal{F}(X) \mid \text{Supp}(s) \subset Z\}
$$
为支撑在 $Z$ 中的截面（见“代数空间的性质”定义
\ref{spaces-properties-definition-support}）。这是一个左正合函子，但一般并不
正合。因此得到导出函子
$$
R\Gamma_Z(X, -) : D(X_\etale) \longrightarrow D(\textit{Ab})
$$
并把支撑在 $Z$ 中的上同调群定义为
$H^q_Z(X, \mathcal{F}) = R^q\Gamma_Z(X, \mathcal{F})$。

\medskip\noindent
设 $\mathcal{I}$ 为 $X_\etale$ 上的单射阿贝尔层，$U \subset X$ 为 $Z$ 的补
开子空间。则限制映射 $\mathcal{I}(X) \to \mathcal{I}(U)$ 为满射
（见“位点上的上同调”引理
\ref{sites-cohomology-lemma-restriction-along-monomorphism-surjective}），
其核为 $\Gamma_Z(X, \mathcal{I})$。立即可得：对 $K \in D(X_\etale)$，
在 $D(\textit{Ab})$ 中有区分三角
$$
R\Gamma_Z(X, K) \to R\Gamma(X, K) \to R\Gamma(U, K) \to R\Gamma_Z(X, K)[1]
$$
从而得到上同调长正合列
$$
\ldots \to H^i_Z(X, K) \to H^i(X, K) \to H^i(U, K) \to
H^{i + 1}_Z(X, K) \to \ldots
$$
其中 $K$ 是 $D(X_\etale)$ 的任意对象。

\medskip\noindent
对 $X_\etale$ 上的阿贝尔层 $\mathcal{F}$，可以考虑{\it 支撑在 $Z$ 中的截面
所成的子层}，记为 $\mathcal{H}_Z(\mathcal{F})$，其定义规则为
$$
\mathcal{H}_Z(\mathcal{F})(U) =
\{s \in \mathcal{F}(U) \mid \text{Supp}(s) \subset U \times_X Z\}
$$
这里使用“代数空间的性质”定义
\ref{spaces-properties-definition-support} 中的截面支撑。利用“代数空间的态射”
引理 \ref{spaces-morphisms-lemma-closed-immersion-push-pull} 的等价，
可把 $\mathcal{H}_Z(\mathcal{F})$ 看成 $Z_\etale$ 上的阿贝尔层。因此得到函子
$$
\textit{Ab}(X_\etale) \longrightarrow \textit{Ab}(Z_\etale),\quad
\mathcal{F} \longmapsto \mathcal{H}_Z(\mathcal{F})
$$
它左正合，但一般并不正合。

\begin{lemma}
\label{lemma-sections-with-support-acyclic}
设 $S$ 为概形，$i : Z \to X$ 为 $S$ 上代数空间的闭浸入，并设
$\mathcal{I}$ 为 $X_\etale$ 上的单射阿贝尔层。则
$\mathcal{H}_Z(\mathcal{I})$ 是 $Z_\etale$ 上的单射阿贝尔层。
\end{lemma}

\begin{proof}
注意，对 $Z_\etale$ 上任意阿贝尔层 $\mathcal{G}$，有
$$
\Hom_Z(\mathcal{G}, \mathcal{H}_Z(\mathcal{F})) =
\Hom_X(i_*\mathcal{G}, \mathcal{F})
$$
因为 $i_*\mathcal{G}$ 的任意截面终究都支撑在 $Z$ 中。由于 $i_*$ 正合
（引理 \ref{lemma-finite-higher-direct-image-zero}），且 $\mathcal{I}$ 在
$X_\etale$ 上单射，可知 $\mathcal{H}_Z(\mathcal{I})$ 在 $Z_\etale$ 上单射。
\end{proof}

\noindent
以
$$
R\mathcal{H}_Z : D(X_\etale) \longrightarrow D(Z_\etale)
$$
表示导出函子。令
$\mathcal{H}^q_Z(\mathcal{F}) = R^q\mathcal{H}_Z(\mathcal{F})$，从而
$\mathcal{H}^0_Z(\mathcal{F}) = \mathcal{H}_Z(\mathcal{F})$。
由上面的引理，有 Grothendieck 谱序列
$$
E_2^{p, q} = H^p(Z, \mathcal{H}^q_Z(\mathcal{F}))
\Rightarrow H^{p + q}_Z(X, \mathcal{F})
$$

\begin{lemma}
\label{lemma-cohomology-with-support-sheaf-on-support}
设 $S$ 为概形，$i : Z \to X$ 为 $S$ 上代数空间的闭浸入，并设
$\mathcal{G}$ 为 $Z_\etale$ 上的单射阿贝尔层。则当 $p > 0$ 时，
$\mathcal{H}^p_Z(i_*\mathcal{G}) = 0$。
\end{lemma}

\begin{proof}
这是因为函子 $i_*$ 正合（引理
\ref{lemma-finite-higher-direct-image-zero}），并把单射阿贝尔层变为单射
阿贝尔层（见“位点上的上同调”引理
\ref{sites-cohomology-lemma-pushforward-injective-flat}）。
\end{proof}

\begin{lemma}
\label{lemma-etale-localization-sheaf-with-support}
设 $S$ 为概形，$f : X \to Y$ 为 $S$ 上代数空间的平展态射。设
$Z \subset Y$ 为闭子空间，且 $f^{-1}(Z) \to Z$ 是代数空间的同构。
% P09-ERR-0012: frozen source places F on X although it is then pulled back from Y.
设 $\mathcal{F}$ 为 $X$ 上的阿贝尔层。则
$$
\mathcal{H}^q_Z(\mathcal{F}) = \mathcal{H}^q_{f^{-1}(Z)}(f^{-1}\mathcal{F})
$$
作为 $Z = f^{-1}(Z)$ 上的阿贝尔层有上述等式，并且
$H^q_Z(Y, \mathcal{F}) = H^q_{f^{-1}(Z)}(X, f^{-1}\mathcal{F})$。
\end{lemma}

\begin{proof}
因为 $f$ 平展，$\mathcal{F}$ 的单射消解拉回为 $f^{-1}\mathcal{F}$ 的单射
消解。因此只需对 $\mathcal{H}_Z(-)$ 检验等式，而这由定义直接得到。
支撑上同调的证明相同。若干细节从略。
\end{proof}

\noindent
设 $S$ 为概形，$X$ 为 $S$ 上的代数空间，$T \subset |X|$ 为闭子集。
以 $D_T(X_\etale)$ 表示 $D(X_\etale)$ 的严格满且饱和的三角子范畴，
其对象的上同调层都支撑在 $T$ 上。

\begin{lemma}
\label{lemma-complexes-with-support-on-closed}
设 $S$ 为概形，$i : Z \to X$ 为 $S$ 上代数空间的闭浸入。映射
$Ri_* = i_* : D(Z_\etale) \to D(X_\etale)$ 诱导等价
$D(Z_\etale) \to D_{|Z|}(X_\etale)$，其拟逆为
$$
% P09-ERR-0013: frozen source alternates D_Z and the defined D_|Z| notation.
i^{-1}|_{D_Z(X_\etale)} = R\mathcal{H}_Z|_{D_{|Z|}(X_\etale)}
$$
\end{lemma}

\begin{proof}
回顾 $i^{-1}$ 与 $i_*$ 是一对正合伴随函子，并且
% P09-ERR-0014: frozen source says identify functor; see ERRATA_CANDIDATES.jsonl.
$i^{-1}i_*$ 同构于阿贝尔层上的 identify 函子。见“代数空间的性质”引理
\ref{spaces-properties-lemma-stalk-pullback} 和“代数空间的态射”引理
\ref{spaces-morphisms-lemma-closed-immersion-push-pull}。因此
$i_* : D(Z_\etale) \to D_Z(X_\etale)$ 全忠实，而 $i^{-1}$ 给出左逆。
另一方面，设 $K$ 为 $D_Z(X_\etale)$ 的对象，并考虑伴随映射
$K \to i_*i^{-1}K$.
由 $i_*$ 与 $i^{-1}$ 的正合性，它在上同调层上诱导伴随映射
$H^n(K) \to i_*i^{-1}H^n(K)$。由于这些上同调层支撑在 $Z$ 上，
这些伴随映射是同构，从而 $D(Z_\etale) \to D_Z(X_\etale)$ 是等价。

\medskip\noindent
最后还要证明：若 $K$ 是 $D_Z(X_\etale)$ 的对象，则
$R\mathcal{H}_Z(K) = i^{-1}K$。可以使用刚刚证明的
$K = i_*i^{-1}K$。为 $i^{-1}K$ 选取一个 K-单射代表
$\mathcal{I}^\bullet$。由于 $i_*$ 是正合函子 $i^{-1}$ 的右伴随，复形
$i_*\mathcal{I}^\bullet$ 是 K-单射的（见“导出范畴”引理
\ref{derived-lemma-adjoint-preserve-K-injectives}）。于是
$R\mathcal{H}_Z(K)$ 由
$\mathcal{H}_Z(i_*\mathcal{I}^\bullet) = \mathcal{I}^\bullet$ 计算，
正合所需。
\end{proof}






\section{维数以上的消失性}
\label{section-vanishing-above-dimension}

\noindent
设 $S$ 为概形，$X$ 为 $S$ 上拟紧拟分离的代数空间。此时 $|X|$ 是谱空间；
见“代数空间的性质”引理
\ref{spaces-properties-lemma-quasi-compact-quasi-separated-spectral}。
此外，$X$ 的维数（定义见“代数空间的性质”定义
\ref{spaces-properties-definition-dimension}）等于 $|X|$ 的 Krull 维数；
见“合宜代数空间”引理
\ref{decent-spaces-lemma-dimension-decent-space}。我们将证明：$X$ 上拟凝聚层
的上同调在超过维数的次数上消失。即便只考虑域上有限型的拟分离代数空间，
这个结果也已经很有意义。

\begin{lemma}
\label{lemma-vanishing-above-dimension}
设 $S$ 为概形，$X$ 为 $S$ 上拟紧拟分离的代数空间，并假设对某个整数 $d$ 有
$\dim(X) \leq d$。
% P09-ERR-0015: frozen source repeats F in the noun phrase; see ERRATA_CANDIDATES.jsonl.
设 $\mathcal{F}$ 为 $X$ 上的拟凝聚层 $\mathcal{F}$。
\begin{enumerate}
\item $H^q(X, \mathcal{F}) = 0$ for $q > d$,
\item 对任意拟紧开子空间 $U \subset X$，
$H^d(X, \mathcal{F}) \to H^d(U, \mathcal{F})$ 为满射；
\item 对任意补集拟紧的闭子空间 $Z \subset X$，当 $q > d$ 时
$H^q_Z(X, \mathcal{F}) = 0$。
\end{enumerate}
\end{lemma}

\begin{proof}
由“代数空间的性质”引理
\ref{spaces-properties-lemma-dimension-decent-invariant-under-etale}，
每个在 $X$ 上平展的代数空间 $Y$ 都满足维数 $\leq d$。若 $Y$ 拟分离，
则由“合宜代数空间”引理
\ref{decent-spaces-lemma-dimension-decent-space}，$Y$ 的维数等于 $|Y|$ 的
Krull 维数。此外，若 $Y$ 是概形，则 $\mathcal{F}$ 在 $Y$ 上的平展上同调
（相应地，支撑在闭子概形上的平展上同调）与 $\mathcal{F}$ 的通常上同调
（相应地，支撑在该闭子概形上的通常上同调）一致。见“下降”命题
\ref{descent-proposition-same-cohomology-quasi-coherent} 和“平展上同调”引理
\ref{etale-cohomology-lemma-cohomology-with-support-quasi-coherent}。
下文将不加说明地使用这些事实。

\medskip\noindent
由“合宜代数空间”引理
\ref{decent-spaces-lemma-filter-quasi-compact-quasi-separated}，
存在整数 $n$ 以及开子空间
$$
\emptyset = U_{n + 1} \subset
U_n \subset U_{n - 1} \subset \ldots \subset U_1 = X
$$
它们具有如下性质：令 $T_p = U_p \setminus U_{p + 1}$，赋予其诱导的既约
子空间结构；则存在拟紧分离概形 $V_p$ 以及平展满态射
$f_p : V_p \to U_p$，使 $f_p^{-1}(T_p) \to T_p$ 为同构。

\medskip\noindent
由于 $U_n = V_n$ 是概形，由开头的说明和“上同调”命题
\ref{cohomology-proposition-cohomological-dimension-spectral}，
$\mathcal{F}$ 在 $U_n$ 上的上同调在次数 $> d$ 时消失。归纳地假设已经证明：
当 $q > d$ 时，$H^q(U_{p + 1}, \mathcal{F}|_{U_{p + 1}}) = 0$。
要推出 $\mathcal{F}$ 在 $U_p$ 上的上同调在次数 $> d$ 时消失，只需证明
当 $q > d$ 时 $H_{T_p}^q(U_p, \mathcal{F})$ 为零。而由引理
\ref{lemma-etale-localization-sheaf-with-support}，有
$$
H^q_{T_p}(U_p, \mathcal{F}) = H^q_{f_p^{-1}(T_p)}(V_p, \mathcal{F})
$$
又因为 $V_p$ 是概形，所需消失性由“上同调”命题
\ref{cohomology-proposition-cohomological-dimension-spectral} 得到。
由此证明了 (1)。

\medskip\noindent
为证明 (2)，设 $U \subset X$ 为拟紧开子空间。考虑开子空间
$U' = U \cup U_n$，并令 $Z = U' \setminus U$。则
$g : U_n \to U'$ 是平展态射，且 $g^{-1}(Z) \to Z$ 为同构。因此由引理
\ref{lemma-etale-localization-sheaf-with-support}，有
$H^q_Z(U', \mathcal{F}) = H^q_Z(U_n, \mathcal{F})$。由于 $U_n$ 是概形，
可应用“上同调”命题
\ref{cohomology-proposition-cohomological-dimension-spectral}，所以它在次数
$> d$ 时消失。于是 $H^d(U', \mathcal{F}) \to H^d(U, \mathcal{F})$
为满射。归纳地假设问题已经归结到 $U$ 包含 $U_{p + 1}$ 的情形。
令 $U' = U \cup U_p$、$Z = U' \setminus U$，并使用平展态射
$f_p : V_p \to U'$；它满足 $f_p^{-1}(Z) \to Z$ 为同构。换言之，再次有
$$
H^q_Z(U', \mathcal{F}) = H^q_{f_p^{-1}(Z)}(V_p, \mathcal{F})
$$
它在次数 $> d$ 时消失。因此
$H^d(U', \mathcal{F}) \to H^d(U, \mathcal{F})$ 为满射。最终到达
$U_1 = X \subset U$ 的阶段，证明完成。

\medskip\noindent
形式论证说明 (2) 蕴含 (3)。
\end{proof}








\section{上同调与基变换（一）}
\label{section-cohomology-and-base-change}

\noindent
设 $S$ 为一个概形。设 $f : X \to Y$ 为 $S$ 上代数空间之间的态射，
并设 $\mathcal{F}$ 为 $X$ 上的拟凝聚层。再设 $g : Y' \to Y$ 为
$S$ 上代数空间之间的态射。记 $X' = X_{Y'} = Y' \times_Y X$ 为
$X$ 的基变换，并记 $f' : X' \to Y'$ 为 $f$ 的基变换。
还把投影记为 $g' : X' \to X$，并令
$\mathcal{F}' = (g')^*\mathcal{F}$。下图表示这一情形：
\begin{equation}
\label{equation-base-change-diagram}
\vcenter{
\xymatrix{
\mathcal{F}' = (g')^*\mathcal{F} &
X' \ar[r]_{g'} \ar[d]_{f'} &
X \ar[d]^f &
\mathcal{F} \\
Rf'_*\mathcal{F}' &
Y' \ar[r]^g &
Y &
Rf_*\mathcal{F}
}
}
\end{equation}
下面是我们所考虑的基变换性质最简单的情形。

\begin{lemma}
\label{lemma-affine-base-change}
设 $S$ 为一个概形。设 $f : X \to Y$ 为 $S$ 上代数空间之间的仿射态射，
并设 $\mathcal{F}$ 为拟凝聚 $\mathcal{O}_X$-模。此时
$f_*\mathcal{F} \cong Rf_*\mathcal{F}$ 是拟凝聚层，而且对于每个图
(\ref{equation-base-change-diagram})，都有
$$
g^*f_*\mathcal{F} = f'_*(g')^*\mathcal{F}.
$$
\end{lemma}

\begin{proof}
由 (\ref{equation-representable-higher-direct-image}) 附近的讨论，
问题归结为概形之间仿射态射的情形；该情形已在《概形的上同调》引理
\ref{coherent-lemma-affine-base-change} 中处理。
\end{proof}

\begin{lemma}[平坦基变换]
\label{lemma-flat-base-change-cohomology}
设 $S$ 为一个概形。考虑 $S$ 上代数空间的笛卡尔图
$$
\xymatrix{
X' \ar[d]_{f'} \ar[r]_{g'} & X \ar[d]^f \\
Y' \ar[r]^g & Y
}
$$
设 $\mathcal{F}$ 为拟凝聚 $\mathcal{O}_X$-模，其拉回为
$\mathcal{F}' = (g')^*\mathcal{F}$。假设 $g$ 平坦，且 $f$ 拟紧、拟分离。
对于任意 $i \geq 0$，
\begin{enumerate}
\item 《站点上的上同调》引理
\ref{sites-cohomology-lemma-base-change-map-flat-case}
中的基变换映射
$$
g^*R^if_*\mathcal{F} \longrightarrow R^if'_*\mathcal{F}',
$$
是同构；
\item 如果 $Y = \Spec(A)$ 且 $Y' = \Spec(B)$，则
$H^i(X, \mathcal{F}) \otimes_A B = H^i(X', \mathcal{F}')$.
\end{enumerate}
\end{lemma}

\begin{proof}
由《空间的态射》引理 \ref{spaces-morphisms-lemma-base-change-flat}，
态射 $g'$ 平坦。注意，$g$ 和 $g'$ 的平坦性等价于相应小 \'etale 带环站点
之间态射的平坦性，见《空间的态射》引理
\ref{spaces-morphisms-lemma-flat-morphism-sites}。因此可以应用
《站点上的上同调》引理
\ref{sites-cohomology-lemma-base-change-map-flat-case}
得到基变换映射
$$
g^*R^pf_*\mathcal{F} \longrightarrow R^pf'_*\mathcal{F}'
$$
为了证明此映射是同构，可以在 $Y'$ 的 \'etale 拓扑下局部地验证。
因此可以假设 $Y$ 和 $Y'$ 都是仿射概形。不妨写成
$Y = \Spec(A)$ 和 $Y' = \Spec(B)$。此时真正要证明的是映射
$$
H^p(X, \mathcal{F}) \otimes_A B \longrightarrow H^p(X_B, \mathcal{F}_B)
$$
是同构，其中 $X_B = \Spec(B) \times_{\Spec(A)} X$，而
$\mathcal{F}_B$ 是 $\mathcal{F}$ 在 $X_B$ 上的拉回。换言之，只需证明 (2)。

\medskip\noindent
固定一个平坦环同态 $A \to B$，并设 $X$ 为 $A$ 上拟紧且拟分离的代数空间。
注意，$g' : X_B \to X$ 是 $\Spec(B) \to \Spec(A)$ 的基变换，因而是仿射的。
故由引理 \ref{lemma-affine-vanishing-higher-direct-images}，
高阶正像 $R^i(g')_*\mathcal{F}_B$ 均为零。因此
$H^p(X_B, \mathcal{F}_B) = H^p(X, g'_*\mathcal{F}_B)$，见
《站点上的上同调》引理 \ref{sites-cohomology-lemma-apply-Leray}。
此外，有
$$
g'_*\mathcal{F}_B = \mathcal{F} \otimes_{\underline{A}} \underline{B}
$$
其中 $\underline{A}$、$\underline{B}$ 表示取值分别为 $A$、$B$ 的常值环层。
事实上，显然存在一个从右端到左端的映射。对于任意 \'etale 于 $X$ 的
仿射概形 $U$，有
\begin{align*}
g'_*\mathcal{F}_B(U) & = \mathcal{F}_B(\Spec(B) \times_{\Spec(A)} U) \\
& =
\Gamma(\Spec(B) \times_{\Spec(A)} U,
(\Spec(B) \times_{\Spec(A)} U \to U)^*\mathcal{F}|_U) \\
& =
B \otimes_A \mathcal{F}(U)
\end{align*}
所以该映射是同构。利用 Lazard 定理，把 $B = \colim M_i$ 写成有限自由
$A$-模 $M_i$ 的滤过余极限，见《代数》定理
\ref{algebra-theorem-lazard}。由此得到
\begin{align*}
H^p(X, g'_*\mathcal{F}_B) &
= H^p(X, \mathcal{F} \otimes_{\underline{A}} \underline{B}) \\
& = H^p(X, \colim_i \mathcal{F} \otimes_{\underline{A}} \underline{M_i}) \\
& = \colim_i H^p(X, \mathcal{F} \otimes_{\underline{A}} \underline{M_i}) \\
& = \colim_i H^p(X, \mathcal{F}) \otimes_A M_i \\
& = H^p(X, \mathcal{F}) \otimes_A \colim_i M_i \\
& = H^p(X, \mathcal{F}) \otimes_A B
\end{align*}
第一个等式来自上面已经证明的
$g'_*\mathcal{F}_B = \mathcal{F} \otimes_{\underline{A}} \underline{B}$。
第二个等式是因为 $\otimes$ 与余极限交换。第三个等式是因为 $X$ 上的
上同调与余极限交换（见引理 \ref{lemma-colimits}）。第四个等式是因为
$M_i$ 有限自由（也就是说，上同调与有限直和交换）。第五个等式是因为
$\otimes$ 与余极限交换。第六个等式来自我们对该系统的选取。
\end{proof}

\begin{lemma}
\label{lemma-etale-pull-push-split}
设 $f : X \to Y$ 为代数空间之间拟紧、分离的 \'etale 态射。则对于任意
拟凝聚 $\mathcal{O}_X$-模 $\mathcal{F}$，映射
$f^*f_*\mathcal{F} \to \mathcal{F}$ 是分裂的。
\end{lemma}

\begin{proof}
考虑笛卡尔图
$$
\xymatrix{
X \times_Y X \ar[r]_-p \ar[d]_q & X \ar[d]^f \\
X \ar[r]^f & Y
}
$$
由引理 \ref{lemma-affine-base-change}，有
$f^*f_*\mathcal{F} = q_*p^*\mathcal{F}$。态射
$\Delta : X \to X \times_Y X$ 是开的（因为它是两个 \'etale 于 $Y$ 的
代数空间之间的态射），又因 $f$ 分离而是闭的。因此可见，
$p^*\mathcal{F}$ 是 $\Delta_*\mathcal{F}$ 与一个支撑在
$\Delta(X)$ 的（开闭）补集上的拟凝聚模的直和。追踪这些映射，读者可验证
$$
\mathcal{F} = q_*\Delta_*\mathcal{F} \to
q_*p^*\mathcal{F} = f^*f_*\mathcal{F} \to
\mathcal{F}
$$
是恒等映射。细节从略。
\end{proof}



\section{局部 Noether 代数空间上的凝聚模}
\label{section-coherent}

\noindent
本节对应于《概形的上同调》\ref{coherent-section-coherent-sheaves} 节。
在《站点上的模》定义 \ref{sites-modules-definition-site-local} 中，
我们已经定义了任意带环拓扑斯上的凝聚模。这里用这一概念定义局部
Noether 代数空间上的凝聚模。虽然也可以在更一般的情形下研究凝聚模，
但我们不打算这样做。

\begin{definition}
\label{definition-coherent}
设 $S$ 为一个概形，$X$ 为 $S$ 上的局部 Noether 代数空间。如果 $X$ 上的
拟凝聚模 $\mathcal{F}$ 按照《站点上的模》定义
\ref{sites-modules-definition-site-local} 的意义，是站点 $X_\etale$ 上的
凝聚 $\mathcal{O}_X$-模，则称 $\mathcal{F}$ 是{\it 凝聚的}。
\end{definition}

\noindent
这一定义与局部 Noether 概形上已有的凝聚模概念相容；见《空间的性质》
\ref{spaces-properties-section-properties-modules} 节中的断言 (5)（或更直接地，
见《下降》引理 \ref{descent-lemma-equivalence-quasi-coherent-properties}）。
因此，从现在起，如果 $X$ 是 $S$ 上的局部 Noether 概形，我们将不再区分
把 $X$ 看作概形时的凝聚模与把 $X$ 看作代数空间时的凝聚模；这与
《空间的性质》\ref{spaces-properties-section-quasi-coherent} 节所讨论的
拟凝聚模范畴的相应等同相容。

\medskip\noindent
在上述约定下，下面的引理给出了局部 Noether 代数空间上凝聚模的一种
直观刻画。

\begin{lemma}
\label{lemma-coherent-Noetherian}
设 $S$ 为一个概形，$X$ 为 $S$ 上的局部 Noether 代数空间，并设
$\mathcal{F}$ 为一个 $\mathcal{O}_X$-模。下列条件等价：
\begin{enumerate}
\item $\mathcal{F}$ 是凝聚的；
\item $\mathcal{F}$ 是有限型拟凝聚 $\mathcal{O}_X$-模；
\item $\mathcal{F}$ 是有限表示的 $\mathcal{O}_X$-模；
\item 对于任意 \'etale 态射 $\varphi : U \to X$，其中 $U$ 是概形，
拉回 $\varphi^*\mathcal{F}$ 是 $U$ 上的凝聚模；
\item 存在满 \'etale 态射 $\varphi : U \to X$，其中 $U$ 是概形，
使得拉回 $\varphi^*\mathcal{F}$ 是 $U$ 上的凝聚模。
\end{enumerate}
特别地，$\mathcal{O}_X$ 是凝聚的，任意可逆 $\mathcal{O}_X$-模是凝聚的；
更一般地，任意有限局部自由 $\mathcal{O}_X$-模都是凝聚的。
\end{lemma}

\begin{proof}
首先，如果 $X$ 是局部 Noether 代数空间，而 $U \to X$ 是 \'etale 态射，
则 $U$ 是局部 Noether 的，见《空间的性质》
\ref{spaces-properties-section-types-properties} 节。于是本引理由
《空间的性质》\ref{spaces-properties-section-properties-modules} 节中的
(1)--(5) 以及局部 Noether 概形上凝聚模的相应结论推出；后者见
《概形的上同调》引理 \ref{coherent-lemma-coherent-Noetherian}。
\end{proof}

\begin{lemma}
\label{lemma-coherent-abelian-Noetherian}
设 $S$ 为一个概形，$X$ 为 $S$ 上的局部 Noether 代数空间。
凝聚 $\mathcal{O}_X$-模所成的范畴是 Abel 范畴。更确切地说，凝聚
$\mathcal{O}_X$-模之间态射的核与余核都是凝聚的。凝聚层的任意扩张也是
凝聚的。
\end{lemma}

\begin{proof}
选取一个概形 $U$ 和满 \'etale 态射 $f : U \to X$。拉回 $f^*$ 等于限制函子，
因而是正合函子，见《空间的性质》公式
(\ref{spaces-properties-equation-restrict-modules})。由引理
\ref{lemma-coherent-Noetherian}，可以通过检验 $f^*\mathcal{F}$ 是否凝聚来
检验 $\mathcal{O}_X$-模 $\mathcal{F}$ 是否凝聚。因此本引理由概形的情形
推出；该情形即《概形的上同调》引理
\ref{coherent-lemma-coherent-abelian-Noetherian}。
\end{proof}

\noindent
凝聚模构成拟凝聚 $\mathcal{O}_X$-模范畴的一个 Serre 子范畴。对于一般
带环拓扑斯上的模，这一结论并不成立。

\begin{lemma}
\label{lemma-coherent-Noetherian-quasi-coherent-sub-quotient}
设 $S$ 为一个概形，$X$ 为 $S$ 上的局部 Noether 代数空间，并设
$\mathcal{F}$ 为凝聚 $\mathcal{O}_X$-模。$\mathcal{F}$ 的任意拟凝聚子模
都是凝聚的，$\mathcal{F}$ 的任意拟凝聚商模也都是凝聚的。
\end{lemma}

\begin{proof}
选取一个概形 $U$ 和满 \'etale 态射 $f : U \to X$。拉回 $f^*$ 等于限制函子，
因而是正合函子，见《空间的性质》公式
(\ref{spaces-properties-equation-restrict-modules})。由引理
\ref{lemma-coherent-Noetherian}，可以通过检验 $f^*\mathcal{H}$ 是否凝聚来
检验 $\mathcal{O}_X$-模 $\mathcal{G}$ 是否凝聚。% P09-ERR-0016
因此本引理由概形的情形推出；该情形即《概形的上同调》引理
\ref{coherent-lemma-coherent-Noetherian-quasi-coherent-sub-quotient}。
\end{proof}

\begin{lemma}
\label{lemma-tensor-hom-coherent}
设 $S$ 为一个概形。设 $X$ 为 $S$ 上的局部 Noether 代数空间，。% P09-ERR-0017
设 $\mathcal{F}$、$\mathcal{G}$ 为凝聚 $\mathcal{O}_X$-模。则
$\mathcal{O}_X$-模 $\mathcal{F} \otimes_{\mathcal{O}_X} \mathcal{G}$ 和
$\SheafHom_{\mathcal{O}_X}(\mathcal{F}, \mathcal{G})$ 都是凝聚的。
\end{lemma}

\begin{proof}
由引理 \ref{lemma-coherent-Noetherian}，这可由概形的相应结论推出，见
《概形的上同调》引理 \ref{coherent-lemma-tensor-hom-coherent}。
\end{proof}

\begin{lemma}
\label{lemma-local-isomorphism}
设 $S$ 为一个概形，$X$ 为 $S$ 上的局部 Noether 代数空间。设
$\mathcal{F}$、$\mathcal{G}$ 为凝聚 $\mathcal{O}_X$-模，且
$\varphi : \mathcal{G} \to \mathcal{F}$ 为 $\mathcal{O}_X$-模同态。
设 $\overline{x}$ 为 $X$ 上位于 $x \in |X|$ 之上的几何点。
\begin{enumerate}
\item 如果 $\mathcal{F}_{\overline{x}} = 0$，则存在 $x$ 的开邻域
$X' \subset X$，使得 $\mathcal{F}|_{X'} = 0$。
\item 如果 $\varphi_{\overline{x}} : \mathcal{G}_{\overline{x}} \to
\mathcal{F}_{\overline{x}}$ 是单射，则存在 $x$ 的开邻域 $X' \subset X$，
使得 $\varphi|_{X'}$ 是单射。
\item 如果 $\varphi_{\overline{x}} : \mathcal{G}_{\overline{x}} \to
\mathcal{F}_{\overline{x}}$ 是满射，则存在 $x$ 的开邻域 $X' \subset X$，
使得 $\varphi|_{X'}$ 是满射。
\item 如果 $\varphi_{\overline{x}} : \mathcal{G}_{\overline{x}} \to
\mathcal{F}_{\overline{x}}$ 是双射，则存在 $x$ 的开邻域 $X' \subset X$，
使得 $\varphi|_{X'}$ 是同构。
\end{enumerate}
\end{lemma}

\begin{proof}
设 $\varphi : U \to X$ 为 \'etale 态射，其中 $U$ 是概形，并设 $u \in U$
为映到 $x$ 的点。由《空间的性质》引理
\ref{spaces-properties-lemma-stalk-quasi-coherent} 及
\ref{spaces-properties-lemma-describe-etale-local-ring}
以及《代数进阶》引理
\ref{more-algebra-lemma-dumb-properties-henselization}，可见
$\varphi_{\overline{x}}$ 是单射、满射或双射，当且仅当
$\varphi_u : \varphi^*\mathcal{F}_u \to \varphi^*\mathcal{G}_u$
具有相应性质。% P09-ERR-0018
因此可以应用本引理的概形版本，得到（必要时缩小 $U$ 后）映射
$\varphi^*\mathcal{F} \to \varphi^*\mathcal{G}$ 是单射、满射或同构。
设 $X' \subset X$ 为与 $|\varphi|(|U|) \subset |X|$ 对应的开子空间，见
《空间的性质》引理 \ref{spaces-properties-lemma-open-subspaces}。由于
$\{U \to X'\}$ 是 \'etale 拓扑的覆盖，便得到 $\varphi|_{X'}$ 如所要求地
为单射、满射或同构。最后，考察映射 $\mathcal{F} \to 0$ 即可看出 (1)
由 (2) 推出。
\end{proof}

\begin{lemma}
\label{lemma-coherent-support-closed}
设 $S$ 为一个概形，$X$ 为 $S$ 上的局部 Noether 代数空间。设
$\mathcal{F}$ 为凝聚 $\mathcal{O}_X$-模，$i : Z \to X$ 为
$\mathcal{F}$ 的概形论支撑，并设 $\mathcal{G}$ 为满足
$i_*\mathcal{G} = \mathcal{F}$ 的拟凝聚 $\mathcal{O}_Z$-模，见
《空间的态射》定义
\ref{spaces-morphisms-definition-scheme-theoretic-support}。则
$\mathcal{G}$ 是凝聚 $\mathcal{O}_Z$-模。
\end{lemma}

\begin{proof}
本引理的陈述是有意义的，因为凝聚模尤其是有限型的。此外，$Z \to X$
是闭浸入，因而是局部有限型的，所以 $Z$ 局部 Noether；见《空间的态射》引理
\ref{spaces-morphisms-lemma-immersion-locally-finite-type} 及
\ref{spaces-morphisms-lemma-locally-finite-type-locally-noetherian}。
最后，由于 $\mathcal{G}$ 是有限型的，引理
\ref{lemma-coherent-Noetherian} 表明它是凝聚 $\mathcal{O}_Z$-模。
\end{proof}

\begin{lemma}
\label{lemma-i-star-equivalence}
设 $S$ 为一个概形，$i : Z \to X$ 为 $S$ 上局部 Noether 代数空间之间的
闭浸入。设 $\mathcal{I} \subset \mathcal{O}_X$ 为定义 $Z$ 的拟凝聚理想层。
函子 $i_*$ 在被 $\mathcal{I}$ 零化的凝聚 $\mathcal{O}_X$-模范畴与
凝聚 $\mathcal{O}_Z$-模范畴之间诱导一个等价。
\end{lemma}

\begin{proof}
由《空间的态射》引理 \ref{spaces-morphisms-lemma-i-star-equivalence}，
此函子全忠实。设 $\mathcal{F}$ 为被 $\mathcal{I}$ 零化的凝聚
$\mathcal{O}_X$-模。由《空间的态射》引理
\ref{spaces-morphisms-lemma-i-star-equivalence}，可将
$\mathcal{F}$ 写成 $i_*\mathcal{G}$，其中 $\mathcal{G}$ 是 $Z$ 上的某个
拟凝聚层。为了检验 $\mathcal{G}$ 凝聚，可以在 \'etale 拓扑下局部地验证
（引理 \ref{lemma-coherent-Noetherian}）。选取一个由概形给出的 \'etale
覆盖后，概形的情形（《概形的上同调》引理
\ref{coherent-lemma-i-star-equivalence}）表明 $\mathcal{G}$ 凝聚。
因此该函子全忠实，证明完成。% P09-ERR-0019
\end{proof}

\begin{lemma}
\label{lemma-finite-pushforward-coherent}
设 $S$ 为一个概形，$f : X \to Y$ 为 $S$ 上代数空间之间的有限态射，且
$Y$ 局部 Noether。设 $\mathcal{F}$ 为凝聚 $\mathcal{O}_X$-模。假设 $f$
有限且 $Y$ 局部 Noether。% P09-ERR-0020
则当 $p > 0$ 时 $R^pf_*\mathcal{F} = 0$，并且 $f_*\mathcal{F}$ 凝聚。
\end{lemma}

\begin{proof}
选取一个概形 $V$ 和满 \'etale 态射 $V \to Y$。则
$V \times_Y X \to V$ 是局部 Noether 概形之间的有限态射。由
(\ref{equation-representable-higher-direct-image})，问题归结为概形的情形；
该情形即《概形的上同调》引理
\ref{coherent-lemma-finite-pushforward-coherent}。
\end{proof}






\section{Noether 空间上的凝聚层}
\label{section-coherent-quasi-compact}

\noindent
本节介绍 Noether 代数空间上凝聚层的一些性质。

\begin{lemma}
\label{lemma-acc-coherent}
设 $S$ 为一个概形，$X$ 为 $S$ 上的 Noether 代数空间，并设
$\mathcal{F}$ 为凝聚 $\mathcal{O}_X$-模。$\mathcal{F}$ 的拟凝聚子模满足
升链条件。换言之，给定任意拟凝聚子模序列
$$
\mathcal{F}_1 \subset \mathcal{F}_2 \subset \ldots \subset \mathcal{F}
$$
则存在某个 $n \geq 0$，使得
$\mathcal{F}_n = \mathcal{F}_{n + 1} = \ldots $。
\end{lemma}

\begin{proof}
选取一个仿射概形 $U$ 和满 \'etale 态射 $U \to X$（见《空间的性质》引理
\ref{spaces-properties-lemma-quasi-compact-affine-cover}）。由
《空间的态射》引理
\ref{spaces-morphisms-lemma-locally-finite-type-locally-noetherian}，
$U$ 是 Noether 概形。如果
$\mathcal{F}_n|_U = \mathcal{F}_{n + 1}|_U = \ldots$，则
$\mathcal{F}_n = \mathcal{F}_{n + 1} = \ldots$。因此结论可由概形的情形
推出，见《概形的上同调》引理 \ref{coherent-lemma-acc-coherent}。
\end{proof}

\begin{lemma}
\label{lemma-power-ideal-kills-sheaf}
设 $S$ 为一个概形，$X$ 为 $S$ 上的 Noether 代数空间。设 $\mathcal{F}$ 为
$X$ 上的凝聚层，并设 $\mathcal{I} \subset \mathcal{O}_X$ 为与闭子空间
$Z \subset X$ 对应的拟凝聚理想层。存在某个 $n \geq 0$ 使
$\mathcal{I}^n\mathcal{F} = 0$，当且仅当（在集合论意义下）
$\text{Supp}(\mathcal{F}) \subset Z$。
\end{lemma}

\begin{proof}
选取一个仿射概形 $U$ 和满 \'etale 态射 $U \to X$（见《空间的性质》引理
\ref{spaces-properties-lemma-quasi-compact-affine-cover}）。由
《空间的态射》引理
\ref{spaces-morphisms-lemma-locally-finite-type-locally-noetherian}，
$U$ 是 Noether 概形。注意，$\mathcal{I}^n\mathcal{F}|_U = 0$ 当且仅当
$\mathcal{I}^n\mathcal{F} = 0$，支撑的条件也同样如此。因此结论可由概形
的情形推出，见《概形的上同调》引理
\ref{coherent-lemma-power-ideal-kills-sheaf}。
\end{proof}

\begin{lemma}[Artin-Rees]
\label{lemma-Artin-Rees}
设 $S$ 为一个概形，$X$ 为 $S$ 上的 Noether 代数空间。设 $\mathcal{F}$ 为
$X$ 上的凝聚层，$\mathcal{G} \subset \mathcal{F}$ 为拟凝聚子层，并设
$\mathcal{I} \subset \mathcal{O}_X$ 为拟凝聚理想层。则存在 $c \geq 0$，
使得对于所有 $n \geq c$，都有
$$
\mathcal{I}^{n - c}(\mathcal{I}^c\mathcal{F} \cap \mathcal{G})
=
\mathcal{I}^n\mathcal{F} \cap \mathcal{G}
$$
\end{lemma}

\begin{proof}
选取一个仿射概形 $U$ 和满 \'etale 态射 $U \to X$（见《空间的性质》引理
\ref{spaces-properties-lemma-quasi-compact-affine-cover}）。由
《空间的态射》引理
\ref{spaces-morphisms-lemma-locally-finite-type-locally-noetherian}，
$U$ 是 Noether 概形。本引理中的等式成立，当且仅当把它限制到 $U$ 后
成立。因此结论可由概形的情形推出，见《概形的上同调》引理
\ref{coherent-lemma-Artin-Rees}。
\end{proof}

\begin{lemma}
\label{lemma-homs-over-open}
设 $S$ 为一个概形，$X$ 为 $S$ 上的 Noether 代数空间。设 $\mathcal{F}$ 为
拟凝聚 $\mathcal{O}_X$-模，$\mathcal{G}$ 为凝聚 $\mathcal{O}_X$-模，并设
$\mathcal{I} \subset \mathcal{O}_X$ 为拟凝聚理想层。记 $Z \subset X$ 为
相应的闭子空间，并令 $U = X \setminus Z$。存在典范同构
$$
\colim_n \Hom_{\mathcal{O}_X}(\mathcal{I}^n\mathcal{G}, \mathcal{F})
\longrightarrow
\Hom_{\mathcal{O}_U}(\mathcal{G}|_U, \mathcal{F}|_U).
$$
特别地，有同构
$$
\colim_n \Hom_{\mathcal{O}_X}(\mathcal{I}^n, \mathcal{F})
\longrightarrow
\Gamma(U, \mathcal{F}).
$$
\end{lemma}

\begin{proof}
设 $W$ 为仿射概形，并设 $W \to X$ 为满 \'etale 态射（见《空间的性质》
引理 \ref{spaces-properties-lemma-quasi-compact-affine-cover}）。令
$R = W \times_X W$。则 $W$ 和 $R$ 都是 Noether 概形，见《空间的态射》
引理 \ref{spaces-morphisms-lemma-locally-finite-type-locally-noetherian}。
因此，由《概形的上同调》引理 \ref{coherent-lemma-homs-over-open}，
该结论对于 $\mathcal{F}$、$\mathcal{G}$、$\mathcal{I}$、$U$、$Z$ 在
$W$ 和 $R$ 上的限制成立。% P09-ERR-0021
由此形式地推出该结论在 $X$ 上成立。
\end{proof}






\section{凝聚层的逐层剥离}
\label{section-devissage}

\noindent
本节对应于《概形的上同调》\ref{coherent-section-devissage} 节。

\begin{lemma}
\label{lemma-prepare-filter-support}
设 $S$ 为一个概形，$X$ 为 $S$ 上的 Noether 代数空间，并设 $\mathcal{F}$
为 $X$ 上的凝聚层。假设
$\text{Supp}(\mathcal{F}) = Z \cup Z'$，其中 $Z$、$Z'$ 都是闭集。
则存在凝聚层的短正合列
$$
0 \to \mathcal{G}' \to \mathcal{F} \to \mathcal{G} \to 0
$$
满足 $\text{Supp}(\mathcal{G}') \subset Z'$ 且
$\text{Supp}(\mathcal{G}) \subset Z$。
\end{lemma}

\begin{proof}
设 $\mathcal{I} \subset \mathcal{O}_X$ 为定义 $Z$ 上约化诱导闭子空间结构的
理想层，见《空间的性质》引理
\ref{spaces-properties-lemma-reduced-closed-subspace}。考虑子层
$\mathcal{G}'_n = \mathcal{I}^n\mathcal{F}$ 以及商层
$\mathcal{G}_n = \mathcal{F}/\mathcal{I}^n\mathcal{F}$。对于每个 $n$，都有
短正合列
$$
0 \to \mathcal{G}'_n \to \mathcal{F} \to \mathcal{G}_n \to 0
$$
对于 $Z' \setminus Z$ 的每个几何点 $\overline{x}$，有
$\mathcal{I}_{\overline{x}} = \mathcal{O}_{X, \overline{x}}$，因而
$\mathcal{G}_{n, \overline{x}} = 0$。所以
$\text{Supp}(\mathcal{G}_n) \subset Z$。注意，$X \setminus Z'$ 是 Noether
代数空间。因此由引理 \ref{lemma-power-ideal-kills-sheaf}，存在 $n$ 使得
$\mathcal{G}'_n|_{X \setminus Z'} =
\mathcal{I}^n\mathcal{F}|_{X \setminus Z'} = 0$。对于这样的 $n$，有
$\text{Supp}(\mathcal{G}'_n) \subset Z'$。取
$\mathcal{G}' = \mathcal{G}'_n$ 和 $\mathcal{G} = \mathcal{G}_n$ 即可。
\end{proof}

\noindent
下文将自由使用《空间的态射》定义
\ref{spaces-morphisms-definition-scheme-theoretic-support} 中所定义的
有限型模的概形论支撑。

\begin{lemma}
\label{lemma-prepare-filter-irreducible}
设 $S$ 为一个概形，$X$ 为 $S$ 上的 Noether 代数空间，并设 $\mathcal{F}$
为 $X$ 上的凝聚层。假设 $\mathcal{F}$ 的概形论支撑是约化的
$Z \subset X$，且 $|Z|$ 不可约。则存在整数 $r > 0$、非零理想层
$\mathcal{I} \subset \mathcal{O}_Z$，以及凝聚层之间的单射
$$
i_*\left(\mathcal{I}^{\oplus r}\right) \to \mathcal{F}
$$
其余核支撑在 $Z$ 的一个真闭子空间上。% P09-ERR-0022
\end{lemma}

\begin{proof}
由假设，存在支撑为 $Z$ 的凝聚 $\mathcal{O}_Z$-模 $\mathcal{G}$，使得
$\mathcal{F} \cong i_*\mathcal{G}$，见引理
\ref{lemma-coherent-support-closed}。因此只需在 $Z = X$ 且
$i = \text{id}$ 的情形下证明本引理。

\medskip\noindent
由《空间的性质》命题
\ref{spaces-properties-proposition-locally-quasi-separated-open-dense-scheme}
存在一个稠密开子空间 $U \subset X$，且 $U$ 是概形。注意，$U$ 是 Noether
整概形。缩小 $U$ 后，可以假设
$\mathcal{F}|_U \cong \mathcal{O}_U^{\oplus r}$（例如由《概形的上同调》
引理 \ref{coherent-lemma-prepare-filter-irreducible}，或由直接的代数论证）。
设 $\mathcal{I} \subset \mathcal{O}_X$ 为拟凝聚理想层，其相应闭子空间是
$U$ 在 $X$ 中的补集（例如见《空间的性质》
\ref{spaces-properties-section-reduced} 节）。由引理
\ref{lemma-homs-over-open}，存在 $n \geq 0$ 以及态射
$\mathcal{I}^n(\mathcal{O}_X^{\oplus r}) \to \mathcal{F}$，它在 $U$ 上恢复
上述同构。由于
$\mathcal{I}^n(\mathcal{O}_X^{\oplus r}) = (\mathcal{I}^n)^{\oplus r}$，
便得到引理所述的映射。该映射是单射：事实上，若 $\sigma$ 是某个概形
$W$ 上 $\mathcal{I}^{\oplus r}$ 的非零截面，其中 $W$ \'etale 于 $X$，
则由于 $X$ 约化，从而 $W$ 约化，$\sigma$ 的支撑包含 $W$ 的一个非空开集。
然而，$(\mathcal{I}^n)^{\oplus r} \to \mathcal{F}$ 的核在一个稠密开集上
为零，所以 $\sigma$ 不可能是该核的截面。
\end{proof}

\begin{lemma}
\label{lemma-coherent-filter}
设 $S$ 为一个概形，$X$ 为 $S$ 上的 Noether 代数空间，并设 $\mathcal{F}$
为 $X$ 上的凝聚层。存在由凝聚子层组成的滤过
$$
0 = \mathcal{F}_0 \subset \mathcal{F}_1 \subset
\ldots \subset \mathcal{F}_m = \mathcal{F}
$$
使得对于每个 $j = 1, \ldots, m$，存在约化闭子空间 $Z_j \subset X$，
其中 $|Z_j|$ 不可约，并存在理想层
$\mathcal{I}_j \subset \mathcal{O}_{Z_j}$，使得
$$
\mathcal{F}_j/\mathcal{F}_{j - 1}
\cong (Z_j \to X)_* \mathcal{I}_j
$$
\end{lemma}

\begin{proof}
考虑集合
$$
\mathcal{T} =
\left\{
\begin{matrix}
T \subset |X|
\text{ 为闭集，且存在凝聚层 }
\mathcal{F} \\
\text{ 满足 }
\text{Supp}(\mathcal{F}) = T
\text{，而本引理对它不成立}
\end{matrix}
\right\}
$$
我们要证明 $\mathcal{T}$ 为空。否则，由于 $|X|$ 是 Noether 的（《空间的性质》
引理 \ref{spaces-properties-lemma-Noetherian-topology}），可以选取
$\mathcal{T}$ 的极小元 $T$。这意味着存在 $X$ 上支撑为 $T$ 的凝聚层
$\mathcal{F}$，而本引理对它不成立。显然 $T \not = \emptyset$，因为支撑
为空的层只能是零层，而本引理对零层成立（取 $m = 0$）。

\medskip\noindent
如果 $T$ 可约，则可写成 $T = Z_1 \cup Z_2$，其中 $Z_1$、$Z_2$ 都是
严格小于 $T$ 的闭集。于是可以应用引理
\ref{lemma-prepare-filter-support}，得到凝聚层的短正合列
$$
0 \to
\mathcal{G}_1 \to
\mathcal{F} \to
\mathcal{G}_2 \to 0
$$
满足 $\text{Supp}(\mathcal{G}_i) \subset Z_i$。由 $T$ 的极小性，每个
$\mathcal{G}_i$ 都有引理陈述中的滤过。考察 $\mathcal{F}$ 上诱导的滤过，
便得到矛盾。因此 $T$ 不可约。

\medskip\noindent
假设 $T$ 不可约。设 $\mathcal{J}$ 为定义 $T$ 上约化诱导闭子空间结构的
理想层，见《空间的性质》引理
\ref{spaces-properties-lemma-reduced-closed-subspace}。由引理
\ref{lemma-power-ideal-kills-sheaf}，存在 $n \geq 0$ 使得
$\mathcal{J}^n\mathcal{F} = 0$。因而得到滤过
$$
0 = \mathcal{I}^n\mathcal{F} \subset \mathcal{I}^{n - 1}\mathcal{F}
\subset \ldots \subset \mathcal{I}\mathcal{F} \subset \mathcal{F}
$$
其每个相继次商都被 $\mathcal{J}$ 零化。% P09-ERR-0023
因此，如果这些次商中的每一个都有引理陈述中的滤过，那么
$\mathcal{F}$ 也有。换言之，可以假设 $\mathcal{J}$ 零化 $\mathcal{F}$。

\medskip\noindent
假设 $T$ 不可约，且 $\mathcal{J}\mathcal{F} = 0$，其中 $\mathcal{J}$ 如上。
则 $\mathcal{F}$ 的概形论支撑为 $T$，见《空间的态射》引理
\ref{spaces-morphisms-lemma-i-star-equivalence}。因此可以应用引理
\ref{lemma-prepare-filter-irreducible}，得到短正合列
$$
0 \to
i_*(\mathcal{I}^{\oplus r}) \to
\mathcal{F} \to
\mathcal{Q} \to 0
$$
其中 $\mathcal{Q}$ 的支撑是 $T$ 的真闭子集。由 $T$ 的极小性，
$\mathcal{Q}$ 有所需类型的滤过。于是显然 $\mathcal{F}$ 也有这样的滤过，
这给出最终的矛盾。
\end{proof}

\begin{lemma}
\label{lemma-property-initial}
设 $S$ 为一个概形，$X$ 为 $S$ 上的 Noether 代数空间，并设
$\mathcal{P}$ 为 $X$ 上凝聚层的一条性质。假设：
\begin{enumerate}
\item 对于凝聚层的任意短正合列
$$
0 \to \mathcal{F}_1 \to \mathcal{F} \to \mathcal{F}_2 \to 0
$$
如果 $\mathcal{F}_i$（$i = 1, 2$）具有性质 $\mathcal{P}$，则
$\mathcal{F}$ 也具有性质 $\mathcal{P}$；
\item 对于每个满足 $|Z|$ 不可约的约化闭子空间 $Z \subset X$，以及每个
拟凝聚理想层 $\mathcal{I} \subset \mathcal{O}_Z$，$i_*\mathcal{I}$
都具有性质 $\mathcal{P}$。
\end{enumerate}
则 $X$ 上的每个凝聚层都具有性质 $\mathcal{P}$。
\end{lemma}

\begin{proof}
首先注意，如果凝聚层 $\mathcal{F}$ 有由凝聚子层组成的滤过
$$
0 = \mathcal{F}_0 \subset \mathcal{F}_1 \subset
\ldots \subset \mathcal{F}_m = \mathcal{F}
$$
并且每个 $\mathcal{F}_i/\mathcal{F}_{i - 1}$ 都具有性质 $\mathcal{P}$，
则 $\mathcal{F}$ 也具有性质 $\mathcal{P}$。这由 $\mathcal{P}$ 的性质 (1)
推出。另一方面，由引理 \ref{lemma-coherent-filter}，可以为任意
$\mathcal{F}$ 构造滤过，使其相继次商如 (2) 所述。因此本引理成立。
\end{proof}

\noindent
下面给出上述引理的一个更实用的变体。

\begin{lemma}
\label{lemma-property-higher-rank-cohomological}
设 $S$ 为一个概形，$X$ 为 $S$ 上的 Noether 代数空间，并设
$\mathcal{P}$ 为 $X$ 上凝聚层的一条性质。假设：
\begin{enumerate}
\item 对于凝聚层的任意短正合列
$$
0 \to \mathcal{F}_1 \to \mathcal{F} \to \mathcal{F}_2 \to 0
$$
如果 $\mathcal{F}_i$（$i = 1, 2$）具有性质 $\mathcal{P}$，则
$\mathcal{F}$ 也具有性质 $\mathcal{P}$；
\item 如果对于某个 $r \geq 1$，$\mathcal{F}^{\oplus r}$ 具有性质
$\mathcal{P}$，则 $\mathcal{F}$ 也具有性质 $\mathcal{P}$；
\item 对于每个满足 $|Z|$ 不可约的约化闭子空间 $i : Z \to X$，
存在 $Z$ 上的凝聚层 $\mathcal{G}$，使得
\begin{enumerate}
\item $\text{Supp}(\mathcal{G}) = Z$；
\item 对于每个非零拟凝聚理想层
$\mathcal{I} \subset \mathcal{O}_Z$，存在拟凝聚子层
$\mathcal{G}' \subset \mathcal{I}\mathcal{G}$，使得
$\text{Supp}(\mathcal{G}/\mathcal{G}')$ 是 $|Z|$ 中的真闭集，并且
$i_*\mathcal{G}'$ 具有性质 $\mathcal{P}$。
\end{enumerate}
\end{enumerate}
则 $X$ 上的每个凝聚层都具有性质 $\mathcal{P}$。
\end{lemma}

\begin{proof}
考虑集合
$$
\mathcal{T} =
\left\{
\begin{matrix}
T \subset |X|
\text{ 为非空闭集，且存在凝聚层 } \\
\mathcal{F}
\text{ 满足 }
\text{Supp}(\mathcal{F}) = T
\text{，而本引理对它不成立}
\end{matrix}
\right\}
$$
我们要证明 $\mathcal{T}$ 为空。否则，由于 $|X|$ 是 Noether 的（《空间的性质》
引理 \ref{spaces-properties-lemma-Noetherian-topology}），可以选取
$\mathcal{T}$ 的极小元 $T$。这意味着存在 $X$ 上支撑为 $T$ 的凝聚层
$\mathcal{F}$，而本引理对它不成立。

\medskip\noindent
如果 $T$ 可约，则可写成 $T = Z_1 \cup Z_2$，其中 $Z_1$、$Z_2$ 都是
严格小于 $T$ 的闭集。于是可以应用引理
\ref{lemma-prepare-filter-support}，得到凝聚层的短正合列
$$
0 \to
\mathcal{G}_1 \to
\mathcal{F} \to
\mathcal{G}_2 \to 0
$$
满足 $\text{Supp}(\mathcal{G}_i) \subset Z_i$。由 $T$ 的极小性，每个
$\mathcal{G}_i$ 都具有性质 $\mathcal{P}$。于是由 (1)，$\mathcal{F}$
具有性质 $\mathcal{P}$，这给出矛盾。

\medskip\noindent
假设 $T$ 不可约。设 $\mathcal{J}$ 为定义 $T$ 上约化诱导闭子空间结构的
理想层，见《空间的性质》引理
\ref{spaces-properties-lemma-reduced-closed-subspace}。由引理
\ref{lemma-power-ideal-kills-sheaf}，存在 $n \geq 0$ 使得
$\mathcal{J}^n\mathcal{F} = 0$。因而得到滤过
$$
0 = \mathcal{J}^n\mathcal{F} \subset \mathcal{J}^{n - 1}\mathcal{F}
\subset \ldots \subset \mathcal{J}\mathcal{F} \subset \mathcal{F}
$$
其每个相继次商都被 $\mathcal{J}$ 零化。因此，如果这些次商中的每一个
都有本引理陈述中的滤过，那么由 (1)，$\mathcal{F}$ 也有。% P09-ERR-0025
换言之，可以假设 $\mathcal{J}$ 零化 $\mathcal{F}$。

\medskip\noindent
假设 $T$ 不可约，且 $\mathcal{J}\mathcal{F} = 0$，其中 $\mathcal{J}$ 如上。
记 $i : Z \to X$ 为与 $\mathcal{J}$ 对应的闭子空间。则对某个凝聚
$\mathcal{O}_Z$-模 $\mathcal{H}$，有 $\mathcal{F} = i_*\mathcal{H}$；见
《空间的态射》引理 \ref{spaces-morphisms-lemma-i-star-equivalence} 以及
引理 \ref{lemma-coherent-support-closed}。设 $\mathcal{G}$ 为满足
(3)(a) 和 (3)(b) 的 $Z$ 上凝聚层。应用引理
\ref{lemma-prepare-filter-irreducible}，得到单射
$$
\mathcal{I}_1^{\oplus r_1} \to \mathcal{H}
\quad\text{以及}\quad
\mathcal{I}_2^{\oplus r_2} \to \mathcal{G}
$$
其余核的支撑是 $Z$ 中的真闭集。因此存在非空开集 $V \subset Z$，使得
$$
\mathcal{H}^{\oplus r_2}_V \cong \mathcal{G}^{\oplus r_1}_V
$$
设 $\mathcal{I} \subset \mathcal{O}_Z$ 为定义 $Z \setminus V$ 的拟凝聚
理想层，我们得到（引理 \ref{lemma-homs-over-open}）映射
$$
\mathcal{I}^n\mathcal{G}^{\oplus r_1} \longrightarrow \mathcal{H}^{\oplus r_2}
$$
它在 $V$ 上是同构。其核支撑在 $Z \setminus V$ 上，因而被
$\mathcal{I}$ 的某个幂零化，见引理
\ref{lemma-power-ideal-kills-sheaf}。因此增大 $n$ 后，可以假设所示映射
是单射，见引理 \ref{lemma-Artin-Rees}。应用 (3)(b)，可找到
$\mathcal{G}' \subset \mathcal{I}^n\mathcal{G}$，使得
$$
(i_*\mathcal{G}')^{\oplus r_1} \longrightarrow
i_*\mathcal{H}^{\oplus r_2} = \mathcal{F}^{\oplus r_2}
$$
是单射，其余核支撑在 $Z$ 的真闭子集上，并且 $i_*\mathcal{G}'$ 具有性质
$\mathcal{P}$。由 (1)，$(i_*\mathcal{G}')^{\oplus r_1}$ 具有性质
$\mathcal{P}$。再由 (1) 以及 $T = |Z|$ 的极小性，
$\mathcal{F}^{\oplus r_2}$ 具有性质 $\mathcal{P}$。最后由 (2)，
$\mathcal{F}$ 具有性质 $\mathcal{P}$，这正是所需的矛盾。
\end{proof}

\begin{lemma}
\label{lemma-property-higher-rank-cohomological-variant}
设 $S$ 为一个概形，$X$ 为 $S$ 上的 Noether 代数空间，并设
$\mathcal{P}$ 为 $X$ 上凝聚层的一条性质。假设：
\begin{enumerate}
\item 对于 $X$ 上凝聚层的任意短正合列，如果其中三个层中的两个具有
性质 $\mathcal{P}$，则第三个也具有性质 $\mathcal{P}$；
\item 如果对于某个 $r \geq 1$，$\mathcal{F}^{\oplus r}$ 具有性质
$\mathcal{P}$，则 $\mathcal{F}$ 也具有性质 $\mathcal{P}$；
\item 对于每个满足 $|Z|$ 不可约的约化闭子空间 $i : Z \to X$，
存在 $X$ 上概形论支撑为 $Z$ 的凝聚层 $\mathcal{G}$，并且
$\mathcal{G}$ 具有性质 $\mathcal{P}$。
\end{enumerate}
则 $X$ 上的每个凝聚层都具有性质 $\mathcal{P}$。
\end{lemma}

\begin{proof}
我们将证明引理 \ref{lemma-property-initial} 的条件 (1) 和 (2) 成立。
条件 (1) 显然成立。为证明 (2)，令
$$
\mathcal{T} =
\left\{
\begin{matrix}
i : Z \to X \text{ 为满足 }|Z|\text{ 不可约的约化闭子空间，并且}\\
\text{ 对某个拟凝聚 }\mathcal{I} \subset \mathcal{O}_Z\text{，}
i_*\mathcal{I}\text{ 不具有 }\mathcal{P}
\end{matrix}
\right\}
$$
如果 $\mathcal{T}$ 非空，则由于 $X$ 是 Noether 的，可以在
$\mathcal{T}$ 中找到极小的 $i : Z \to X$。下面证明这会导致矛盾。

\medskip\noindent
设 $\mathcal{G}$ 为假设 (3) 中假定存在的、概形论支撑为 $Z$ 的层。% P09-ERR-0026
设 $\varphi : i_*\mathcal{I}^{\oplus r} \to \mathcal{G}$ 如引理
\ref{lemma-prepare-filter-irreducible} 所述。设
$$
0 = \mathcal{F}_0 \subset \mathcal{F}_1 \subset
\ldots \subset \mathcal{F}_m = \Coker(\varphi)
$$
为引理 \ref{lemma-coherent-filter} 所述的滤过。由 $Z$ 的极小性和假设 (1)，
$\Coker(\varphi)$ 具有性质 $\mathcal{P}$。由于 $\varphi$ 是单射，再次使用
假设 (1) 可知 $i_*\mathcal{I}^{\oplus r}$ 具有性质 $\mathcal{P}$。
使用假设 (2)，便得到 $i_*\mathcal{I}$ 具有性质 $\mathcal{P}$。

\medskip\noindent
最后，若 $\mathcal{J} \subset \mathcal{O}_Z$ 是另一个拟凝聚理想层，令
$\mathcal{K} = \mathcal{I} \cap \mathcal{J}$，并考虑短正合列
$$
0 \to \mathcal{K} \to \mathcal{I} \to \mathcal{I}/\mathcal{K} \to 0
\quad
\text{以及}
\quad
0 \to \mathcal{K} \to \mathcal{J} \to \mathcal{J}/\mathcal{K} \to 0
$$
沿用上述论证并使用 $Z$ 的极小性，可见
$i_*\mathcal{I}/\mathcal{K}$ 和 $i_*\mathcal{J}/\mathcal{K}$
具有性质 $\mathcal{P}$。% P09-ERR-0027
因此由假设 (1)，先得到 $i_*\mathcal{K}$ 具有性质 $\mathcal{P}$，再得到
$i_*\mathcal{J}$ 具有性质 $\mathcal{P}$。换言之，$Z$ 不是
$\mathcal{T}$ 的元素，这正是所需的矛盾。
\end{proof}






\section{凝聚模的极限}
\label{section-limits}

\noindent
（局部 Noether 代数空间上的）凝聚模的余极限通常并不凝聚。但它是拟凝聚的，
因为代数空间上拟凝聚模的任意余极限都是拟凝聚的，见《空间的性质》引理
\ref{spaces-properties-lemma-properties-quasi-coherent}。反过来，如果该
代数空间是 Noether 的，那么每个拟凝聚模都是凝聚模的滤过余极限。

\begin{lemma}
\label{lemma-directed-colimit-coherent}
设 $S$ 为一个概形，$X$ 为 $S$ 上的 Noether 代数空间。每个拟凝聚
$\mathcal{O}_X$-模都是其凝聚子模的滤过余极限。
\end{lemma}

\begin{proof}
设 $\mathcal{F}$ 为拟凝聚 $\mathcal{O}_X$-模。如果
$\mathcal{G}, \mathcal{H} \subset \mathcal{F}$ 是凝聚
$\mathcal{O}_X$-子模，则 $\mathcal{G} \oplus \mathcal{H} \to \mathcal{F}$
的像是另一个同时包含二者的凝聚 $\mathcal{O}_X$-子模（见引理
\ref{lemma-coherent-abelian-Noetherian} 和
\ref{lemma-coherent-Noetherian-quasi-coherent-sub-quotient}）。由此可见，
该系统是有向的。因此现在只需证明 $\mathcal{F}$ 可以写成凝聚模的滤过
余极限；一旦如此，取这些模在 $\mathcal{F}$ 中的像，便可断定这样的
凝聚子模足够多。

\medskip\noindent
设 $U$ 为仿射概形，$U \to X$ 为满 \'etale 态射。令
$R = U \times_X U$，从而照例有 $X = U/R$。由《空间的性质》命题
\ref{spaces-properties-proposition-quasi-coherent}
可见 $\QCoh(\mathcal{O}_X) = \QCoh(U, R, s, t, c)$。因此问题归结为证明
$\QCoh(U, R, s, t, c)$ 的相应结论。于是本结论由更一般的《群胚》引理
\ref{groupoids-lemma-colimit-coherent} 推出。
\end{proof}

\begin{lemma}
\label{lemma-direct-colimit-finite-presentation}
设 $S$ 为一个概形，$f : X \to Y$ 为 $S$ 上代数空间之间的仿射态射，且
$Y$ 是 Noether 的。则每个拟凝聚 $\mathcal{O}_X$-模都是有限表示
$\mathcal{O}_X$-模的滤过余极限。
\end{lemma}

\begin{proof}
设 $\mathcal{F}$ 为拟凝聚 $\mathcal{O}_X$-模。写成
$f_*\mathcal{F} = \colim \mathcal{H}_i$，其中 $\mathcal{H}_i$ 是凝聚
$\mathcal{O}_Y$-模；见引理 \ref{lemma-directed-colimit-coherent}。由引理
\ref{lemma-coherent-Noetherian}，各 $\mathcal{H}_i$ 都是有限表示的
$\mathcal{O}_Y$-模。因此 $f^*\mathcal{H}_i$ 是有限表示的
$\mathcal{O}_X$-模，见《空间的性质》
\ref{spaces-properties-section-properties-modules} 节。我们断言映射
$$
\colim f^*\mathcal{H}_i = f^*f_*\mathcal{F} \to \mathcal{F}
$$
是满射，因为已假设 $f$ 仿射，即，选取一个概形 $V$ 和满 \'etale 态射
$V \to Y$。% P09-ERR-0030
令 $U = X \times_Y V$。则 $U$ 是概形，$f' : U \to V$ 仿射，且
$U \to X$ 是满 \'etale 态射。由《空间的性质》引理
\ref{spaces-properties-lemma-pushforward-etale-base-change-modules}
可见 $f'_*(\mathcal{F}|_U) = f_*\mathcal{F}|_V$，拉回也同样如此。因此
$f^*f_*\mathcal{F} \to \mathcal{F}$ 在 $U$ 上的限制是映射
$$
f^*f_*\mathcal{F}|_U = (f')^*(f_*\mathcal{F})|_V) =
(f')^*f'_*(\mathcal{F}|_U) \to \mathcal{F}|_U
$$
它是满射，因为 $f'$ 是概形之间的仿射态射。% P09-ERR-0031
所以该断言成立。

\medskip\noindent
由此得出，$X$ 上每个拟凝聚模都是有限表示模的某个滤过余极限的商。
特别地，$\mathcal{F}$ 是某个映射的余核
$$
\colim_{j \in J} \mathcal{G}_j \longrightarrow \colim_{i \in I} \mathcal{H}_i
$$
其中 $\mathcal{G}_j$ 和 $\mathcal{H}_i$ 都有限表示。注意，对于每个
$j \in I$，存在 $i \in I$ 以及态射
$\alpha : \mathcal{G}_j \to \mathcal{H}_i$，使得% P09-ERR-0032
$$
\xymatrix{
\mathcal{G}_j \ar[r]_\alpha \ar[d] & \mathcal{H}_i \ar[d] \\
\colim_{j \in J} \mathcal{G}_j \ar[r] &
\colim_{i \in I} \mathcal{H}_i
}
$$
交换，见引理 \ref{lemma-finite-presentation-quasi-compact-colimit}。
在这种情形下，$\Coker(\alpha)$ 是有限表示的 $\mathcal{O}_X$-模，并带有
一个映射 $\Coker(\alpha) \to \mathcal{F}$。考虑由上述三元组
$(i, j, \alpha)$ 组成的集合 $K$。规定
$(i, j, \alpha) \leq (i', j', \alpha')$，当且仅当
$i \leq i'$、$j \leq j'$，且下图
$$
\xymatrix{
\mathcal{G}_j \ar[r]_\alpha \ar[d] & \mathcal{H}_i \ar[d] \\
\mathcal{G}_{j'} \ar[r]^{\alpha'} &
\mathcal{H}_{i'}
}
$$
交换。由上可知，$K$ 是有向偏序集，
$$
\mathcal{F} = \colim_{(i, j, \alpha) \in K} \Coker(\alpha),
$$
于是结论得证。
\end{proof}






\section{上同调的消失性}
\label{section-vanishing}

\noindent
本节证明：如果一个拟紧且拟分离的代数空间上所有拟凝聚层的高阶上同调
都消失，则该代数空间仿射。我们将通过一系列引理来证明；一旦命题
\ref{proposition-vanishing-affine} 得证，这些引理便都会失去独立作用。

\begin{situation}
\label{situation-vanishing}
这里 $S$ 是一个概形，$X$ 是 $S$ 上拟紧且拟分离的代数空间，并满足下述
性质：对于每个拟凝聚 $\mathcal{O}_X$-模 $\mathcal{F}$，都有
$H^1(X, \mathcal{F}) = 0$。令 $A = \Gamma(X, \mathcal{O}_X)$。
\end{situation}

\noindent
我们希望证明典范态射
$$
p : X \longrightarrow \Spec(A)
$$
（见《空间的性质》引理
\ref{spaces-properties-lemma-morphism-to-affine-scheme}）是同构。
如果 $M$ 是 $A$-模，我们用 $M \otimes_A \mathcal{O}_X$ 表示拟凝聚模
$p^*\tilde M$。

\begin{lemma}
\label{lemma-vanishing-compute}
在情形 \ref{situation-vanishing} 中，对于 $A$-模 $M$，有
$p_*(M \otimes_A \mathcal{O}_X) = \widetilde{M}$ 且
$\Gamma(X, M \otimes_A \mathcal{O}_X) = M$。
\end{lemma}

\begin{proof}
由《空间的态射》引理 \ref{spaces-morphisms-lemma-pushforward}，
$p_*(M \otimes_A \mathcal{O}_X)$ 是 $\Spec(A)$ 上的拟凝聚模，故等式
$p_*(M \otimes_A \mathcal{O}_X) = \widetilde{M}$ 可由等式
$\Gamma(X, M \otimes_A \mathcal{O}_X) = M$ 推出。由引理
\ref{lemma-colimits}，注意到
$\Gamma(X, \bigoplus_{i \in I} \mathcal{O}_X) =
\bigoplus_{i \in I} A$。因此本引理对自由模成立。选取短正合列
$F_1 \to F_0 \to M$，其中 $F_0,F_1$ 是自由 $A$-模。由于
$H^1(X, -)$ 为零，全局截面函子右正合。此外，拉回 $p^*$ 也右正合。
所以
$$
\Gamma(X, F_1 \otimes_A \mathcal{O}_X) \to
\Gamma(X, F_0 \otimes_A \mathcal{O}_X) \to
\Gamma(X, M \otimes_A \mathcal{O}_X) \to 0
$$
正合。结论随即成立。
\end{proof}

\noindent
下面的引理表明，对 $X \to \Spec(A)$ 沿
$\Spec(A') \to \Spec(A)$ 作基变换后，情形 \ref{situation-vanishing}
仍然保持。

\begin{lemma}
\label{lemma-vanishing-base-change}
在情形 \ref{situation-vanishing} 中。
\begin{enumerate}
\item 给定代数空间之间的仿射态射 $X' \to X$，对于每个拟凝聚
$\mathcal{O}_{X'}$-模 $\mathcal{F}'$，都有
$H^1(X', \mathcal{F}') = 0$。
\item 给定 $A$-代数 $A'$，令 $X' = X \times_{\Spec(A)} \Spec(A')$，
则态射 $X' \to X$ 仿射，且 $\Gamma(X', \mathcal{O}_{X'}) = A'$。
\end{enumerate}
\end{lemma}

\begin{proof}
(1) 由引理 \ref{lemma-affine-vanishing-higher-direct-images} 以及 Leray
谱序列（《站点上的上同调》引理 \ref{sites-cohomology-lemma-Leray}）推出。
设 $A \to A'$ 如 (2) 所述。由于仿射态射在基变换下保持（《空间的态射》
引理 \ref{spaces-morphisms-lemma-base-change-affine}），且仿射概形之间的
态射是仿射的，所以 $X' \to X$ 仿射。由引理
\ref{lemma-affine-base-change}，有
$(X' \to X)_*\mathcal{O}_{X'} = A' \otimes_A \mathcal{O}_X$，因此
$$
\Gamma(X', \mathcal{O}_{X'}) =
\Gamma(X, (X' \to X)_*\mathcal{O}_{X'}) =
\Gamma(X, A' \otimes_A \mathcal{O}_X) = A'
$$
由引理 \ref{lemma-vanishing-compute} 可得所需等式。
\end{proof}

\begin{lemma}
\label{lemma-vanishing-separate-closed}
在情形 \ref{situation-vanishing} 中，设 $Z_0, Z_1 \subset |X|$ 为不交的
闭子集。则存在 $a \in A$，使得 $Z_0 \subset V(a)$ 且
$Z_1 \subset V(a - 1)$。
\end{lemma}

\begin{proof}
赋予 $Z_0$、$Z_1$ 约化诱导子空间结构（《空间的性质》定义
\ref{spaces-properties-definition-reduced-induced-space}），并将相应闭浸入
记为 $i_0 : Z_0 \to X$ 和 $i_1 : Z_1 \to X$。由于
$Z_0 \cap Z_1 = \emptyset$，可见拟凝聚 $\mathcal{O}_X$-模的典范映射
$$
\mathcal{O}_X
\longrightarrow
i_{0, *}\mathcal{O}_{Z_0} \oplus i_{1, *}\mathcal{O}_{Z_1}
$$
是满射（考察几何点处的茎即可）。由于 $H^1(X, -)$ 在此映射的核上为零，
所诱导的全局截面映射是满射。因此可以找到 $a \in A$，它映到右端的
全局截面 $(0, 1)$。
\end{proof}

\begin{lemma}
\label{lemma-vanishing-injective}
在情形 \ref{situation-vanishing} 中，态射 $p : X \to \Spec(A)$ 是泛单的。
\end{lemma}

\begin{proof}
设 $A \to k$ 为环同态，其中 $k$ 是域。只需证明
$\Spec(k) \times_{\Spec(A)} X$ 至多有一个点（见《空间的态射》引理
\ref{spaces-morphisms-lemma-universally-injective-local}）。利用引理
\ref{lemma-vanishing-base-change}，可以假设 $A$ 是域，而我们必须证明
$|X|$ 至多有一个点。

\medskip\noindent
把 $X$ 看作 $\Spec(k)$ 上的代数空间，并对任意扩张 $K/k$，用 $X(K)$
表示 $X$ 的 $K$-值点，见《空间的态射》
\ref{spaces-morphisms-section-points-fields} 节。如果 $K/k$ 是超越次数
足够大的代数闭域扩张，则 $X(K) \to |X|$ 是满射，见《空间的态射》引理
\ref{spaces-morphisms-lemma-large-enough}。因此用 $K$ 替换 $k$ 后，只需证明
（在 $A = k$ 的情形下）$X(k)$ 是单点集。

\medskip\noindent
设 $x, x' \in X(k)$。由《良态空间》引理
\ref{decent-spaces-lemma-algebraic-residue-field-extension-closed-point}，
$x$ 和 $x'$ 是 $|X|$ 的闭点。因此，如果 $x \not = x'$，由引理
\ref{lemma-vanishing-separate-closed}，$x$ 和 $x'$ 映到 $\Spec(k)$ 的不同点。
% P09-ERR-0034
由此得到 $x = x'$，如所要求。
\end{proof}

\begin{lemma}
\label{lemma-vanishing-separated}
在情形 \ref{situation-vanishing} 中，态射 $p : X \to \Spec(A)$ 是分离的。
\end{lemma}

\begin{proof}
由《良态空间》引理
\ref{decent-spaces-lemma-there-is-a-scheme-integral-over}
可以找到一个概形 $Y$ 和满整态射 $Y \to X$。整态射是仿射的，所以可以
应用引理 \ref{lemma-vanishing-base-change}，得到对于每个拟凝聚
$\mathcal{O}_Y$-模 $\mathcal{G}$，都有 $H^1(Y, \mathcal{G}) = 0$。
由于 $Y \to X$ 拟紧且 $X$ 拟紧，$Y$ 也拟紧。由于 $Y$ 是概形，可以应用
《概形的上同调》引理
\ref{coherent-lemma-quasi-compact-h1-zero-covering}，得到 $Y$ 仿射，因而
$Y$ 分离。注意，整态射是仿射且泛闭的，见《空间的态射》引理
\ref{spaces-morphisms-lemma-integral-universally-closed}。由《空间的态射》
引理 \ref{spaces-morphisms-lemma-image-universally-closed-separated}，
可见 $X$ 是分离代数空间。
\end{proof}

\begin{proposition}
\label{proposition-vanishing-affine}
\begin{slogan}
代数空间情形下的 Serre 仿射性判别法。
\end{slogan}
一个拟紧且拟分离的代数空间仿射，当且仅当拟凝聚层的所有高阶上同调群
都消失。更确切地说，情形 \ref{situation-vanishing} 中的任意代数空间
都是仿射概形。
\end{proposition}

\begin{proof}
选取一个仿射概形 $U = \Spec(B)$ 和满 \'etale 态射
$\varphi : U \to X$。令 $R = U \times_X U$。由于 $p$ 分离（引理
\ref{lemma-vanishing-separated}），$R$ 是
$U \times_{\Spec(A)} U = \Spec(B \otimes_A B)$ 的闭子概形。因此
$R = \Spec(C)$ 也是仿射的，而且环同态
$$
B \otimes_A B \longrightarrow C
$$
是满射。照例把两个映射记为 $s, t : B \to C$。选取
$g_1, \ldots, g_m \in B$，使得 $s(g_1), \ldots, s(g_m)$ 在
$t : B \to C$ 上生成 $C$（这是可能的，因为 $t : B \to C$ 有限表示，
且所示映射是满射）。于是 $g_1, \ldots, g_m$ 给出
$\varphi_*\mathcal{O}_U$ 的全局截面，并且映射
$$
\mathcal{O}_X[z_1, \ldots, z_n] \longrightarrow \varphi_*\mathcal{O}_U,
\quad
z_j \longmapsto g_j
$$
是满射：限制到 $U$ 上即可检验。% P09-ERR-0033
事实上，由引理 \ref{lemma-flat-base-change-cohomology}，
$\varphi^*\varphi_*\mathcal{O}_U = t_*\mathcal{O}_R$，因此恰好得到
$s(g_i)$ 在 $t : B \to C$ 上生成 $C$ 的条件。由于该映射的核的
$H^1$ 消失，可见
$$
\Gamma(X, \mathcal{O}_X[x_1, \ldots, x_n]) =
A[x_1, \ldots, x_n] \longrightarrow
\Gamma(X, \varphi_*\mathcal{O}_U) = \Gamma(U, \mathcal{O}_U) = B
$$
是满射。因此 $B$ 是有限型 $A$-代数，从而 $X \to \Spec(A)$ 有限型且分离。
由引理 \ref{lemma-vanishing-injective} 以及《空间的态射》引理
\ref{spaces-morphisms-lemma-locally-quasi-finite}，它还是局部拟有限的。
于是由《空间的态射》引理
\ref{spaces-morphisms-lemma-locally-quasi-finite-separated-representable}，
$X \to \Spec(A)$ 可表，从而 $X$ 是概形。最后，应用《概形的上同调》引理
\ref{coherent-lemma-quasi-compact-h1-zero-covering}，可知 $X$ 仿射，
因而等于 $\Spec(A)$。
\end{proof}

\begin{lemma}
\label{lemma-Noetherian-h1-zero}
设 $S$ 为一个概形，$X$ 为 $S$ 上的 Noether 代数空间。假设对于每个凝聚
$\mathcal{O}_X$-模 $\mathcal{F}$，都有 $H^1(X, \mathcal{F}) = 0$。
则 $X$ 是仿射概形。
\end{lemma}

\begin{proof}
由引理 \ref{lemma-directed-colimit-coherent} 和 \ref{lemma-colimits}，
该假设推出对于每个拟凝聚 $\mathcal{O}_X$-模 $\mathcal{F}$，都有
$H^1(X, \mathcal{F}) = 0$。于是由命题
\ref{proposition-vanishing-affine}，$X$ 仿射。
\end{proof}

\begin{lemma}
\label{lemma-Noetherian-h1-zero-invertible}
设 $S$ 为一个概形，$X$ 为 $S$ 上的 Noether 代数空间，并设
$\mathcal{L}$ 为可逆 $\mathcal{O}_X$-模。假设对于每个凝聚
$\mathcal{O}_X$-模 $\mathcal{F}$，存在 $n \geq 1$ 使得
$H^1(X, \mathcal{F} \otimes_{\mathcal{O}_X} \mathcal{L}^{\otimes n}) = 0$。
则 $X$ 是概形，而且 $\mathcal{L}$ 在 $X$ 上是丰沛的。
\end{lemma}

\begin{proof}
设 $s \in H^0(X, \mathcal{L}^{\otimes d})$ 为全局截面。设 $U \subset X$
为 $s$ 生成 $\mathcal{L}^{\otimes d}$ 的开子空间。特别地，有
$\mathcal{L}^{\otimes d}|_U \cong \mathcal{O}_U$。我们断言 $U$ 仿射。

\medskip\noindent
断言的证明。我们将证明，对于每个拟凝聚 $\mathcal{O}_U$-模
$\mathcal{F}$，都有 $H^1(U, \mathcal{F}) = 0$。由命题
\ref{proposition-vanishing-affine}，这将证明该断言。记
$j : U \to X$ 为包含态射。由于 $j$ 在 \'etale 拓扑下局部仿射（由
《态射》引理 \ref{morphisms-lemma-affine-s-open}），《空间的态射》引理
\ref{spaces-morphisms-lemma-affine-local} 表明 $j$ 仿射。因此
$$
H^1(U, \mathcal{F}) = H^1(X, j_*\mathcal{F})
$$
由引理 \ref{lemma-affine-vanishing-higher-direct-images}（以及《站点上的
上同调》引理 \ref{sites-cohomology-lemma-apply-Leray}）成立。把
$j_*\mathcal{F} = \colim \mathcal{F}_i$ 写成凝聚 $\mathcal{O}_X$-模的
滤过余极限，见引理 \ref{lemma-directed-colimit-coherent}。则
$$
H^1(X, j_*\mathcal{F}) = \colim H^1(X, \mathcal{F}_i)
$$
由引理 \ref{lemma-colimits} 成立。因此只需证明 $H^1(X, \mathcal{F}_i)$
在 $H^1(U, j^*\mathcal{F}_i)$ 中的像为零。由假设，存在 $n \geq 1$，使得
$$
H^1(X,
\mathcal{F}_i \otimes_{\mathcal{O}_X}
(\mathcal{O}_X \oplus \mathcal{L} \oplus \ldots
\oplus \mathcal{L}^{\otimes d - 1})
\otimes_{\mathcal{O}_X} \mathcal{L}^{\otimes n}) = 0
$$
因此存在 $a \geq 0$，使得
$H^1(X, \mathcal{F}_i \otimes_{\mathcal{O}_X} \mathcal{L}^{\otimes ad}) = 0$。
另一方面，映射
$$
s^a : \mathcal{F}_i \longrightarrow
\mathcal{F}_i \otimes_{\mathcal{O}_X} \mathcal{L}^{\otimes ad}
$$
限制到 $U$ 后是同构。考察交换图
$$
\xymatrix{
H^1(X, \mathcal{F}_i) \ar[r] \ar[d]_{s^a} & H^1(U, j^*\mathcal{F}_i)
\ar[d]^{\cong} \\
H^1(X, \mathcal{F}_i \otimes_{\mathcal{O}_X} \mathcal{L}^{\otimes ad}) \ar[r] &
H^1(U,
j^*(\mathcal{F}_i \otimes_{\mathcal{O}_X} \mathcal{L}^{\otimes ad}))
}
$$
可知映射 $H^1(X, \mathcal{F}_i) \to H^1(U, j^*\mathcal{F}_i)$ 为零，
断言得证。

\medskip\noindent
设 $x \in |X|$ 为闭点。由《良态空间》引理
\ref{decent-spaces-lemma-decent-space-closed-point}，可以用闭浸入
$i : \Spec(k) \to X$ 表示 $x$（这里还用到了拟分离代数空间是良态的，
见《良态空间》\ref{decent-spaces-section-reasonable-decent} 节）。因此
$\mathcal{O}_X \to i_*\mathcal{O}_{\Spec(k)}$ 是满射。设
$\mathcal{I} \subset \mathcal{O}_X$ 为其核，并选取 $d \geq 1$，使得
$H^1(X, \mathcal{I} \otimes_{\mathcal{O}_X} \mathcal{L}^{\otimes d}) = 0$。
则
$$
H^0(X, \mathcal{L}^{\otimes d}) \to
H^0(X,
i_*\mathcal{O}_{\Spec(k)} \otimes_{\mathcal{O}_X} \mathcal{L}^{\otimes d}) =
H^0(\Spec(k), i^*\mathcal{L}^{\otimes d}) \cong k
$$
由上同调长正合列，该映射是满射。因此存在
$s \in H^0(X, \mathcal{L}^{\otimes d})$，使得 $x \in U$，其中 $U$ 是上面
与 $s$ 对应的开子空间。由上述断言，$x$ 位于 $X$ 的概形轨迹中
（见《空间的性质》引理 \ref{spaces-properties-lemma-subscheme}）。

\medskip\noindent
为了得出 $X$ 是概形，只需证明 $|X|$ 中任何包含所有闭点的开子集都等于
$|X|$。这由 $|X|$ 是 Noether 拓扑空间这一事实推出，见《空间的性质》
引理 \ref{spaces-properties-lemma-Noetherian-sober}。最后，当 $X$ 是概形时，
可以应用《概形的上同调》引理
\ref{coherent-lemma-quasi-compact-h1-zero-invertible}，得到
$\mathcal{L}$ 是丰沛的。
\end{proof}









\section{有限态射与仿射性}
\label{section-finite-affine}

\noindent
本节对应于《概形的上同调》\ref{coherent-section-finite-affine} 节。

\begin{lemma}
\label{lemma-finite-morphism-Noetherian}
设 $S$ 为一个概形，$f : Y \to X$ 为 $S$ 上代数空间之间的态射。假设
$f$ 有限且为满射，并且 $X$ 局部 Noether。设 $i : Z \to X$ 为闭浸入。
记 $i' : Z' \to Y$ 为 $Z$ 的逆像（《空间的态射》
\ref{spaces-morphisms-section-closed-immersions} 节），并记
$f' : Z' \to Z$ 为诱导态射。则
$\mathcal{G} = f'_*\mathcal{O}_{Z'}$ 是支撑为 $Z$ 的凝聚
$\mathcal{O}_Z$-模。
\end{lemma}

\begin{proof}
注意，$f'$ 是 $f$ 的基变换，因而由《空间的态射》引理
\ref{spaces-morphisms-lemma-base-change-surjective} 和
\ref{spaces-morphisms-lemma-base-change-integral}，它有限且为满射。
由《空间的态射》引理
\ref{spaces-morphisms-lemma-locally-finite-type-locally-noetherian}（以及闭浸入
和有限态射都是有限型的这一事实），$Y$、$Z$、$Z'$ 都局部 Noether。
由引理 \ref{lemma-finite-pushforward-coherent}，$\mathcal{G}$ 是凝聚
$\mathcal{O}_Z$-模。$\mathcal{G}$ 的支撑在 $|Z|$ 中闭，见《空间的态射》
引理 \ref{spaces-morphisms-lemma-support-finite-type}。因此，如果
$\mathcal{G}$ 的支撑不等于 $|Z|$，则用一个开子空间替换 $X$ 后，可以
假设 $\mathcal{G} = 0$ 而 $Z \not = \emptyset$。这将意味着
$f'_*\mathcal{O}_{Z'} = 0$。特别地，截面
$1 \in \Gamma(Z', \mathcal{O}_{Z'}) = \Gamma(Z, f'_*\mathcal{O}_{Z'})$
将为零，从而 $Z' = \emptyset$ 是空代数空间。这不可能，因为
$Z' \to Z$ 是满射。
\end{proof}

\begin{lemma}
\label{lemma-affine-morphism-projection-ideal}
设 $S$ 为一个概形，$f : Y \to X$ 为 $S$ 上代数空间之间的态射。
设 $\mathcal{F}$ 为 $Y$ 上的拟凝聚层，$\mathcal{I}$ 为 $X$ 上的拟凝聚
理想层。如果 $f$ 仿射，则
$\mathcal{I}f_*\mathcal{F} = f_*(f^{-1}\mathcal{I}\mathcal{F})$
（记号的含义见证明）。
\end{lemma}

\begin{proof}
上述记号的含义如下。由于 $f^{-1}$ 是正合函子，
$f^{-1}\mathcal{I}$ 是 $f^{-1}\mathcal{O}_X$ 的理想层。通过
$Y_\etale$ 上的映射
$f^\sharp : f^{-1}\mathcal{O}_X \to \mathcal{O}_Y$，它作用在
$\mathcal{F}$ 上。于是 $f^{-1}\mathcal{I}\mathcal{F}$ 是由形如 $as$ 的
局部截面的和生成的子层，其中 $a$ 是 $f^{-1}\mathcal{I}$ 的局部截面，
$s$ 是 $\mathcal{F}$ 的局部截面。它是 $\mathcal{F}$ 的拟凝聚
$\mathcal{O}_Y$-子模，因为它也等于自然映射
$f^*\mathcal{I} \otimes_{\mathcal{O}_Y} \mathcal{F} \to \mathcal{F}$ 的像。

\medskip\noindent
说明这些记号后，证明便很直接。事实上，该问题在 $X$ 上是 \'etale 局部的，
因而可以假设 $X$ 是仿射概形。在这种情形下，结论由概形的相应结论推出，
见《概形的上同调》引理
\ref{coherent-lemma-affine-morphism-projection-ideal}。
\end{proof}

\begin{lemma}
\label{lemma-image-affine-finite-morphism-affine-Noetherian}
设 $S$ 为一个概形，$f : Y \to X$ 为 $S$ 上代数空间之间的态射。假设：
\begin{enumerate}
\item $f$ 有限；
\item $f$ 为满射；
\item $Y$ 仿射；
\item $X$ 是 Noether 的。
\end{enumerate}
则 $X$ 仿射。
\end{lemma}

\begin{proof}
我们将证明，在本引理的假设下，对于任意凝聚 $\mathcal{O}_X$-模
$\mathcal{F}$，都有 $H^1(X, \mathcal{F}) = 0$。由引理
\ref{lemma-directed-colimit-coherent} 和 \ref{lemma-colimits}，这又推出
对于每个拟凝聚 $\mathcal{O}_X$-模 $\mathcal{F}$，都有
$H^1(X, \mathcal{F}) = 0$。于是由命题
\ref{proposition-vanishing-affine}，$X$ 仿射。

\medskip\noindent
设 $\mathcal{P}$ 为 $X$ 上凝聚层 $\mathcal{F}$ 的下述性质：
$$
\mathcal{P}(\mathcal{F}) \Leftrightarrow H^1(X, \mathcal{F}) = 0.
$$
我们将应用引理 \ref{lemma-property-higher-rank-cohomological}，因此必须为
$\mathcal{P}$ 验证该引理的 (1)、(2)、(3)。性质 (1) 由与层的短正合列
相伴的上同调长正合列推出。性质 (2) 来自 $H^1(X, -)$ 是可加函子。
为证明 (3)，设 $i : Z \to X$ 为满足 $|Z|$ 不可约的约化闭子空间。
设 $i' : Z' \to Y$ 和 $f' : Z' \to Z$ 如引理
\ref{lemma-finite-morphism-Noetherian} 中所述，并令
$\mathcal{G} = f'_*\mathcal{O}_{Z'}$。我们断言 $\mathcal{G}$ 满足引理
\ref{lemma-property-higher-rank-cohomological} 的性质 (3)(a) 和 (3)(b)，
这将完成证明。性质 (3)(a) 已在引理
\ref{lemma-finite-morphism-Noetherian} 中证明。为证明 (3)(b)，设
$\mathcal{I}$ 为 $Z$ 上的非零拟凝聚理想层。记
$\mathcal{I}' \subset \mathcal{O}_{Z'}$ 为拟凝聚理想
$(f')^{-1}\mathcal{I} \mathcal{O}_{Z'}$，也就是
$(f')^*\mathcal{I} \to \mathcal{O}_{Z'}$ 的像。由引理
\ref{lemma-affine-morphism-projection-ideal}，有
$f_*\mathcal{I}' = \mathcal{I} \mathcal{G}$。% P09-ERR-0035
我们断言其共同值
$\mathcal{G}' = \mathcal{I} \mathcal{G} = f'_*\mathcal{I}'$
满足 (3)(b) 中的条件。首先，显然 $\mathcal{G}/\mathcal{G}'$ 的支撑包含在
$\mathcal{O}_Z/\mathcal{I}$ 的支撑中；后者是 $|Z|$ 的真子空间，因为
$\mathcal{I}$ 是约化不可约代数空间 $Z$ 上的非零理想层。态射 $f'$ 仿射，
故由引理 \ref{lemma-affine-vanishing-higher-direct-images}，
$R^1f'_*\mathcal{I}' = 0$。由于 $Z'$ 仿射（它是仿射概形的闭子概形），
有 $H^1(Z', \mathcal{I}') = 0$。因此 Leray 谱序列（取《站点上的上同调》
引理 \ref{sites-cohomology-lemma-apply-Leray} 的形式）推出
$H^1(Z, f'_*\mathcal{I}') = 0$。由于 $i : Z \to X$ 仿射，
$R^1i_*f'_*\mathcal{I}' = 0$；再次应用 Leray，得到
$H^1(X, i_*f'_*\mathcal{I}') = 0$。换言之，如所要求，有
$H^1(X, i_*\mathcal{G}') = 0$。
\end{proof}








\section{Chow 引理的弱形式}
\label{section-weak-chow}

\noindent
本节快速证明下面的引理，以帮助我们证明恰当代数空间上凝聚模上同调的
基本结果。

\begin{lemma}
\label{lemma-weak-chow}
设 $A$ 为一个环，$X$ 为 $\Spec(A)$ 上的代数空间，其结构态射
$X \to \Spec(A)$ 分离且为有限型。则存在恰当满态射 $X' \to X$，
其中 $X'$ 是 $\Spec(A)$ 上 H-拟射影的概形。
\end{lemma}

\begin{proof}
设 $W$ 为仿射概形，$f : W \to X$ 为满 \'etale 态射。存在整数 $d$，使得
$f$ 的所有几何纤维都至多有 $d$ 个点（因为 $X$ 是分离代数，因而合理，
见《良态空间》引理 \ref{decent-spaces-lemma-bounded-fibres}）。
% P09-ERR-0036
取最小的 $d$，便得到非空开集 $U \subset X$，使得
$f^{-1}(U) \to U$ 是次数为 $d$ 的有限 \'etale 态射，见《良态空间》引理
\ref{decent-spaces-lemma-quasi-compact-reasonable-stratification}。令
$$
V \subset W \times_X W \times_X \ldots \times_X W
$$
为所有对角线的补集（纤维积中有 $d$ 个因子）。由于 $W \to X$ 分离，
对角态射 $W \to W \times_X W$ 是闭浸入。由于 $W \to X$ 是 \'etale 的，
该对角态射也是开浸入，见《空间的态射》引理
\ref{spaces-morphisms-lemma-etale-unramified} 和
\ref{spaces-morphisms-lemma-diagonal-unramified-morphism}。因此这些对角线是
拟紧概形 $W \times_X \ldots \times_X W$ 的开闭子概形。特别地，$V$
是拟紧概形。选取开浸入 $W \subset Y$，其中 $Y$ 在 $A$ 上 H-射影
（这是可能的，因为 $W$ 仿射且在 $A$ 上为有限型；例如可以使用《态射》引理
\ref{morphisms-lemma-quasi-affine-finite-type-over-S} 和
\ref{morphisms-lemma-H-quasi-projective-open-H-projective}）。令
$$
Z \subset Y \times_A Y \times_A \ldots \times_A Y
$$
为复合态射
$V \to W \times_X \ldots \times_X W \to Y \times_A \ldots \times_A Y$
的概形论像。由于 $V$ 拟紧且 $Y \times_A \ldots \times_A Y$ 分离，
该态射拟紧。注意，$V \to Y \times_A \ldots \times_A Y$ 是浸入，故
$V \to Z$ 是开浸入，见《态射》引理
\ref{morphisms-lemma-quasi-compact-immersion}。投影态射给出 $d$ 个态射
$g_i : Z \to Y$。由于 $Y$ 在 $A$ 上射影，这些态射 $g_i$ 都是射影的，
见《态射》\ref{morphisms-section-projective} 节中的材料。令
$$
X' = \bigcup g_i^{-1}(W) \subset Z
$$
存在态射 $X' \to X$，它在 $g_i^{-1}(W)$ 上的限制是复合态射
$g_i^{-1}(W) \to W \to X$。事实上，这些态射在 $V$ 上相同，因而由
《空间的态射》引理 \ref{spaces-morphisms-lemma-equality-of-morphisms}，
它们在 $g_i^{-1}(W) \cap g_j^{-1}(W)$ 上相同。断言：态射
$X' \to X$ 恰当。

\medskip\noindent
如果断言成立，则对 $d$ 归纳即可证明本引理。事实上，按构造，$X'$ 在
$\Spec(A)$ 上 H-拟射影。由于 $V$ 满射到 $U$，$X' \to X$ 的像包含开集
$U$。记 $T$ 为 $X \setminus U$ 上的约化诱导代数空间结构。则
$T \times_X W$ 是 $W$ 的闭子概形，因而仿射。此外，态射
$T \times_X W \to T$ 是 \'etale 的，且每个几何纤维少于 $d$ 个点。
由归纳假设，存在恰当满态射 $T' \to T$，其中 $T'$ 是
$\Spec(A)$ 上 H-拟射影的概形。由于 $T$ 是 $X$ 的闭子空间，
$T' \to X$ 是恰当态射。因此取恰当满态射
$X' \amalg T' \to X$ 即得本引理。

\medskip\noindent
断言的证明。按构造，态射 $X' \to X$ 分离且为有限型。我们将为态射
$V \to X'$ 和 $X' \to X$ 验证《空间的态射》引理
\ref{spaces-morphisms-lemma-refined-valuative-criterion-universally-closed}
的条件 (1)--(4)。条件 (1) 和 (2) 已在上面验证。条件 (3) 成立，因为
$X' \to X$ 分离（其源是分离代数空间）。因此只需对图
$$
\xymatrix{
\Spec(K) \ar[r] \ar[d] & V \ar[d] \\
\Spec(R) \ar[r] & X
}
$$
检验到 $X'$ 的可提升性，其中 $R$ 是分式域为 $K$ 的赋值环。注意，上方
水平态射由 $W$ 的 $d$ 个两两不同的 $K$-值点
$w_1, \ldots, w_d$ 给出。事实上，它们恰好是该图所给点
$x \in X(K)$ 的全部逆像。由于 $W \to X$ 为满射，必要时用赋值环扩张
替换 $R$ 后，可以把态射 $\Spec(R) \to X$ 提升为态射
$w : \Spec(R) \to W$，见《空间的态射》引理
\ref{spaces-morphisms-lemma-lift-valuation-ring-through-flat-morphism}。
由于 $w_1, \ldots, w_d$ 是 $x$ 的全部逆像，我们看到
$w|_{\Spec(K)}$ 等于其中某一个，比如 $w_i$。% P09-ERR-0037
于是得到交换图
$$
\xymatrix{
\Spec(K) \ar[r] \ar[d] &  Z \ar[d]_{g_i}\\
\Spec(R) \ar[r]^w & Y
}
$$
由射影态射 $g_i$ 的恰当性赋值判别法，可以把 $w$ 提升为
$z : \Spec(R) \to Z$，见《态射》引理
\ref{morphisms-lemma-locally-projective-proper} 以及《概形》命题
\ref{schemes-proposition-characterize-universally-closed}。$z$ 的像位于
$g_i^{-1}(W) \subset X'$ 中，证明完成。
\end{proof}







\section{Noether 赋值判别法}
\label{section-Noetherian-valuative-criterion}

\noindent
我们使用离散赋值环证明恰当性赋值判别法的一个版本。更精确（因而也更
技术化）的版本见《空间的极限》
\ref{spaces-limits-section-Noetherian-valuative-criterion} 节。

\begin{lemma}
\label{lemma-check-separated-dvr}
设 $S$ 为一个概形，$f : X \to Y$ 为 $S$ 上代数空间之间的态射。假设：
\begin{enumerate}
\item $Y$ 局部 Noether；
\item $f$ 局部有限型且拟分离；
\item 对于每个交换图
$$
\xymatrix{
\Spec(K) \ar[r] \ar[d] & X \ar[d] \\
\Spec(A) \ar[r] \ar@{-->}[ru] & Y
}
$$
其中 $A$ 是离散赋值环，$K$ 是其分式域，使该图交换的虚线箭头至多有一个。
\end{enumerate}
则 $f$ 分离。
\end{lemma}

\begin{proof}
必须证明对角态射 $\Delta : X \to X \times_Y X$ 是闭浸入。我们已经知道
$\Delta$ 可表、分离，是单态射且局部有限型，见《空间的态射》引理
\ref{spaces-morphisms-lemma-properties-diagonal}。选取仿射概形 $U$ 和
\'etale 态射 $U \to X \times_Y X$。令
$V = X \times_{\Delta, X \times_Y X} U$。只需证明 $V \to U$ 是闭浸入
（《空间的态射》引理 \ref{spaces-morphisms-lemma-closed-immersion-local}）。
由于 $X \times_Y X$ 在 $Y$ 上局部有限型，$U$ 是 Noether 的（使用
《空间的态射》引理
\ref{spaces-morphisms-lemma-composition-finite-type},
\ref{spaces-morphisms-lemma-base-change-finite-type} 和
\ref{spaces-morphisms-lemma-locally-finite-type-locally-noetherian}）。
由于 $\Delta$ 可表，$V$ 是概形。此外，因为 $f$ 拟分离，$V$ 拟紧。
所以 $V \to U$ 为有限型。考虑概形态射的交换图
$$
\xymatrix{
\Spec(K) \ar[r] \ar[d] & V \ar[d] \\
\Spec(A) \ar[r] \ar@{-->}[ru] & U
}
$$
其中 $A$ 是分式域为 $K$ 的离散赋值环。可以把复合态射
$\Spec(A) \to U \to X \times_Y X$ 解释为一对态射
$a, b : \Spec(A) \to X$；它们复合到 $Y$ 后相同，且限制到
$\Spec(K)$ 后相等。因此假设 (3) 保证 $a = b$，从而得到图中的虚线箭头。
由《极限》引理 \ref{limits-lemma-Noetherian-dvr-valuative-proper}，
$V \to U$ 恰当。换言之，$\Delta$ 恰当。由于 $\Delta$ 是单态射，
《\'Etale 态射》引理 \ref{etale-lemma-characterize-closed-immersion}
表明 $\Delta$ 是闭浸入，如所要求。
\end{proof}

\begin{lemma}
\label{lemma-check-proper-dvr}
设 $S$ 为一个概形，$f : X \to Y$ 为 $S$ 上代数空间之间的态射。假设：
\begin{enumerate}
\item $Y$ 局部 Noether；
\item $f$ 有限型且拟分离；
\item 对于每个交换图
$$
\xymatrix{
\Spec(K) \ar[r] \ar[d] & X \ar[d] \\
\Spec(A) \ar[r] \ar@{-->}[ru] & Y
}
$$
其中 $A$ 是离散赋值环，$K$ 是其分式域，存在唯一的虚线箭头使该图交换。
\end{enumerate}
则 $f$ 恰当。
\end{lemma}

\begin{proof}
由引理 \ref{lemma-check-separated-dvr}，$f$ 分离，所以只需证明 $f$ 泛闭。
为此可以在 $Y$ 上 \'etale 局部地验证（《空间的态射》引理
\ref{spaces-morphisms-lemma-universally-closed-local}）。因此可以假设
$Y = \Spec(A)$ 是 Noether 仿射概形。按照 Chow 引理的弱形式（引理
\ref{lemma-weak-chow}）选取 $X' \to X$。我们断言
$X' \to \Spec(A)$ 泛闭。由《空间的态射》引理
\ref{spaces-morphisms-lemma-image-proper-is-proper}，该断言蕴含本引理。
为证明该断言，根据《极限》引理
\ref{limits-lemma-check-universally-closed-Noetherian}，只需证明在每个实线
部分交换的图
$$
\xymatrix{
\Spec(K) \ar[r] \ar[d] & X' \ar[r] & X \ar[d] \\
\Spec(A) \ar[rr] \ar@{-->}[ru]^a \ar@{-->}[rru]_b & & Y
}
$$
中，可以找到虚线箭头 $a$，其中 $A$ 是分式域为 $K$ 的离散赋值环。
% P09-ERR-0038
由假设可以找到虚线箭头 $b$。于是态射
$X' \times_{X, b} \Spec(A) \to \Spec(A)$ 是概形之间的恰当态射；
由概形态射的赋值判别法，可以把 $b$ 提升为所需态射 $a$。
\end{proof}

\begin{remark}[完备离散赋值环的变体]
\label{remark-variant}
在引理 \ref{lemma-check-separated-dvr} 和 \ref{lemma-check-proper-dvr} 中，
只需考虑完备离散赋值环。确切地说，在引理
\ref{lemma-check-separated-dvr} 中，可以把条件 (3) 替换为下述条件：
给定任意交换图
$$
\xymatrix{
\Spec(K) \ar[r] \ar[d] & X \ar[d] \\
\Spec(A) \ar[r] \ar@{-->}[ru] & Y
}
$$
其中 $A$ 是分式域为 $K$ 的完备离散赋值环，使该图交换的虚线箭头至多
有一个。事实上，给定引理 \ref{lemma-check-separated-dvr} (3) 中的任意图，
其完备化 $A^\wedge$ 是离散赋值环（《代数进阶》引理
\ref{more-algebra-lemma-completion-dvr}），而箭头
$\Spec(A^\wedge) \to X$ 的唯一性蕴含箭头 $\Spec(A) \to X$ 的唯一性；
例如这由《空间的性质》命题
\ref{spaces-properties-proposition-sheaf-fpqc} 推出。
类似地，在引理 \ref{lemma-check-proper-dvr} 中，可以把条件 (3) 替换为
下述条件：给定任意交换图
$$
\xymatrix{
\Spec(K) \ar[r] \ar[d] & X \ar[d] \\
\Spec(A) \ar[r] & Y
}
$$
其中 $A$ 是分式域为 $K$ 的完备离散赋值环，则存在完备离散赋值环的扩张
$A \subset A'$，诱导分式域扩张 $K \subset K'$，并存在唯一箭头
$\Spec(A') \to X$ 使图
$$
\xymatrix{
\Spec(K') \ar[r] \ar[d] & \Spec(K) \ar[r] & X \ar[d] \\
\Spec(A') \ar[r] \ar[rru] & \Spec(A) \ar[r] & Y
}
$$
交换。事实上，给定引理 \ref{lemma-check-proper-dvr} (3) 中的任意图，
只要对于离散赋值环的{\it 任意}扩张 $A \subset B$，存在交换图
$$
\xymatrix{
\Spec(L) \ar[r] \ar[d] & \Spec(K) \ar[r] & X \ar[d] \\
\Spec(B) \ar[r] \ar[rru] & \Spec(A) \ar[r] & Y
}
$$
就能推出存在箭头 $\Spec(A) \to X$ 嵌入该图。这已在《空间的态射》引理
\ref{spaces-morphisms-lemma-push-down-solution} 中证明。事实上，从这些考察
可以看出，只要某类离散赋值环满足：任给一个离散赋值环，都有它的一个
扩张属于该类，那么只在该类离散赋值环上的图中寻找虚线箭头就已足够。
例如，可以取剩余域代数闭的完备离散赋值环。
\end{remark}







\section{凝聚层的高阶正像}
\label{section-proper-pushforward}

\noindent
本节证明一个基本事实：凝聚层在恰当态射下的高阶正像是凝聚的。首先证明
一个辅助引理。

\begin{lemma}
\label{lemma-kill-by-twisting}
设 $S$ 为一个概形。考虑 $S$ 上代数空间的交换图
$$
\xymatrix{
X \ar[r]_i \ar[rd]_f & \mathbf{P}^n_Y \ar[d] \\
& Y
}
$$
假设 $i$ 是闭浸入，且 $Y$ 是 Noether 的。令
$\mathcal{L} = i^*\mathcal{O}_{\mathbf{P}^n_Y}(1)$。设 $\mathcal{F}$ 为
$X$ 上的凝聚模。则存在整数 $d_0$，使得对于所有 $d \geq d_0$ 和
所有 $p > 0$，都有
$R^pf_*(\mathcal{F} \otimes_{\mathcal{O}_X} \mathcal{L}^{\otimes d}) = 0$。
\end{lemma}

\begin{proof}
可以在 $Y$ 上 \'etale 局部地检验
$R^pf_*(\mathcal{F} \otimes \mathcal{L}^{\otimes d})$ 是否为零，见公式
(\ref{equation-representable-higher-direct-image})。因此可以假设 $Y$ 是
Noether 环的谱。在这种情形下，$X$ 是概形，结论由《概形的上同调》引理
\ref{coherent-lemma-kill-by-twisting} 推出。
\end{proof}

\begin{lemma}
\label{lemma-proper-pushforward-coherent}
设 $S$ 为一个概形，$f : X \to Y$ 为 $S$ 上代数空间之间的恰当态射，
且 $Y$ 局部 Noether。设 $\mathcal{F}$ 为凝聚 $\mathcal{O}_X$-模。
则对于所有 $i \geq 0$，$R^if_*\mathcal{F}$ 都是凝聚
$\mathcal{O}_Y$-模。
\end{lemma}

\begin{proof}
首先，由《空间的态射》引理
\ref{spaces-morphisms-lemma-locally-finite-type-locally-noetherian}，
$X$ 是局部 Noether 代数空间，故引理的陈述有意义。此外，计算
$R^if_*\mathcal{F}$ 与 $Y$ 上的 \'etale 局部化交换（《空间的性质》引理
\ref{spaces-properties-lemma-pushforward-etale-base-change-modules}），而且
可以在 $Y$ 上 \'etale 局部地检验 $R^if_*\mathcal{F}$ 是否凝聚（引理
\ref{lemma-coherent-Noetherian}）。因此可以假设
$Y = \Spec(A)$ 是 Noether 仿射概形。

\medskip\noindent
假设 $Y = \Spec(A)$ 是仿射概形。注意，$f$ 局部有限表示（《空间的态射》
引理 \ref{spaces-morphisms-lemma-noetherian-finite-type-finite-presentation}）。
因此它有限表示，从而 $X$ 是 Noether 的（《空间的态射》引理
\ref{spaces-morphisms-lemma-finite-presentation-noetherian}）。于是引理
\ref{lemma-property-higher-rank-cohomological-variant} 可应用于 $X$ 的
凝聚模范畴。对于 $X$ 上的凝聚层 $\mathcal{F}$，规定：$\mathcal{P}$
成立，当且仅当 $R^if_*\mathcal{F}$ 是 $\Spec(A)$ 上的凝聚模。
下面证明引理 \ref{lemma-property-higher-rank-cohomological-variant} 的
条件 (1)、(2)、(3) 对此性质成立，从而完成本引理的证明。

\medskip\noindent
验证条件 (1)。设
$$
0 \to \mathcal{F}_1 \to \mathcal{F}_2 \to \mathcal{F}_3 \to 0
$$
为 $X$ 上凝聚层的短正合列。考虑高阶正像的长正合列
$$
R^{p - 1}f_*\mathcal{F}_3 \to
R^pf_*\mathcal{F}_1 \to
R^pf_*\mathcal{F}_2 \to
R^pf_*\mathcal{F}_3 \to
R^{p + 1}f_*\mathcal{F}_1
$$
显然，如果三个层 $\mathcal{F}_i$ 中的两个具有性质 $\mathcal{P}$，
则第三个的高阶正像在这个正合复形中夹在两个凝聚层之间。因此由引理
\ref{lemma-coherent-abelian-Noetherian} 和
\ref{lemma-coherent-Noetherian-quasi-coherent-sub-quotient}，这些高阶正像
也凝聚。所以第三个层也具有性质 $\mathcal{P}$。

\medskip\noindent
验证条件 (2)。这立即来自等式
$R^if_*(\mathcal{F}_1 \oplus \mathcal{F}_2) =
R^if_*\mathcal{F}_1 \oplus R^if_*\mathcal{F}_2$，以及凝聚模的直和项凝聚
这一事实（见上面引用的引理）。

\medskip\noindent
验证条件 (3)。设 $i : Z \to X$ 为闭浸入，其中 $Z$ 约化且 $|Z|$ 不可约。
令 $g = f \circ i : Z \to \Spec(A)$。设 $\mathcal{G}$ 为 $Z$ 上概形论
支撑等于 $Z$ 的凝聚模，并且对于所有 $p$，$R^pg_*\mathcal{G}$ 凝聚。
则 $\mathcal{F} = i_*\mathcal{G}$ 是 $X$ 上概形论支撑为 $Z$ 的凝聚模，
且 $R^pf_*\mathcal{F} = R^pg_*\mathcal{G}$。为看出这一点，使用 Leray
谱序列（《站点上的上同调》引理
\ref{sites-cohomology-lemma-relative-Leray}），并使用引理
\ref{lemma-affine-vanishing-higher-direct-images} 以及闭浸入仿射这一事实
（《空间的态射》引理 \ref{spaces-morphisms-lemma-closed-immersion-affine}），
得到当 $q > 0$ 时 $R^qi_*\mathcal{G} = 0$。因此问题归结为寻找
$Z$ 上支撑等于 $Z$ 的凝聚层
$\mathcal{G}$，使得对于所有 $p$，$R^pg_*\mathcal{G}$ 凝聚。

\medskip\noindent
对态射 $Z \to \Spec(A)$ 应用引理 \ref{lemma-weak-chow}，得到图
$$
\xymatrix{
Z \ar[rd]_g & Z' \ar[d]^-{g'} \ar[l]^\pi \ar[r]_i & \mathbf{P}^n_A \ar[dl] \\
& \Spec(A) &
}
$$
其中 $\pi : Z' \to Z$ 恰当且为满射，$i$ 是浸入。由于
$Z \to \Spec(A)$ 恰当，由《空间的态射》引理
\ref{spaces-morphisms-lemma-composition-proper}，$g'$ 恰当。因此 $i$ 是
闭浸入（《空间的态射》引理
\ref{spaces-morphisms-lemma-universally-closed-permanence} 和
\ref{spaces-morphisms-lemma-immersion-when-closed}）。由此可知态射
$i' = (i, \pi) : \mathbf{P}^n_A \times_{\Spec(A)} Z' = \mathbf{P}^n_Z$ 为
闭浸入（《空间的态射》引理
\ref{spaces-morphisms-lemma-semi-diagonal}）。% P09-ERR-0039
令
$$
\mathcal{L} =
i^*\mathcal{O}_{\mathbf{P}^n_A}(1) =
(i')^*\mathcal{O}_{\mathbf{P}^n_Z}(1)
$$
可以分别对 $\mathcal{L}$ 与 $\pi$、$\mathcal{L}$ 与 $g'$ 应用引理
\ref{lemma-kill-by-twisting}。因此对于所有 $d \gg 0$ 和 $p > 0$，有
$R^p\pi_*\mathcal{L}^{\otimes d} = 0$ 且
$R^p(g')_*\mathcal{L}^{\otimes d} = 0$。令
$\mathcal{G} = \pi_*\mathcal{L}^{\otimes d}$。由 Leray 谱序列（《站点上的
上同调》引理 \ref{sites-cohomology-lemma-relative-Leray}），有
$$
E_2^{p, q} = R^pg_* R^q\pi_*\mathcal{L}^{\otimes d}
\Rightarrow
R^{p + q}(g')_*\mathcal{L}^{\otimes d}
$$
按 $d$ 的选取，$E_2^{p, q}$ 中仅 $q = 0$ 的项可能非零，而
$R^{p + q}(g')_*\mathcal{L}^{\otimes d}$ 中仅 $p = q = 0$ 的项可能非零。
这推出当 $p > 0$ 时 $R^pg_*\mathcal{G} = 0$，并且
$g_*\mathcal{G} = (g')_*\mathcal{L}^{\otimes d}$。应用《概形的上同调》
引理 \ref{coherent-lemma-locally-projective-pushforward}，可知
$g_*\mathcal{G} = (g')_*\mathcal{L}^{\otimes d}$ 凝聚。

\medskip\noindent
还必须检验 $\mathcal{G}$ 的支撑是 $Z$。这来自
$\mathcal{L}^{\otimes d}$ 有许多全局截面这一事实。下面详述。由于
$\mathcal{O}_{\mathbf{P}^n}(d)$ 对所有 $d \geq 0$ 都由全局截面生成，
$\mathcal{L}^{\otimes d}$ 也如此。选取 $z \in Z'$，使它映到 $Z$ 的一般点
$\xi$；由于 $\pi$ 为满射，这是可以做到的。（注意，$Z$ 确实有一般点：
$|Z|$ 不可约，且 $Z$ 是 Noether 的，因而拟分离；所以由《空间的性质》
引理 \ref{spaces-properties-lemma-quasi-separated-sober}，$|Z|$ 是 sober
拓扑空间。）选取在 $z$ 处不消失的
$s \in \Gamma(Z', \mathcal{L}^{\otimes d})$。由于
$\Gamma(Z, \mathcal{G}) = \Gamma(Z', \mathcal{L}^{\otimes d})$，可以把
$s$ 看作 $\mathcal{G}$ 的全局截面。选取 $Z'$ 上位于 $z$ 之上的几何点
$\overline{z}$，并把 $Z$ 上相应的几何点记为
$\overline{\xi} = g' \circ \overline{z}$。% P09-ERR-0040
伴随映射
$$
(g')^*\mathcal{G} =
(g')^*g'_*\mathcal{L}^{\otimes d} \longrightarrow \mathcal{L}^{\otimes d}
$$
诱导茎之间的映射
$\mathcal{G}_{\overline{\xi}} \to \mathcal{L}_{\overline{z}}$，见
《空间的性质》引理
\ref{spaces-properties-lemma-stalk-pullback-quasi-coherent}。此外，该伴随
映射把 $s$ 的拉回（把 $s$ 看作 $\mathcal{G}$ 的截面）送到 $s$
（把 $s$ 看作 $\mathcal{L}^{\otimes d}$ 的截面）。因此，$s$ 在下列箭头
源端向量空间中的像
$$
\mathcal{G}_{\overline{\xi}} \otimes \kappa(\overline{\xi})
\longrightarrow
\mathcal{L}^{\otimes d}_{\overline{z}} \otimes \kappa(\overline{z})
$$
非零，因为按 $s$ 的选取，它在箭头目标端的像非零。因此 $\xi$ 位于
$\mathcal{G}$ 的支撑中（《空间的态射》引理
\ref{spaces-morphisms-lemma-support-finite-type}）。由于 $|Z|$ 不可约且
$Z$ 约化，可知 $\mathcal{G}$ 的概形论支撑是整个 $Z$，如所要求。
\end{proof}

\begin{lemma}
\label{lemma-proper-over-affine-cohomology-finite}
设 $A$ 为 Noether 环，$f : X \to \Spec(A)$ 为代数空间的恰当态射，
并设 $\mathcal{F}$ 为凝聚 $\mathcal{O}_X$-模。则对于所有 $i \geq 0$，
$H^i(X, \mathcal{F})$ 是有限 $A$-模。
\end{lemma}

\begin{proof}
这正是引理 \ref{lemma-proper-pushforward-coherent} 的仿射情形。事实上，
由引理 \ref{lemma-higher-direct-image}，$R^if_*\mathcal{F}$ 是拟凝聚层，
因而是与 $A$-模
$\Gamma(\Spec(A), R^if_*\mathcal{F}) = H^i(X, \mathcal{F})$ 相伴的拟凝聚层。
该等式由《站点上的上同调》引理
\ref{sites-cohomology-lemma-apply-Leray} 以及仿射概形上拟凝聚模的高阶
上同调群消失（《概形的上同调》引理
\ref{coherent-lemma-quasi-coherent-affine-cohomology-zero}）推出。由引理
\ref{lemma-coherent-Noetherian}，$R^if_*\mathcal{F}$ 是凝聚层，当且仅当
$H^i(X, \mathcal{F})$ 是有限型 $A$-模。因此引理
\ref{lemma-proper-pushforward-coherent} 给出结论。
\end{proof}

\begin{lemma}
\label{lemma-graded-finiteness}
设 $A$ 为 Noether 环，$B$ 为有限生成的分次 $A$-代数，
$f : X \to \Spec(A)$ 为代数空间的恰当态射。令
$\mathcal{B} = f^*\widetilde B$。设 $\mathcal{F}$ 为有限型拟凝聚分次
$\mathcal{B}$-模。对于每个 $p \geq 0$，分次 $B$-模
$H^p(X, \mathcal{F})$ 都是有限 $B$-模。
\end{lemma}

\begin{proof}
为证明这一点，考虑纤维积图
$$
\xymatrix{
X' = \Spec(B) \times_{\Spec(A)} X
\ar[r]_-\pi \ar[d]_{f'} &
X \ar[d]^f \\
\Spec(B) \ar[r] &
\Spec(A)
}
$$
注意，$f'$ 是恰当态射，见《空间的态射》引理
\ref{spaces-morphisms-lemma-base-change-proper}。此外，$B$ 是有限生成
$A$-代数，因而是 Noether 的（《代数》引理
\ref{algebra-lemma-Noetherian-permanence}）。这推出 $X'$ 是 Noether
代数空间（《空间的态射》引理
\ref{spaces-morphisms-lemma-finite-presentation-noetherian}）。由《空间的态射》
引理 \ref{spaces-morphisms-lemma-affine-equivalence-algebras}，$X'$ 是拟凝聚
$\mathcal{O}_X$-代数 $\mathcal{B}$ 的相对谱。由于 $\mathcal{F}$ 是拟凝聚
$\mathcal{B}$-模，存在唯一的拟凝聚 $\mathcal{O}_{X'}$-模
$\mathcal{F}'$，使得 $\pi_*\mathcal{F}' = \mathcal{F}$，见《空间的态射》
引理 \ref{spaces-morphisms-lemma-affine-equivalence-modules}。由于
$\mathcal{F}$ 作为 $\mathcal{B}$-模是有限型的，$\mathcal{F}'$ 是有限型
$\mathcal{O}_{X'}$-模（细节从略）。换言之，$\mathcal{F}'$ 是凝聚
$\mathcal{O}_{X'}$-模（引理 \ref{lemma-coherent-Noetherian}）。由于态射
$\pi : X' \to X$ 仿射，有
$$
H^p(X, \mathcal{F}) = H^p(X', \mathcal{F}')
$$
这由引理 \ref{lemma-affine-vanishing-higher-direct-images} 和《站点上的上同调》
引理 \ref{sites-cohomology-lemma-apply-Leray} 推出。因此本引理由引理
\ref{lemma-proper-over-affine-cohomology-finite} 推出。
\end{proof}









\section{丰沛可逆层与上同调}
\label{section-ample-cohomology}

\noindent
下面用扭转后的上同调消失性，给出仿射基上恰当代数空间的丰沛性判别法。

\begin{lemma}
\label{lemma-vanshing-gives-ample}
设 $R$ 为 Noether 环，$X$ 为 $R$ 上的恰当代数空间，并设 $\mathcal{L}$
为可逆 $\mathcal{O}_X$-模。下列条件等价：
\begin{enumerate}
\item $X$ 是概形，且 $\mathcal{L}$ 在 $X$ 上丰沛；
\item 对于每个凝聚 $\mathcal{O}_X$-模 $\mathcal{F}$，存在
$n_0 \geq 0$，使得对于所有 $n \geq n_0$ 和 $p > 0$，都有
$H^p(X, \mathcal{F} \otimes \mathcal{L}^{\otimes n}) = 0$；
\item 对于每个凝聚 $\mathcal{O}_X$-模 $\mathcal{F}$，存在
$n \geq 1$，使得
$H^1(X, \mathcal{F} \otimes \mathcal{L}^{\otimes n}) = 0$。
\end{enumerate}
\end{lemma}

\begin{proof}
(1) $\Rightarrow$ (2) 由《概形的上同调》引理
\ref{coherent-lemma-vanshing-gives-ample} 推出。(2) $\Rightarrow$ (3) 显然。
(3) $\Rightarrow$ (1) 即引理 \ref{lemma-Noetherian-h1-zero-invertible}。
\end{proof}

\begin{lemma}
\label{lemma-surjective-finite-morphism-ample}
设 $R$ 为 Noether 环，$f : Y \to X$ 为 $R$ 上恰当代数空间之间的态射。
设 $\mathcal{L}$ 为可逆 $\mathcal{O}_X$-模。假设 $f$ 有限且为满射。
下列条件等价：
\begin{enumerate}
\item $X$ 是概形且 $\mathcal{L}$ 丰沛；
\item $Y$ 是概形且 $f^*\mathcal{L}$ 丰沛。
\end{enumerate}
\end{lemma}

\begin{proof}
假设 (1) 成立。有限态射（由概形）可表，所以 $Y$ 是概形，见《空间的态射》
引理 \ref{spaces-morphisms-lemma-integral-local}。于是 (2) 由
《概形的上同调》引理
\ref{coherent-lemma-surjective-finite-morphism-ample} 推出。

\medskip\noindent
假设 (2) 成立。对凝聚 $\mathcal{O}_X$-模 $\mathcal{F}$，定义下述性质
$P$：存在 $n_0$，使得对于所有 $n \geq n_0$ 和 $p > 0$，都有
$H^p(X, \mathcal{F} \otimes \mathcal{L}^{\otimes n}) = 0$。我们将证明
任意凝聚 $\mathcal{O}_X$-模 $\mathcal{F}$ 都具有 $P$；由引理
\ref{lemma-vanshing-gives-ample}，这蕴含 $\mathcal{L}$ 丰沛。我们将应用
引理 \ref{lemma-property-higher-rank-cohomological}，因此必须为 $P$ 验证
该引理的 (1)、(2)、(3)。性质 (1) 来自与层的短正合列相伴的上同调
长正合列，以及与可逆层作张量积是正合函子这一事实。性质 (2) 来自
$H^p(X, -)$ 是可加函子。

\medskip\noindent
为证明 (3)，设 $i : Z \to X$ 为满足 $|Z|$ 不可约的约化闭子空间。
设 $i' : Z' \to Y$ 和 $f' : Z' \to Z$ 如引理
\ref{lemma-finite-morphism-Noetherian} 中所述，并令
$\mathcal{G} = f'_*\mathcal{O}_{Z'}$。我们断言 $\mathcal{G}$ 满足引理
\ref{lemma-property-higher-rank-cohomological} 的性质 (3)(a) 和 (3)(b)，
这将完成证明。性质 (3)(a) 已在引理
\ref{lemma-finite-morphism-Noetherian} 中证明。为证明 (3)(b)，设
$\mathcal{I}$ 为 $Z$ 上的非零拟凝聚理想层。记
$\mathcal{I}' \subset \mathcal{O}_{Z'}$ 为拟凝聚理想
$(f')^{-1}\mathcal{I} \mathcal{O}_{Z'}$，也就是
$(f')^*\mathcal{I} \to \mathcal{O}_{Z'}$ 的像。由引理
\ref{lemma-affine-morphism-projection-ideal}，有
$f_*\mathcal{I}' = \mathcal{I} \mathcal{G}$。% P09-ERR-0041
我们断言其共同值
$\mathcal{G}' = \mathcal{I} \mathcal{G} = f'_*\mathcal{I}'$ 满足 (3)(b)
中的条件。首先，显然 $\mathcal{G}/\mathcal{G}'$ 的支撑包含在
$\mathcal{O}_Z/\mathcal{I}$ 的支撑中；后者是 $|Z|$ 的真子空间，因为
$\mathcal{I}$ 是约化不可约代数空间 $Z$ 上的非零理想层。回忆 $f'_*$、
$i_*$、$i'_*$ 把凝聚模变为凝聚模，见引理
\ref{lemma-finite-pushforward-coherent} 和 \ref{lemma-i-star-equivalence}。
由于 $Y$ 是概形且 $\mathcal{L}$ 丰沛，由引理
\ref{lemma-vanshing-gives-ample}，存在 $n_0$ 使得% P09-ERR-0042
$$
H^p(Y, i'_*\mathcal{I}' \otimes_{\mathcal{O}_Y} f^*\mathcal{L}^{\otimes n}) = 0
$$
对于 $n \geq n_0$ 和 $p > 0$ 成立。于是得到
\begin{align*}
H^p(X, i_*\mathcal{G}' \otimes_{\mathcal{O}_X} \mathcal{L}^{\otimes n})
& =
H^p(Z, \mathcal{G'} \otimes_{\mathcal{O}_Z} i^*\mathcal{L}^{\otimes n}) \\
& =
H^p(Z,
f'_*\mathcal{I}' \otimes_{\mathcal{O}_Z} i^*\mathcal{L}^{\otimes n})) \\
& =
H^p(Z, f'_*(\mathcal{I}' \otimes_{\mathcal{O}_{Z'}}
(f')^*i^*\mathcal{L}^{\otimes n})) \\
& =
H^p(Z, f'_*(\mathcal{I}' \otimes_{\mathcal{O}_{Z'}}
(i')^*f^*\mathcal{L}^{\otimes n})) \\
& =
H^p(Z',
\mathcal{I}' \otimes_{\mathcal{O}_{Z'}} (i')^*f^*\mathcal{L}^{\otimes n})) \\
& =
H^p(Y, i'_*\mathcal{I}' \otimes_{\mathcal{O}_Y} f^*\mathcal{L}^{\otimes n}) = 0
\end{align*}
% P09-ERR-0043
这里使用了投影公式和 Leray 谱序列（见《站点上的上同调》
\ref{sites-cohomology-section-projection-formula} 节和
\ref{sites-cohomology-section-leray} 节），以及引理
\ref{lemma-finite-higher-direct-image-zero}。这验证了引理
\ref{lemma-property-higher-rank-cohomological} 的性质 (3)(b)，如所要求。
\end{proof}










\section{形式函数定理}
\label{section-theorem-formal-functions}

\noindent
本节对应于《概形的上同调》
\ref{coherent-section-theorem-formal-functions} 节。建议读者先阅读该节。

\begin{situation}
\label{situation-formal-functions}
这里 $A$ 是 Noether 环，$I \subset A$ 是理想。此外，
$f : X \to \Spec(A)$ 是代数空间的恰当态射，$\mathcal{F}$ 是 $X$ 上的
凝聚层。
\end{situation}

\noindent
在这一情形下，用 $I^n\mathcal{F}$ 表示 $\mathcal{F}$ 的拟凝聚子模：
它作为 $\mathcal{O}_X$-模由 $\mathcal{F}$ 的局部截面与 $I^n$ 中元素的
乘积生成。换言之，它是映射
$f^*\widetilde{I} \otimes_{\mathcal{O}_X} \mathcal{F} \to \mathcal{F}$ 的像。
% P09-ERR-0044

\begin{lemma}
\label{lemma-cohomology-powers-ideal-times-F}
在情形 \ref{situation-formal-functions} 中。令
$B = \bigoplus_{n \geq 0} I^n$。则对于每个 $p \geq 0$，分次 $B$-模
$\bigoplus_{n \geq 0} H^p(X, I^n\mathcal{F})$ 是有限 $B$-模。
\end{lemma}

\begin{proof}
令 $\mathcal{B} = \bigoplus I^n\mathcal{O}_X = f^*\widetilde{B}$。
则 $\bigoplus I^n\mathcal{F}$ 是有限型分次 $\mathcal{B}$-模。因此结论由
引理 \ref{lemma-graded-finiteness} 推出。
\end{proof}

\begin{lemma}
\label{lemma-cohomology-powers-ideal-application}
在情形 \ref{situation-formal-functions} 中。对于每个 $p \geq 0$，
存在整数 $c \geq 0$，使得：
\begin{enumerate}
\item 对于所有 $n \geq c$，乘法映射
$I^{n - c} \otimes H^p(X, I^c\mathcal{F}) \to H^p(X, I^n\mathcal{F})$
是满射；
\item 对于所有 $n \geq 0$、$m \geq c$，映射
$H^p(X, I^{n + m}\mathcal{F}) \to H^p(X, I^n\mathcal{F})$ 的像包含在
子模 $I^{m - c} H^p(X, I^n\mathcal{F})$ 中。
\end{enumerate}
\end{lemma}

\begin{proof}
由引理 \ref{lemma-cohomology-powers-ideal-times-F}，可以找到
$d_1, \ldots, d_t \geq 0$ 以及
$x_i \in H^p(X, I^{d_i}\mathcal{F})$，使得
$\bigoplus_{n \geq 0} H^p(X, I^n\mathcal{F})$ 作为
$B = \bigoplus_{n \geq 0} I^n$-模由 $x_1, \ldots, x_t$ 生成。取
$c = \max\{d_i\}$。(1) 显然成立。对于 (2)，令
$b = \max(0, n - c)$。考虑 $A$-模的交换图
$$
\xymatrix{
I^{n + m - c - b} \otimes I^b \otimes
H^p(X, I^c\mathcal{F}) \ar[r] \ar[d] &
I^{n + m - c} \otimes H^p(X, I^c\mathcal{F}) \ar[r] &
H^p(X, I^{n + m}\mathcal{F}) \ar[d] \\
I^{n + m - c - b} \otimes H^p(X, I^n\mathcal{F}) \ar[rr] & &
H^p(X, I^n\mathcal{F})
}
$$
由本引理的 (1)，当 $n + m \geq c$ 时，水平箭头的复合是满射。
另一方面，显然 $n + m - c - b \geq m - c$。所以 (2) 成立。
\end{proof}

\begin{lemma}
\label{lemma-ML-cohomology-powers-ideal}
在情形 \ref{situation-formal-functions} 中。固定 $p \geq 0$。
\begin{enumerate}
\item 存在 $c_1 \geq 0$，使得对于所有 $n \geq c_1$，有
$$
\Ker(
H^p(X, \mathcal{F}) \to H^p(X, \mathcal{F}/I^n\mathcal{F})
)
\subset
I^{n - c_1}H^p(X, \mathcal{F}).
$$
\item 逆系统
$$
\left(H^p(X, \mathcal{F}/I^n\mathcal{F})\right)_{n \in \mathbf{N}}
$$
满足 Mittag-Leffler 条件（见《同调》定义
\ref{homology-definition-Mittag-Leffler}）。
\item 事实上，对于任意 $p$ 和 $n$，存在 $c_2(n) \geq n$，使得
$$
\Im(H^p(X, \mathcal{F}/I^k\mathcal{F})
\to H^p(X, \mathcal{F}/I^n\mathcal{F}))
=
\Im(H^p(X, \mathcal{F})
\to H^p(X, \mathcal{F}/I^n\mathcal{F}))
$$
对于所有 $k \geq c_2(n)$ 成立。
\end{enumerate}
\end{lemma}

\begin{proof}
令 $c_1 = \max\{c_p, c_{p + 1}\}$，其中 $c_p,c_{p +1}$ 是引理
\ref{lemma-cohomology-powers-ideal-application} 分别对 $H^p$ 和
$H^{p + 1}$ 给出的整数。我们将在 (1)、(2)、(3) 的证明中使用这个常数。

\medskip\noindent
先证明 (1)。考虑短正合列
$$
0 \to I^n\mathcal{F} \to \mathcal{F} \to \mathcal{F}/I^n\mathcal{F} \to 0
$$
由上同调长正合列可见
$$
\Ker(
H^p(X, \mathcal{F}) \to H^p(X, \mathcal{F}/I^n\mathcal{F})
)
=
\Im(
H^p(X, I^n\mathcal{F}) \to H^p(X, \mathcal{F})
)
$$
因此按 $c_1$ 的选取，当 $n \geq c_1$ 时，它包含在
$I^{n - c_1}H^p(X, \mathcal{F})$ 中。

\medskip\noindent
注意，按 Mittag-Leffler 条件的定义，(3) 蕴含 (2)。

\medskip\noindent
下面证明 (3)。在余下的证明中固定 $n$。考虑交换图
$$
\xymatrix{
0 \ar[r] &
I^n\mathcal{F} \ar[r] &
\mathcal{F} \ar[r] &
\mathcal{F}/I^n\mathcal{F} \ar[r] &
0 \\
0 \ar[r] &
I^{n + m}\mathcal{F} \ar[r] \ar[u] &
\mathcal{F} \ar[r] \ar[u] &
\mathcal{F}/I^{n + m}\mathcal{F} \ar[r] \ar[u] &
0
}
$$
它诱导出下面的交换图
$$
\xymatrix{
H^p(X, I^n\mathcal{F}) \ar[r] &
H^p(X, \mathcal{F}) \ar[r] &
H^p(X, \mathcal{F}/I^n\mathcal{F}) \ar[r]_\delta &
H^{p + 1}(X, I^n\mathcal{F}) \\
H^p(X, I^{n + m}\mathcal{F}) \ar[r] \ar[u] &
H^p(X, \mathcal{F}) \ar[r] \ar[u]^1 &
H^p(X, \mathcal{F}/I^{n + m}\mathcal{F}) \ar[r] \ar[u] &
H^{p + 1}(X, I^{n + m}\mathcal{F}) \ar[u]^a
}
$$
如果 $m \geq c_1$，则 $a$ 的像包含在
$I^{m - c_1} H^{p + 1}(X, I^n\mathcal{F})$ 中。由 Artin--Rees 引理
（见《代数》引理 \ref{algebra-lemma-map-AR}），存在整数 $c_3(n)$，使得
$$
I^N H^{p + 1}(X, I^n\mathcal{F}) \cap \Im(\delta)
\subset
\delta\left(I^{N - c_3(n)}H^p(X, \mathcal{F}/I^n\mathcal{F})\right)
$$
对于所有 $N \geq c_3(n)$ 成立。由于
$H^p(X, \mathcal{F}/I^n\mathcal{F})$ 被 $I^n$ 零化，可见如果
$m \geq c_3(n) + c_1 + n$，则
$$
\Im(H^p(X, \mathcal{F}/I^{n + m}\mathcal{F})
\to H^p(X, \mathcal{F}/I^n\mathcal{F}))
=
\Im(H^p(X, \mathcal{F})
\to H^p(X, \mathcal{F}/I^n\mathcal{F}))
$$
换言之，取 $c_2(n) = c_3(n) + c_1 + n$，(3) 成立。
\end{proof}

\begin{theorem}[形式函数定理]
\label{theorem-formal-functions}
在情形 \ref{situation-formal-functions} 中，固定 $p \geq 0$。映射系统
$$
H^p(X, \mathcal{F})/I^nH^p(X, \mathcal{F})
\longrightarrow
H^p(X, \mathcal{F}/I^n\mathcal{F})
$$
诱导极限之间的同构
$$
H^p(X, \mathcal{F})^\wedge
\longrightarrow
\lim_n H^p(X, \mathcal{F}/I^n\mathcal{F})
$$
其中左端是 $A$-模 $H^p(X, \mathcal{F})$ 关于理想 $I$ 的完备化，见
《代数》\ref{algebra-section-completion} 节。此外，对于极限拓扑，
这实际上是同胚。
\end{theorem}

\begin{proof}
事实上，这立即由引理 \ref{lemma-ML-cohomology-powers-ideal} 推出。
下面写出细节。令 $M = H^p(X, \mathcal{F})$ 且
$M_n = H^p(X, \mathcal{F}/I^n\mathcal{F})$，并记
$N_n = \Im(M \to M_n)$。由《同调》
\ref{homology-section-inverse-systems} 节对极限的描述，有
$$
\lim_n M_n
=
\{(x_n) \in \prod M_n \mid \varphi_i(x_n) = x_{n - 1}, \ n = 2, 3, \ldots\}
$$
选取元素 $x = (x_n) \in \lim_n M_n$。% P09-ERR-0045
由引理 \ref{lemma-ML-cohomology-powers-ideal} (3)，对于所有 $n$ 都有
$x_n \in N_n$，因为按定义，对于所有 $m$，$x_n$ 是某个
$x_{n + m} \in M_{n + m}$ 的像。由引理
\ref{lemma-ML-cohomology-powers-ideal} (1)，约化映射存在分解
$$
M \to N_n \to M/I^{n - c_1}M
$$
对于 $n \geq c_1$，记 $y_n \in M/I^{n - c_1}M$ 为 $x_n$ 的像。
当 $n' \geq n$ 时，复合映射 $M \to M_{n'} \to M_n$ 是给定的映射
$M \to M_n$，所以在典范映射
$M/I^{n' - c_1}M \to M/I^{n - c_1}M$ 下，$y_{n'}$ 映到 $y_n$。
因此 $y = (y_{n + c_1})$ 定义了 $\lim_n M/I^nM$ 的一个元素。
我们略去验证 $y$ 在本引理的映射
$$
M^\wedge = \lim_n M/I^nM \longrightarrow \lim_n M_n
$$
下映到 $x$，也略去关于拓扑的验证。
\end{proof}

\begin{lemma}
\label{lemma-spell-out-theorem-formal-functions}
设 $A$ 为一个环，$I \subset A$ 为理想。假设 $A$ 是 Noether 的并关于
$I$ 完备。设 $f : X \to \Spec(A)$ 为代数空间的恰当态射，
$\mathcal{F}$ 为 $X$ 上的凝聚层。则
$$
H^p(X, \mathcal{F}) = \lim_n H^p(X, \mathcal{F}/I^n\mathcal{F})
$$
对于所有 $p \geq 0$ 成立。
\end{lemma}

\begin{proof}
这是形式函数定理（定理 \ref{theorem-formal-functions}）在完备 Noether
基环情形下的改述。事实上，在这种情形下，$A$-模
$H^p(X, \mathcal{F})$ 是有限的（引理
\ref{lemma-proper-over-affine-cohomology-finite}），因而是 $I$-进完备的
（《代数》引理 \ref{algebra-lemma-completion-tensor}）。所以左端不必再
作完备化。
\end{proof}

\begin{lemma}
\label{lemma-formal-functions-stalk}
设 $S$ 为一个概形，$f : X \to Y$ 为 $S$ 上代数空间之间的态射，
$\mathcal{F}$ 为 $X$ 上的拟凝聚层。假设：
\begin{enumerate}
\item $Y$ 局部 Noether；
\item $f$ 恰当；
\item $\mathcal{F}$ 凝聚。
\end{enumerate}
设 $\overline{y}$ 为 $Y$ 的几何点。考虑纤维
$X_1 = X_{\overline{y}}$ 的“无穷小邻域”
$$
\xymatrix{
X_n =
\Spec(\mathcal{O}_{Y, \overline{y}}/\mathfrak m_{\overline{y}}^n) \times_Y X
\ar[r]_-{i_n} \ar[d]_{f_n} &
X \ar[d]^f \\
\Spec(\mathcal{O}_{Y, \overline{y}}/\mathfrak m_{\overline{y}}^n)
\ar[r]^-{c_n} & Y
}
$$
并令 $\mathcal{F}_n = i_n^*\mathcal{F}$。则有
$$
\left(R^pf_*\mathcal{F}\right)_{\overline{y}}^\wedge
\cong
\lim_n H^p(X_n, \mathcal{F}_n)
$$
作为 $\mathcal{O}_{Y, \overline{y}}^\wedge$-模的同构。
\end{lemma}

\begin{proof}
这只是形式函数定理（定理 \ref{theorem-formal-functions}）一个特殊情形的
改述。下面详述。注意，$\mathcal{O}_{Y, \overline{y}}$ 是 Noether 局部环，
见《空间的性质》引理
\ref{spaces-properties-lemma-Noetherian-local-ring-Noetherian}。考虑典范态射
$c : \Spec(\mathcal{O}_{Y, \overline{y}}) \to Y$。它等同局部环，因而是
平坦态射。记 $f' : X' \to \Spec(\mathcal{O}_{Y, \overline{y}})$ 为 $f$
到这个局部环上的基变换。由引理 \ref{lemma-flat-base-change-cohomology}，
有 $c^*R^pf_*\mathcal{F} = R^pf'_*\mathcal{F}'$。此外，对于所有
$n \geq 1$，有典范等同 $X_n = X'_n$。

\medskip\noindent
因此可以假设 $Y = \Spec(A)$，其中 $A$ 是极大理想为 $\mathfrak m$ 的
严格 Hensel Noether 局部环，并且 $\overline{y} \to Y$ 等于
$\Spec(A/\mathfrak m) \to Y$。于是
$$
\left(R^pf_*\mathcal{F}\right)_{\overline{y}} =
\Gamma(Y, R^pf_*\mathcal{F}) =
H^p(X, \mathcal{F})
$$
因为 $(Y, \overline{y})$ 是 $\overline{y}$ 的 \'etale 邻域范畴中的始对象。
各态射 $c_n$ 都是闭浸入，所以其基变换 $i_n$ 也都是闭浸入。注意
$i_{n, *}\mathcal{F}_n = i_{n, *}i_n^*\mathcal{F}
= \mathcal{F}/\mathfrak m^n\mathcal{F}$。由 $i_n$ 的 Leray 谱序列以及
引理 \ref{lemma-finite-pushforward-coherent}，可见
$$
H^p(X_n, \mathcal{F}_n) =
H^p(X, i_{n, *}\mathcal{F}) =
H^p(X, \mathcal{F}/\mathfrak m^n\mathcal{F})
$$
% P09-ERR-0046
因此确实可以应用形式函数定理来计算引理陈述中的极限，结论得证。
\end{proof}

\noindent
下面的引理将在后一个引理中推广到维数 $>0$ 的纤维。

\begin{lemma}
\label{lemma-higher-direct-images-zero-finite-fibre}
设 $S$ 为一个概形，$f : X \to Y$ 为 $S$ 上代数空间之间的态射，
$\overline{y}$ 为 $Y$ 的几何点。假设：
\begin{enumerate}
\item $Y$ 局部 Noether；
\item $f$ 恰当；
\item $X_{\overline{y}}$ 的底拓扑空间是离散的。
\end{enumerate}
则对于 $X$ 上的任意凝聚层 $\mathcal{F}$ 和所有 $p > 0$，都有
$(R^pf_*\mathcal{F})_{\overline{y}} = 0$。
\end{lemma}

\begin{proof}
设 $\kappa(\overline{y})$ 为局部环 $\mathcal{O}_{Y, \overline{y}}$ 的剩余域。
如引理 \ref{lemma-formal-functions-stalk} 中那样，令
$X_{\overline{y}} = X_1 = \Spec(\kappa(\overline{y})) \times_Y X$。
由《空间的态射》引理 \ref{spaces-morphisms-lemma-quasi-finite-at-point}，
态射 $f : X \to Y$ 在 $X \to Y$ 位于 $\overline{y}$ 上的纤维的每个点处
拟有限。因此 $X_{\overline{y}} \to \overline{y}$ 分离且拟有限。于是由
《空间的态射》命题
\ref{spaces-morphisms-proposition-locally-quasi-finite-separated-over-scheme}，
$X_{\overline{y}}$ 是概形。由于它拟紧，其底拓扑空间是有限离散空间。
再由《概形》引理 \ref{schemes-lemma-scheme-finite-discrete-affine}，它是
仿射概形。由引理
\ref{lemma-image-affine-finite-morphism-affine-Noetherian}，各代数空间
$X_n$ 也都是仿射概形。此外，每个 $X_n$ 的底拓扑与 $X_1$ 的相同。
% P09-ERR-0047
因此对于所有 $p > 0$，有 $H^p(X_n, \mathcal{F}_n) = 0$。由引理
\ref{lemma-formal-functions-stalk}，可见
$(R^pf_*\mathcal{F})_{\overline{y}}^\wedge = 0$。由引理
\ref{lemma-proper-pushforward-coherent}，$R^pf_*\mathcal{F}$ 凝聚，因而
$R^pf_*\mathcal{F}_{\overline{y}}$ 是有限
$\mathcal{O}_{Y, \overline{y}}$-模。% P09-ERR-0048
由《代数》引理 \ref{algebra-lemma-completion-tensor}，这推出
$(R^pf_*\mathcal{F})_{\overline{y}} = 0$。
\end{proof}

\begin{lemma}
\label{lemma-higher-direct-images-zero-above-dimension-fibre}
\begin{slogan}
对于恰当映射，次数大于纤维维数的高阶正像之茎为零。
\end{slogan}
设 $S$ 为一个概形，$f : X \to Y$ 为 $S$ 上代数空间之间的态射，
$\overline{y}$ 为 $Y$ 的几何点。假设：
\begin{enumerate}
\item $Y$ 局部 Noether；
\item $f$ 恰当；
\item $\dim(X_{\overline{y}}) = d$.
\end{enumerate}
则对于 $X$ 上的任意凝聚层 $\mathcal{F}$ 和所有 $p > d$，都有
$(R^pf_*\mathcal{F})_{\overline{y}} = 0$。
\end{lemma}

\begin{proof}
设 $\kappa(\overline{y})$ 为局部环 $\mathcal{O}_{Y, \overline{y}}$ 的剩余域。
如引理 \ref{lemma-formal-functions-stalk} 中那样，令
$X_{\overline{y}} = X_1 = \Spec(\kappa(\overline{y})) \times_Y X$。
此外，每个无穷小邻域 $X_n$ 的底拓扑空间都与 $X_{\overline{y}}$ 的相同。
因此由引理 \ref{lemma-vanishing-above-dimension}，对于所有 $p > d$，有
$H^p(X_n, \mathcal{F}_n) = 0$。进而由引理
\ref{lemma-formal-functions-stalk}，当 $p > d$ 时
$(R^pf_*\mathcal{F})_{\overline{y}}^\wedge = 0$。由引理
\ref{lemma-proper-pushforward-coherent}，$R^pf_*\mathcal{F}$ 凝聚，因而
$R^pf_*\mathcal{F}_{\overline{y}}$ 是有限
$\mathcal{O}_{Y, \overline{y}}$-模。% P09-ERR-0048
由《代数》引理 \ref{algebra-lemma-completion-tensor}，这推出
$(R^pf_*\mathcal{F})_{\overline{y}} = 0$。
\end{proof}






\section{形式函数定理的应用}
\label{section-applications-formal-functions}

\noindent
我们将在需要时在这里加入更多内容。

\begin{lemma}
\label{lemma-characterize-finite}
（更一般的版本见《空间的态射进阶》引理
\ref{spaces-more-morphisms-lemma-characterize-finite}。）
设 $S$ 为一个概形，$f : X \to Y$ 为 $S$ 上代数空间之间的态射。
假设 $Y$ 局部 Noether。下列条件等价：
\begin{enumerate}
\item $f$ 有限；
\item $f$ 恰当，并且对于每个态射 $\Spec(k) \to Y$，其中 $k$ 是域，
$|X_k|$ 都是离散空间。
\end{enumerate}
\end{lemma}

\begin{proof}
由《空间的态射》引理 \ref{spaces-morphisms-lemma-finite-proper}，有限态射
恰当。由《空间的态射》引理
\ref{spaces-morphisms-lemma-finite-quasi-finite}，有限态射拟有限。
拟有限态射的纤维 $X_k$ 离散，见《空间的态射》引理
\ref{spaces-morphisms-lemma-locally-quasi-finite}。因此有限态射恰当且具有
离散纤维 $X_k$。

\medskip\noindent
假设 $f$ 恰当且具有离散纤维 $X_k$。我们要证明 $f$ 有限。事实上，只需
证明 $f$ 仿射。如果 $f$ 仿射，则由《空间的态射》引理
\ref{spaces-morphisms-lemma-integral-universally-closed}，$f$ 是整态射；
随后由《空间的态射》引理 \ref{spaces-morphisms-lemma-finite-integral}，
$f$ 有限。

\medskip\noindent
为证明 $f$ 仿射，可以假设 $Y$ 仿射；目标是证明 $X$ 也仿射。由于 $f$
恰当，$X$ 分离且拟紧。我们将证明，对于任意凝聚
$\mathcal{O}_X$-模 $\mathcal{F}$，都有 $H^1(X, \mathcal{F}) = 0$。
由引理 \ref{lemma-directed-colimit-coherent} 和 \ref{lemma-colimits}，
这推出对于每个拟凝聚 $\mathcal{O}_X$-模 $\mathcal{F}$，都有
$H^1(X, \mathcal{F}) = 0$。进而由命题
\ref{proposition-vanishing-affine}，$X$ 仿射。由引理
\ref{lemma-higher-direct-images-zero-finite-fibre}，对于 $Y$ 的所有几何点，
$R^1f_*\mathcal{F}$ 的茎都为零。换言之，$R^1f_*\mathcal{F} = 0$。
因此由 $f$ 的 Leray 谱序列，
$H^1(X , \mathcal{F}) = H^1(Y, f_*\mathcal{F})$。由于 $Y$ 仿射且
$f_*\mathcal{F}$ 拟凝聚（《空间的态射》引理
\ref{spaces-morphisms-lemma-pushforward}），《概形的上同调》引理
\ref{coherent-lemma-quasi-coherent-affine-cohomology-zero} 给出
$H^1(Y, f_*\mathcal{F}) = 0$。所以 $H^1(X, \mathcal{F}) = 0$，如所要求。
\end{proof}

\noindent
作为推论，得到下面的有用结果。

\begin{lemma}
\label{lemma-proper-finite-fibre-finite-in-neighbourhood}
（更一般的版本见《空间的态射进阶》引理
\ref{spaces-more-morphisms-lemma-proper-finite-fibre-finite-in-neighbourhood}。）
设 $S$ 为一个概形，$f : X \to Y$ 为 $S$ 上代数空间之间的态射，
$\overline{y}$ 为 $Y$ 的几何点。假设：
\begin{enumerate}
\item $Y$ 局部 Noether；
\item $f$ 恰当；
\item $|X_{\overline{y}}|$ 有限。
\end{enumerate}
则存在 $\overline{y}$ 的开邻域 $V \subset Y$，使得
$f|_{f^{-1}(V)} : f^{-1}(V) \to V$ 有限。
\end{lemma}

\begin{proof}
由《空间的态射》引理 \ref{spaces-morphisms-lemma-quasi-finite-at-point}，
态射 $f$ 在 $X$ 中位于 $\overline{y}$ 上的所有几何点处拟有限。由
《空间的态射》引理
\ref{spaces-morphisms-lemma-locally-finite-type-quasi-finite-part}，$f$
拟有限的点所成集合是开子空间 $U \subset X$。令 $Z = X \setminus U$。
则 $\overline{y} \not \in f(Z)$。由于 $f$ 恰当，集合 $f(Z) \subset Y$ 闭。
选取 $\overline{y}$ 的任意开邻域 $V \subset Y$，使得
$Z \cap V = \emptyset$。% P09-ERR-0049
则 $f^{-1}(V) \to V$ 局部拟有限且恰当。因此
$f^{-1}(V) \to V$ 具有离散纤维 $X_k$（《空间的态射》引理
\ref{spaces-morphisms-lemma-locally-quasi-finite}），这些纤维拟紧，因而有限。
所以由引理 \ref{lemma-characterize-finite}，$f^{-1}(V) \to V$ 有限。
\end{proof}




\input{chapters}

\bibliography{my}
\bibliographystyle{amsalpha}

\end{document}
